% Further Curve Sketching and Inequalities — Further Mathematics Upper Sixth — Cours signé OnBuch+
% Official MINESEC syllabus. Compiler avec : tectonic FMA07-further-curve-sketching-and-inequalities.tex
\newcommand{\DOCMATIERE}{Further Mathematics}
\newcommand{\DOCNIVEAU}{Upper Sixth}
\newcommand{\DOCMODULE}{Upper Sixth — Further mathematics}
\newcommand{\DOCLECON}{7}
\newcommand{\DOCTITRE}{Further Curve Sketching and Inequalities}
% OnBuch+ — préambule « cours signés » ENRICHI (compilé avec tectonic / XeLaTeX)
% Système visuel « Pop » d'OnBuch+ : fond crème (jamais blanc), bordures encre
% épaisses, ombres « dures » décalées, titres Archivo Black en capitales.
% Version MATHÉMATIQUES : graphes (pgfplots/tikz), boîtes pédagogiques
% (définition, propriété, méthode, attention, le savais-tu, exercice résolu,
% exercice, TP, bilan), page de garde et signature OnBuch+ ancrée en bas.
%
% Macros attendues dans main.tex AVANT \input{preamble} :
%   \DOCMATIERE  \DOCNIVEAU  \DOCMODULE  \DOCLECON (numéro)  \DOCTITRE
\documentclass[11pt]{article}

\usepackage{fontspec}
\usepackage[english]{babel}
\usepackage[a4paper,margin=2cm,top=3.4cm,bottom=2.8cm,headheight=1.4cm,footskip=1.3cm]{geometry}
\usepackage{amsmath,amssymb,amsfonts}
\usepackage{mathtools} % \xmapsto, \coloneqq... souvent utilisés par les modèles
\usepackage{cancel} % simplifications barrées (bilans de forces)
% \abs{z} (module, valeur absolue) et \norm{u} (norme), utilisés par les modèles
\providecommand{\abs}{}\let\abs\relax\DeclarePairedDelimiter{\abs}{\lvert}{\rvert}
\providecommand{\norm}{}\let\norm\relax\DeclarePairedDelimiter{\norm}{\lVert}{\rVert}
\usepackage[table]{xcolor} % [table] : fournit aussi \rowcolors (alternance de lignes)
\usepackage{pagecolor}
\usepackage{tikz}
\usetikzlibrary{arrows.meta,positioning,calc,shapes.geometric,shapes.misc,shapes.arrows,decorations.pathmorphing,decorations.pathreplacing,decorations.markings,patterns,shadows,mindmap,trees,fit,backgrounds,matrix,intersections,angles,quotes}
\usepackage{pgfplots}
\usepackage[version=4]{mhchem} % équations nucléaires \ce{^{235}_{92}U}
\usepackage[european,siunitx]{circuitikz} % schémas électriques
\pgfplotsset{compat=1.18}
\usepgfplotslibrary{fillbetween,groupplots} % aires entre courbes (calcul intégral)
\usepackage{enumitem}
\usepackage{fancyhdr}
% [most] charge toutes les bibliothèques tcolorbox (listings, minted, raster,
% documentation...) et gonfle inutilement la mémoire TeX sur un document long
% (~300 boîtes) : on ne charge que ce qui est réellement utilisé.
\usepackage[skins,breakable]{tcolorbox}
\usepackage{graphicx}
\usepackage{colortbl}
\usepackage{booktabs}
\usepackage{tabularx}
\usepackage{diagbox} % case d'en-tête diagonale des tableaux à double entrée
% Colonnes extensibles centrée / gauche / droite (C, L, R), souvent utilisées par les modèles
\newcolumntype{C}{>{\centering\arraybackslash}X}
\newcolumntype{L}{>{\raggedright\arraybackslash}X}
\newcolumntype{R}{>{\raggedleft\arraybackslash}X}
\usepackage{array}
\usepackage{multicol}
\usepackage{siunitx}
% parse-numbers=false : accepte aussi \SI{2,5 \times 10^{-3}}{...}
\sisetup{locale=UK,output-decimal-marker={.},group-minimum-digits=4,parse-numbers=false}
\DeclareSIUnit\ohm{\ensuremath{\symup{\Omega}}} % U+2126 absent de la police
\DeclareSIUnit\division{div} % réglages d'oscilloscope (ms/div, V/div)
\DeclareSIUnit\tr{tr} % tours (vitesse de rotation, ex. \SI{1500}{\tr\per\minute})
\usepackage{qrcode}
% \degree : commande fréquemment utilisée par les modèles pour un angle ou une
% température en dehors de \SI{}{} (non fournie par les paquets chargés).
\providecommand{\degree}{^{\circ}}
% \qed (fin de démonstration), utilisé par les modèles sans amsthm
\providecommand{\qed}{\hfill$\square$}
% \barre : double barre des valeurs interdites dans un tableau de variations (tkz-tab)
\providecommand{\barre}{\ensuremath{\|}}
% environnement proof (amsthm non chargé) utilisé par les modèles
\ifdefined\proof\else\newenvironment{proof}[1][Proof]{\par\noindent\textit{#1.}\ }{\qed\par}\fi
\usepackage{lastpage}
\usepackage{hyperref}
\hypersetup{hidelinks}

% ============================================================
% Palette « Pop » OnBuch+ — mêmes hex que la classe P (plus_ui.dart)
% ============================================================
\definecolor{poporange}{HTML}{F59321}
\definecolor{popdark}{HTML}{E07A0C}
\definecolor{popcream}{HTML}{FFF6E8}
\definecolor{popink}{HTML}{1C1714}
\definecolor{poppurple}{HTML}{7A5AE0}
\definecolor{popgreen}{HTML}{2E9E6A}
\definecolor{poppink}{HTML}{E0457B}
\definecolor{popblue}{HTML}{2F7DE1}
\definecolor{popgold}{HTML}{C98A12}
\definecolor{popmuted}{HTML}{8A7F76}
\definecolor{popline}{HTML}{EADFD2}
\definecolor{popgray}{HTML}{8A7F76}
\colorlet{popgrey}{popgray}
% Alias tolérants (le modèle nomme parfois les couleurs différemment)
\colorlet{popwhite}{white}
\colorlet{popblack}{popink}
\colorlet{popred}{poppink}
\colorlet{popyellow}{popgold}
% Teintes claires pour fonds de tableaux / zones
\colorlet{popblueL}{popblue!10!white}
\colorlet{poporangeL}{poporange!14!white}
\colorlet{popgreenL}{popgreen!12!white}
\colorlet{poppurpleL}{poppurple!12!white}
\colorlet{poppinkL}{poppink!12!white}
\colorlet{popgoldL}{popgold!14!white}

\pagecolor{popcream}

% ============================================================
% Typographie
% ============================================================
\defaultfontfeatures{Ligatures=TeX}
\setmainfont{PlusJakartaSans-Medium.ttf}[Path=../../../fonts/,BoldFont=PlusJakartaSans-Bold.ttf]
% Maths en sans-serif (Fira Math) avec chiffres et lettres droites de Plus
% Jakarta, pour que les formules se fondent dans le texte.
\usepackage[math-style=ISO,bold-style=ISO]{unicode-math}
\setmathfont{FiraMath-Regular.otf}[Path=../../../fonts/]
\setmathfont{PlusJakartaSans-Medium.ttf}[Path=../../../fonts/,range={up/num,up/{Latin,latin},\mathup}]
% (icomma retiré : décimales avec point en anglais)
% Caractères mathématiques Unicode absents des polices de texte (Plus Jakarta,
% Archivo Black) : rendus par leur équivalent mathématique, en texte comme en
% formule — sinon ils s'affichent en carré vide (ex. « Arithmétique dans ℤ »).
\usepackage{newunicodechar}
\newunicodechar{ℝ}{\ensuremath{\mathbb{R}}}
\newunicodechar{ℤ}{\ensuremath{\mathbb{Z}}}
\newunicodechar{ℂ}{\ensuremath{\mathbb{C}}}
\newunicodechar{ℕ}{\ensuremath{\mathbb{N}}}
\newunicodechar{ℚ}{\ensuremath{\mathbb{Q}}}
\newunicodechar{𝔻}{\ensuremath{\mathbb{D}}}
\newunicodechar{α}{\ensuremath{\alpha}}
\newunicodechar{β}{\ensuremath{\beta}}
\newunicodechar{γ}{\ensuremath{\gamma}}
\newunicodechar{δ}{\ensuremath{\delta}}
\newunicodechar{ε}{\ensuremath{\varepsilon}}
\newunicodechar{θ}{\ensuremath{\theta}}
\newunicodechar{λ}{\ensuremath{\lambda}}
\newunicodechar{μ}{\ensuremath{\mu}}
\newunicodechar{σ}{\ensuremath{\sigma}}
\newunicodechar{τ}{\ensuremath{\tau}}
\newunicodechar{φ}{\ensuremath{\varphi}}
\newunicodechar{ψ}{\ensuremath{\psi}}
\newunicodechar{ω}{\ensuremath{\omega}}
\newunicodechar{Σ}{\ensuremath{\Sigma}}
% majuscules produites par \MakeUppercase dans les titres (α-aminés -> Α)
\newunicodechar{Α}{\ensuremath{\alpha}}
\newunicodechar{Β}{\ensuremath{\beta}}
\newunicodechar{Γ}{\ensuremath{\Gamma}}
\newunicodechar{Δ}{\ensuremath{\Delta}}
\newunicodechar{Ε}{\ensuremath{\varepsilon}}
\newunicodechar{Θ}{\ensuremath{\Theta}}
\newunicodechar{Λ}{\ensuremath{\Lambda}}
\newunicodechar{Μ}{\ensuremath{\mu}}
\newunicodechar{Τ}{\ensuremath{\tau}}
\newunicodechar{Φ}{\ensuremath{\Phi}}
\newunicodechar{Ψ}{\ensuremath{\Psi}}
\newunicodechar{Ω}{\ensuremath{\Omega}}
\newunicodechar{↦}{\ensuremath{\mapsto}}
\newunicodechar{∈}{\ensuremath{\in}}
\newunicodechar{∉}{\ensuremath{\notin}}
\newunicodechar{∧}{\ensuremath{\wedge}}
\newunicodechar{⇔}{\ensuremath{\Leftrightarrow}}
\newunicodechar{⟺}{\ensuremath{\Longleftrightarrow}}
\newunicodechar{⇒}{\ensuremath{\Rightarrow}}
\newunicodechar{⟹}{\ensuremath{\Longrightarrow}}
\newunicodechar{∘}{\ensuremath{\circ}}
\newunicodechar{∪}{\ensuremath{\cup}}
\newunicodechar{∩}{\ensuremath{\cap}}
\newunicodechar{≡}{\ensuremath{\equiv}}
\newunicodechar{⊂}{\ensuremath{\subset}}
\newunicodechar{⊥}{\ensuremath{\perp}}
\newunicodechar{∃}{\ensuremath{\exists}}
\newunicodechar{∀}{\ensuremath{\forall}}
\newunicodechar{∛}{\ensuremath{\sqrt[3]{\,}}}
\newunicodechar{∜}{\ensuremath{\sqrt[4]{\,}}}
\newunicodechar{⃗}{\ensuremath{{}^{\to}}}
\newunicodechar{ⁿ}{\ifmmode^{n}\else\textsuperscript{\ensuremath{n}}\fi}
\newunicodechar{ˣ}{\ifmmode^{x}\else\textsuperscript{\ensuremath{x}}\fi}
\newunicodechar{ᵇ}{\ifmmode^{b}\else\textsuperscript{\ensuremath{b}}\fi}
\newunicodechar{ʳ}{\ifmmode^{r}\else\textsuperscript{\ensuremath{r}}\fi}
\newunicodechar{ᵘ}{\ifmmode^{u}\else\textsuperscript{\ensuremath{u}}\fi}
\newunicodechar{ʸ}{\ifmmode^{y}\else\textsuperscript{\ensuremath{y}}\fi}
\newunicodechar{ᵃ}{\ifmmode^{a}\else\textsuperscript{\ensuremath{a}}\fi}
\newunicodechar{ᵉ}{\ifmmode^{e}\else\textsuperscript{\ensuremath{e}}\fi}
\newunicodechar{ᵏ}{\ifmmode^{k}\else\textsuperscript{\ensuremath{k}}\fi}
\newunicodechar{ᵖ}{\ifmmode^{p}\else\textsuperscript{\ensuremath{p}}\fi}
\newunicodechar{ᴺ}{\ifmmode^{N}\else\textsuperscript{\ensuremath{N}}\fi}
\newunicodechar{ᶜ}{\ifmmode^{c}\else\textsuperscript{\ensuremath{c}}\fi}
\newunicodechar{ʰ}{\ifmmode^{h}\else\textsuperscript{\ensuremath{h}}\fi}
\newunicodechar{ᵠ}{\ifmmode^{q}\else\textsuperscript{\ensuremath{q}}\fi}
\newunicodechar{ₙ}{\ifmmode_{n}\else\textsubscript{\ensuremath{n}}\fi}
\newunicodechar{ₐ}{\ifmmode_{a}\else\textsubscript{\ensuremath{a}}\fi}
\newunicodechar{ᵢ}{\ifmmode_{i}\else\textsubscript{\ensuremath{i}}\fi}
\newunicodechar{ᵣ}{\ifmmode_{r}\else\textsubscript{\ensuremath{r}}\fi}
\newunicodechar{ₖ}{\ifmmode_{k}\else\textsubscript{\ensuremath{k}}\fi}
\newunicodechar{ᵦ}{\ifmmode_{\beta}\else\textsubscript{\ensuremath{\beta}}\fi}
\newunicodechar{ₓ}{\ifmmode_{x}\else\textsubscript{\ensuremath{x}}\fi}
\newfontfamily\popdisplay{ArchivoBlack-Regular.ttf}[Path=../../../fonts/]
\newfontfamily\poplogo{SpaceGrotesk-Bold.ttf}[Path=../../../fonts/]
\newfontfamily\popsemi{PlusJakartaSans-SemiBold.ttf}[Path=../../../fonts/]
\newfontfamily\popxbold{PlusJakartaSans-ExtraBold.ttf}[Path=../../../fonts/]
\newfontfamily\popxboldls{PlusJakartaSans-ExtraBold.ttf}[Path=../../../fonts/,LetterSpace=3.0]

\setlist[enumerate]{leftmargin=1.7em,itemsep=5pt,topsep=4pt}
\setlist[itemize]{leftmargin=1.7em,itemsep=4pt,topsep=4pt,label={\color{poporange}\textbullet}}
\setlength{\parindent}{0pt}
\setlength{\parskip}{5pt}
\renewcommand{\arraystretch}{1.25}

% pgfplots : style Pop par défaut
\pgfplotsset{
  every axis/.append style={axis line style={popink,line width=1pt},tick style={popink},
    grid style={popline,line width=0.5pt},label style={font=\small},tick label style={font=\footnotesize}},
  popcurve/.style={poporange,line width=1.8pt,no markers,smooth},
}
% Même style disponible dans un \draw TikZ simple (hors pgfplots)
\tikzset{popcurve/.style={poporange,line width=1.8pt}}
\tikzset{popgrid/.style={popline,very thin}}
\colorlet{popcurve}{poporange} % le nom du style est parfois utilisé comme couleur

% ============================================================
% Pill / squiggle
% ============================================================
\newcommand{\poppill}[3]{%
  \tikz[baseline=-0.55ex]\node[draw=popink,line width=1.2pt,rounded corners=2.6pt,%
    fill=#2,inner xsep=6.5pt,inner ysep=3pt]{%
    \popxboldls\fontsize{7}{7}\selectfont\color{#3}\MakeUppercase{#1}};%
}
\newcommand{\popsquiggle}[1]{%
  \tikz[baseline=-0.4ex]{
    \draw[#1,line width=2.6pt,line cap=round]
      plot [smooth] coordinates {(0,0.09) (0.34,0.01) (0.68,0.21) (1.02,0.01) (1.36,0.10)};
  }%
}
% Mot-clé surligné (marqueur orange sous le texte)
\newcommand{\cle}[1]{{\popxbold\color{popdark}#1}}

% ============================================================
% En-tête / pied
% ============================================================
\pagestyle{fancy}
\fancyhf{}
\renewcommand{\headrulewidth}{0pt}
\renewcommand{\footrulewidth}{0pt}
\lhead{\raisebox{-0.15ex}{\poplogo\fontsize{13}{13}\selectfont\color{popink}OnBuch\color{poporange}+}}
\rhead{\poppill{\DOCMATIERE\ \textbullet\ \DOCNIVEAU\ \textbullet\ Chapter \DOCLECON}{white}{popink}}
\lfoot{\popxbold\fontsize{7.6}{7.6}\selectfont\color{popmuted}\MakeUppercase{Signed course OnBuch+}\ \textbullet\ {\popsemi\DOCTITRE}}
\rfoot{\tikz[baseline=-0.6ex]\node[rounded corners=8.5pt,fill=popink,minimum height=17pt,minimum width=17pt,inner xsep=5pt,inner sep=0]{\popxbold\fontsize{8}{8}\selectfont\color{popcream}\thepage\,/\,\pageref*{LastPage}};}

% ============================================================
% Titres
% ============================================================
\newcommand{\chaptitre}[2]{%
  \begin{tcolorbox}[enhanced,colback=poporange,colframe=popink,boxrule=2.6pt,arc=11pt,
      drop shadow=popink,left=15pt,right=15pt,top=12pt,bottom=11pt,before skip=3pt,after skip=18pt]
    \poppill{#2}{popink}{popcream}\par\vspace{9pt}
    {\raggedright\hyphenpenalty=10000\exhyphenpenalty=10000\popdisplay\fontsize{22}{24}\selectfont\color{popcream}\MakeUppercase{#1}\par}
  \end{tcolorbox}
}
% Section numérotée : pastille orange avec le numéro + titre Archivo
\newcounter{popsec}
\newcommand{\coursec}[1]{%
  \stepcounter{popsec}\setcounter{popsub}{0}%
  \par\vspace{16pt}\noindent
  \begin{minipage}{\linewidth}
  \tikz[baseline=-0.7ex]\node[draw=popink,line width=1.6pt,fill=poporange,rounded corners=4pt,minimum size=22pt,inner sep=0]{\popdisplay\fontsize{12}{12}\selectfont\color{popcream}\thepopsec};%
  \hspace{8pt}{\popdisplay\fontsize{15}{17}\selectfont\color{popink}\MakeUppercase{#1}}\par
  \vspace{4pt}\noindent{\color{poporange}\rule{36pt}{3.6pt}}
  \end{minipage}\par\nopagebreak\vspace{8pt}
}
\newcounter{popsub}
\newcommand{\courssub}[1]{%
  \stepcounter{popsub}%
  \par\vspace{9pt}\noindent{\popxbold\fontsize{11.5}{13}\selectfont\color{popink}{\color{poporange}\thepopsec.\thepopsub}\hspace{6pt}#1}\par\nopagebreak\vspace{4pt}
}

% ============================================================
% Boîtes pédagogiques — même recette : fond blanc, bordure épaisse colorée,
% ombre dure encre, badge plein. Titre optionnel [#1].
% ============================================================
\tcbset{popbase/.style={breakable,enhanced,colback=white,boxrule=2.3pt,arc=9pt,
  shadow={3pt}{-3pt}{0pt}{popink},left=14pt,right=14pt,top=11pt,bottom=11pt,before skip=11pt,after skip=15pt,
  fontupper=\popsemi\fontsize{10.6}{14.5}\selectfont\color{popink},
  fontlower=\popsemi\fontsize{10.6}{14.5}\selectfont\color{popink}}}
% Coche dessinée (le glyphe ✓ n'existe pas dans les polices du document)
\newcommand{\popcheck}[1]{\tikz[baseline=-0.5ex]\draw[#1,line width=1.6pt,line cap=round,line join=round](-0.13,0.0)--(-0.04,-0.09)--(0.13,0.1);}
\providecommand{\coche}{\popcheck{popgreen}}
\providecommand{\resultat}[1]{\par\smallskip\noindent\fcolorbox{popgreen}{popgreenL}{\parbox{\dimexpr\linewidth-2\fboxsep-2\fboxrule\relax}{\textbf{Result.} #1}}\par\smallskip}
\newcommand{\popboxhead}[3]{\poppill{#1}{#2}{white}\ifx\relax#3\relax\else\hspace{6pt}{\popxbold\fontsize{10.6}{12}\selectfont\color{popink}#3}\fi\par\vspace{7pt}}

\newtcolorbox{aretenir}[1][]{popbase,colframe=popgreen,before upper={\popboxhead{Key points}{popgreen}{#1}}}
\newtcolorbox{exemplebox}[1][]{popbase,colframe=poporange,before upper={\popboxhead{Exemple}{poporange}{#1}}}
\newtcolorbox{definition}[1][]{popbase,colframe=popblue,before upper={\popboxhead{Definition}{popblue}{#1}}}
\newtcolorbox{propriete}[1][]{popbase,colframe=poppurple,before upper={\popboxhead{Property}{poppurple}{#1}}}
\newtcolorbox{methode}[1][]{popbase,colframe=popgold,before upper={\popboxhead{Method}{popgold}{#1}}}
\newtcolorbox{attention}[1][]{popbase,colframe=poppink,before upper={\popboxhead{Warning}{poppink}{#1}}}
\newtcolorbox{experience}[1][]{popbase,colframe=popblue,colback=popblueL,before upper={\popboxhead{Experiment}{popblue}{#1}}}
% « Le savais-tu ? » : carte violette pleine (contexte, culture, Cameroun)
\newtcolorbox{savaistu}[1][]{popbase,colback=poppurple,colframe=popink,
  fontupper=\popsemi\fontsize{10.4}{14.2}\selectfont\color{white},
  before upper={\poppill{Did you know?}{white}{poppurple}\ifx\relax#1\relax\else\hspace{6pt}{\popxbold\color{white}#1}\fi\par\vspace{7pt}}}
% Exercice résolu : énoncé en haut, \tcblower puis la solution sur fond crème
\newcounter{popexo}\newcounter{popexr}
\newtcolorbox{exoresolu}[1][]{popbase,colframe=poporange,colbacklower=popcream,
  segmentation style={popink,line width=1.2pt,dashed},
  before upper={\stepcounter{popexr}\popboxhead{Worked example \thepopexr}{poporange}{#1}},
  before lower={\poppill{Solution}{popink}{popcream}\par\vspace{6pt}}}
% Exercice (non résolu en place) — la correction est dans la partie Corrigés
\newtcolorbox{exercice}[1][]{popbase,colframe=popink,
  before upper={\stepcounter{popexo}\popboxhead{Exercise \thepopexo}{popink}{#1}}}
% Corrigé
\newtcolorbox{corrige}[1][]{popbase,colframe=popgreen,colback=popgreenL,
  before upper={\popboxhead{Solution}{popgreen}{#1}}}
% Objectifs / prérequis / bilan
\newtcolorbox{objectifs}{popbase,colframe=popink,colback=poporangeL,
  before upper={\popboxhead{Chapter objectives}{popink}{}}}
\newtcolorbox{prerequis}{popbase,colframe=popink,colback=popblueL,
  before upper={\popboxhead{Prerequisites}{popblue}{}}}
\newtcolorbox{bilan}{popbase,colframe=popink,colback=white,boxrule=2.8pt,
  before upper={\popboxhead{Fiche bilan}{poporange}{L'essentiel à connaître pour l'examen}}}
% Figure encadrée avec légende
\newtcolorbox{popfigure}[1][]{enhanced,breakable=false,colback=white,colframe=popline,boxrule=1.4pt,arc=8pt,
  left=10pt,right=10pt,top=10pt,bottom=8pt,before skip=10pt,after skip=12pt,halign=center,
  lower separated=false,halign lower=center,
  fontlower=\popsemi\fontsize{9}{11.5}\selectfont\color{popmuted},
  #1}
\newcommand{\legende}[1]{\par\vspace{6pt}{\popsemi\fontsize{9}{11.5}\selectfont\color{popmuted}#1}}

% ============================================================
% Signature OnBuch+
% ============================================================
\newcommand{\poplogoinline}[1]{{\poplogo\fontsize{#1}{#1}\selectfont\color{popink}OnBuch\color{poporange}+}}
\newcommand{\signatureonbuch}{%
  \begin{tcolorbox}[enhanced,colback=popink,colframe=popink,boxrule=2.2pt,arc=10pt,
      drop shadow=poporange,left=14pt,right=14pt,top=10pt,bottom=10pt,before skip=0pt,after skip=0pt]
    \tikz[baseline=-0.7ex]\node[circle,fill=popgreen,draw=popcream,line width=1.3pt,minimum size=22pt,inner sep=0]{\popcheck{white}};%
    \hspace{9pt}\begin{minipage}[c]{0.62\linewidth}
      {\popxbold\fontsize{12}{13}\selectfont\color{popcream}Signed\ \textbullet\ {\poplogo OnBuch\color{poporange}+}}\par\vspace{2pt}
      {\raggedright\popsemi\fontsize{8}{10.5}\selectfont\color{popline}\MakeUppercase{OnBuch+ teaching team}\\\MakeUppercase{Follows the official MINESEC syllabus}\par}
    \end{minipage}\hfill
    {\poplogo\fontsize{22}{22}\selectfont\color{popcream}OnBuch\color{poporange}+}
  \end{tcolorbox}%
}

% ============================================================
% Page de garde
% #1 titre · #2 matière · #3 niveau · #4 module · #5 n° leçon · #6 durée ·
% #7 description / objectifs (texte ou itemize)
% ============================================================
\newcommand{\pagegarde}[7]{%
  \thispagestyle{empty}
  \newgeometry{a4paper,margin=1.7cm,top=1.6cm,bottom=1.4cm}%
  \begin{tikzpicture}[remember picture,overlay]
    \fill[poporange!18] ([xshift=-3.2cm,yshift=-3.6cm]current page.north east) circle (4.6cm);
    \fill[poppurple!14] ([xshift=2.4cm,yshift=7.6cm]current page.south west) circle (3.8cm);
  \end{tikzpicture}%
  {\poplogo\fontsize{30}{30}\selectfont\color{popink}OnBuch\color{poporange}+}\hfill
  \raisebox{0.6ex}{\poppill{Signed course \textbullet\ Premium}{popink}{popcream}}\par
  \vspace{1pt}\popsquiggle{poppurple}\par
  \vspace{8pt}
  \begin{tcolorbox}[enhanced,colback=poporange,colframe=popink,boxrule=2.8pt,arc=13pt,
      drop shadow=popink,left=18pt,right=18pt,top=12pt,bottom=13pt,after skip=10pt]
    \poppill{#2 \textbullet\ #3}{popink}{popcream}\hspace{5pt}\poppill{#4}{white}{popink}\par\vspace{8pt}
    {\popdisplay\fontsize{14}{14}\selectfont\color{popink}CHAPTER #5}\par\vspace{4pt}
    {\raggedright\hyphenpenalty=10000\exhyphenpenalty=10000\popdisplay\fontsize{25}{27}\selectfont\color{popcream}\MakeUppercase{#1}\par}
    \vspace{8pt}
    {\popxbold\fontsize{9.5}{9.5}\selectfont\color{popink}\MakeUppercase{Suggested time: #6}}
  \end{tcolorbox}
  \begin{tcolorbox}[enhanced,colback=white,colframe=popink,boxrule=2.2pt,arc=10pt,
      drop shadow=popink,left=16pt,right=16pt,top=10pt,bottom=10pt,after skip=8pt]
    \poppill{In this chapter}{poppurple}{white}\par\vspace{5pt}
    \popsemi\fontsize{9.3}{12.6}\selectfont\color{popink}#7
  \end{tcolorbox}
  \noindent\begin{minipage}[t]{0.487\linewidth}
    \begin{tcolorbox}[enhanced,colback=popgreen,colframe=popink,boxrule=2.2pt,arc=10pt,
        drop shadow=popink,left=11pt,right=11pt,top=10pt,bottom=10pt,height=3.1cm,valign=center]
      \tcbox[colback=white,colframe=popink,boxrule=1.3pt,arc=5pt,boxsep=0pt,nobeforeafter,
        left=3pt,right=3pt,top=3pt,bottom=3pt,valign=center]{\qrcode[height=1.9cm]{https://play.google.com/store/apps/details?id=com.luvvix.onbuch&hl=en}}%
      \hspace{8pt}\begin{minipage}[c]{0.5\linewidth}
        {\popdisplay\fontsize{11}{12.5}\selectfont\color{white}\MakeUppercase{Download OnBuch}}\par\vspace{4pt}
        {\raggedright\popsemi\fontsize{8.4}{11}\selectfont\color{white}Free \textbullet\ Google Play\\AI tutor, past papers, results\par}
      \end{minipage}
    \end{tcolorbox}
  \end{minipage}\hfill
  \begin{minipage}[t]{0.487\linewidth}
    \begin{tcolorbox}[enhanced,colback=poppurple,colframe=popink,boxrule=2.2pt,arc=10pt,
        drop shadow=popink,left=11pt,right=11pt,top=10pt,bottom=10pt,height=3.1cm,valign=center]
      \tikz[baseline=-0.7ex]\node[circle,fill=white,draw=popink,line width=1.3pt,minimum size=1.6cm,inner sep=0]{\popdisplay\fontsize{24}{24}\selectfont\color{poppurple}+};%
      \hspace{8pt}\begin{minipage}[c]{0.6\linewidth}
        {\popdisplay\fontsize{11}{12.5}\selectfont\color{white}\MakeUppercase{Go OnBuch+}}\par\vspace{4pt}
        {\raggedright\popsemi\fontsize{8.4}{11}\selectfont\color{white}All signed courses, unlimited AI tutor, PDF sheets — OnBuch+ tab in the app\par}
      \end{minipage}
    \end{tcolorbox}
  \end{minipage}
  \vspace{8pt}
  \popcontenu
  \vfill
  \signatureonbuch
  \restoregeometry
  \newpage
}

% Rangée « Ce cours contient » (page de garde)
\newcommand{\poptile}[3]{% #1 couleur #2 gros texte #3 légende
  \begin{tcolorbox}[enhanced,colback=#1,colframe=popink,boxrule=2pt,arc=9pt,shadow={2.5pt}{-2.5pt}{0pt}{popink},
      left=6pt,right=6pt,top=9pt,bottom=9pt,halign=center,valign=center,height=1.75cm,width=\linewidth,nobeforeafter]
    {\popdisplay\fontsize{12.5}{13}\selectfont\color{white}\MakeUppercase{#2}}\par\vspace{3pt}
    {\centering\popsemi\fontsize{7.8}{9.5}\selectfont\color{white}#3\par}
  \end{tcolorbox}}
\newcommand{\popcontenu}{%
  \par\noindent{\popxboldls\fontsize{7.5}{7.5}\selectfont\color{popmuted}THIS SIGNED COURSE CONTAINS}\par\vspace{7pt}
  \noindent\begin{minipage}[t]{0.235\linewidth}\poptile{popblue}{Course}{complete, illustrated, step by step}\end{minipage}\hfill
  \begin{minipage}[t]{0.235\linewidth}\poptile{poporange}{Methods}{and worked examples}\end{minipage}\hfill
  \begin{minipage}[t]{0.235\linewidth}\poptile{poppink}{Exercises}{exam style + solutions}\end{minipage}\hfill
  \begin{minipage}[t]{0.235\linewidth}\poptile{popgreen}{Summary sheet}{the essentials to remember}\end{minipage}\par}

% Sommaire
\newtcolorbox{sommaire}{enhanced,colback=white,colframe=popink,boxrule=2.2pt,arc=10pt,
  drop shadow=popink,left=18pt,right=18pt,top=13pt,bottom=13pt,before skip=6pt,after skip=20pt,
  before upper={\poppill{Contents}{poppurple}{white}\par\vspace{9pt}\popsemi\fontsize{10.6}{14.5}\selectfont\color{popink}}}

% Signature de fin de cours — poussée tout en bas de la dernière page
\newcommand{\findecours}{%
  \par\vfill\nopagebreak
  \begin{center}\popsquiggle{poporange}\end{center}\vspace{4pt}
  \signatureonbuch
}
\AtBeginDocument{\def\llbracket{\mathopen{[\mkern-3.5mu[}}\def\rrbracket{\mathclose{]\mkern-3.5mu]}}} % intervalles d'entiers (glyphe ⟦ absent de la police)
% Glyphes absents de Fira Math : redéfinis à partir de glyphes disponibles
\AtBeginDocument{%
  \def\setminus{\mathbin{\backslash}}\let\smallsetminus\setminus
  \def\vdots{\mathord{\vbox{\baselineskip3.2pt\lineskiplimit0pt\kern4pt\hbox{.}\hbox{.}\hbox{.}}}}%
  \def\ddots{\mathinner{\mkern1mu\raise6.5pt\hbox{.}\mkern2mu\raise3.6pt\hbox{.}\mkern2mu\raise0.7pt\hbox{.}\mkern1mu}}%
  \def\triangle{\mathord{\scalebox{0.9}{$\symup{\Delta}$}}}%
  \def\nabla{\mathord{\raisebox{\depth}{\scalebox{1}[-1]{$\symup{\Delta}$}}}}%
  \def\checkmark{\popcheck{popdark}}%
  \def\star{\mathbin{\ast}}\def\bigstar{\mathord{\ast}}%
  \def\diamond{\mathbin{\scalebox{0.6}{$\Diamond$}}}%
  \def\bigcup{\mathop{\mathchoice{\raisebox{-0.3em}{\scalebox{1.7}{$\cup$}}}{\scalebox{1.25}{$\cup$}}{\cup}{\cup}}\displaylimits}%
  \def\bigcap{\mathop{\mathchoice{\raisebox{-0.3em}{\scalebox{1.7}{$\cap$}}}{\scalebox{1.25}{$\cap$}}{\cap}{\cap}}\displaylimits}%
  \def\lesseqqgtr{\mathrel{\lessgtr}}\def\bigoplus{\mathop{\oplus}\displaylimits}%
}
\providecommand{\ssi}{if and only if} % abréviation fréquente
\AtBeginDocument{\providecommand{\wideparen}{\overparen}} % arc de cercle

% Fonctions usuelles que les modèles utilisent sans que amsmath les fournisse
\providecommand{\arsinh}{\operatorname{arsinh}}\providecommand{\arcosh}{\operatorname{arcosh}}
\providecommand{\artanh}{\operatorname{artanh}}\providecommand{\arsech}{\operatorname{arsech}}
\providecommand{\arcsch}{\operatorname{arcsch}}\providecommand{\arcoth}{\operatorname{arcoth}}
\providecommand{\arcsinh}{\operatorname{arsinh}}\providecommand{\arccosh}{\operatorname{arcosh}}
\providecommand{\arctanh}{\operatorname{artanh}}\providecommand{\sech}{\operatorname{sech}}
\providecommand{\csch}{\operatorname{csch}}\providecommand{\arcsec}{\operatorname{arcsec}}
\providecommand{\arccsc}{\operatorname{arccsc}}\providecommand{\arccot}{\operatorname{arccot}}
\providecommand{\sgn}{\operatorname{sgn}}\providecommand{\lcm}{\operatorname{lcm}}
\providecommand{\Var}{\operatorname{Var}}\providecommand{\Cov}{\operatorname{Cov}}

%% FIGFIX-BEGIN v4 — correctifs globaux des figures (docs/onbuch-generation/tools/figfix)
\usepackage{adjustbox}\usepackage{etoolbox}
% toute figure TikZ/pgfplots plus large que la ligne est réduite à la largeur disponible
\BeforeBeginEnvironment{tikzpicture}{\begin{adjustbox}{max width=\linewidth}}
\AfterEndEnvironment{tikzpicture}{\end{adjustbox}}
% légendes pgfplots lisibles par-dessus les courbes
\pgfplotsset{every axis/.append style={legend style={fill=white,fill opacity=0.92,text opacity=1,draw=black!15,inner sep=3pt}}}
% étiquettes d'arêtes (syntaxe edge["..."]) avec fond blanc
\tikzset{every edge quotes/.append style={fill=white,inner sep=1.2pt}}
%% FIGFIX-END
\begin{document}

\pagegarde{Further Curve Sketching and Inequalities}{Further Mathematics}{Upper Sixth}{Upper Sixth — Further mathematics}{7}{12 periods}{This chapter completes the study of functions begun in Lower Sixth: domains, limits, asymptotes, stationary points, concavity and symmetry are combined into a systematic method for sketching curves. It also develops algebraic and graphical techniques for solving absolute-value inequalities, a favourite GCE Advanced Level topic.
\vspace{4pt}
\begin{multicols}{2}\small
\begin{itemize}
\item By the end of this chapter I can determine the domain of definition of any algebraic, rational, irrational or composite function.
\item By the end of this chapter I can compute limits at the bounds of the domain, including infinite limits and limits at infinity.
\item By the end of this chapter I can find the range of values of a parameter for which a function (or equation) exists or has real solutions.
\item By the end of this chapter I can identify vertical, horizontal and oblique asymptotes and the intercepts of a curve.
\item By the end of this chapter I can locate stationary points and determine their nature using the first or second derivative.
\item By the end of this chapter I can construct a complete table of variation from the sign of the derivative.
\end{itemize}\end{multicols}}

\begin{sommaire}
\begin{multicols}{2}
\begin{enumerate}
\item Domain of Definition and Limits at the Bounds
\item Asymptotes, Intercepts and Infinite Branches
\item Stationary Points and Variation of a Function
\item Concavity, Points of Inflexion and Symmetry
\item Algebraic Solutions of Absolute-Value Inequalities
\item Graphical Solutions of Inequalities and Complete Curve Sketching
\item Integration activity
\item Methods and worked examples
\item Exercises
\item Solutions to the exercises
\item Summary sheet
\end{enumerate}
\end{multicols}
\end{sommaire}

\begin{objectifs}
\begin{itemize}
\item By the end of this chapter I can determine the domain of definition of any algebraic, rational, irrational or composite function.
\item By the end of this chapter I can compute limits at the bounds of the domain, including infinite limits and limits at infinity.
\item By the end of this chapter I can find the range of values of a parameter for which a function (or equation) exists or has real solutions.
\item By the end of this chapter I can identify vertical, horizontal and oblique asymptotes and the intercepts of a curve.
\item By the end of this chapter I can locate stationary points and determine their nature using the first or second derivative.
\item By the end of this chapter I can construct a complete table of variation from the sign of the derivative.
\item By the end of this chapter I can study the position of a curve relative to a line and describe infinite branches.
\item By the end of this chapter I can determine concavity and points of inflexion using the second derivative.
\item By the end of this chapter I can solve absolute-value inequalities both algebraically and graphically.
\item By the end of this chapter I can detect centres and axes of symmetry and produce a full, justified sketch of a curve.
\end{itemize}
\end{objectifs}

\begin{prerequis}
\begin{itemize}
\item Differentiation of polynomial, rational and composite functions (Lower Sixth).
\item Basic limits and continuity of a function at a point (Lower Sixth).
\item Solving linear and quadratic equations and inequalities.
\item The modulus (absolute value) function and its graph.
\item Equation of a straight line and coordinate geometry basics.
\end{itemize}
\end{prerequis}

\newpage

\coursec{Domain of Definition and Limits at the Bounds}

Imagine a mobile money agent in Bamenda. Her daily profit, modelled by a rational function of the number of transactions she processes, rises steeply at first, then levels off, and never exceeds a certain ceiling no matter how hard she works. Why does the profit curve flatten towards a horizontal line, and where exactly does it bend? How can she know for which numbers of transactions her profit stays above a fixed threshold --- that is, how can she solve an inequality involving absolute values of her daily balance? Curve sketching and inequalities give the complete answer: they let us predict the \emph{whole} behaviour of a function before plotting a single point. This chapter turns you into a mathematical detective, able to draw any reasonable curve and to read off solutions of inequalities from it.

Every investigation of a curve begins with two fundamental questions: \emph{Where is the function defined?} and \emph{What happens at the edges of that territory?} That is exactly the business of this section.

\courssub{The Domain of Definition of a Function}

\textbf{Intuition first.} A function is a machine: you feed it a number $x$, and it produces a number $f(x)$. But no machine accepts \emph{every} input. You cannot divide by zero, you cannot take the square root of a negative quantity (in $\mathbb{R}$), and you cannot take the logarithm of a negative number or of zero. The \cle{domain of definition} is simply the list of all inputs the machine will accept.

\begin{definition}[Domain of definition]
Let $f$ be a real-valued function of a real variable. The \cle{domain of definition} of $f$, written $D_f$, is the set of all real numbers $x$ for which $f(x)$ exists and is a real number:
\[ D_f = \{\, x \in \mathbb{R} \;:\; f(x) \text{ exists in } \mathbb{R} \,\}. \]
\end{definition}

\begin{propriete}[Standard domain conditions]
When searching for $D_f$, apply the following rules.
\begin{enumerate}
\item A \textbf{denominator} must be \textbf{non-zero}: in $\dfrac{g(x)}{h(x)}$, we require $h(x) \neq 0$.
\item The \textbf{radicand of an even root} must be \textbf{non-negative}: in $\sqrt{u(x)}$, we require $u(x) \geq 0$.
\item (\emph{Useful extra for the GCE}) The \textbf{argument of a logarithm} must be \textbf{strictly positive}: in $\ln\bigl(u(x)\bigr)$, we require $u(x) > 0$.
\end{enumerate}
\end{propriete}

\begin{attention}[Square root in a denominator]
A very common mistake: for $\dfrac{1}{\sqrt{u(x)}}$, students write only $u(x) \geq 0$. This is wrong! The radicand must be non-negative \emph{and} the denominator must be non-zero, so the correct condition is $u(x) > 0$ \textbf{strictly}. For example, $f(x) = \dfrac{1}{\sqrt{x-2}}$ has domain $]2, +\infty[$, not $[2, +\infty[$.
\end{attention}

\begin{methode}[Finding the domain of definition]
\begin{enumerate}
\item \textbf{List the constraints}: identify every denominator, every even root, every logarithm in the expression of $f$.
\item \textbf{Solve each constraint} separately (equations or inequalities in $x$).
\item \textbf{Intersect} the solution sets: $x$ must satisfy \emph{all} constraints simultaneously.
\item \textbf{Write the answer as a union of intervals}, using brackets $[\,]$ for included bounds and $]\,[$ for excluded bounds.
\end{enumerate}
\end{methode}

\textbf{Domains of combined functions.} If $f$ and $g$ have domains $D_f$ and $D_g$, then:
\begin{itemize}
\item the \textbf{sum} $f+g$ and the \textbf{product} $fg$ are defined on $D_f \cap D_g$;
\item the \textbf{quotient} $\dfrac{f}{g}$ is defined on $\{\, x \in D_f \cap D_g : g(x) \neq 0 \,\}$;
\item the \textbf{composite} $g \circ f$ is defined on $\{\, x \in D_f : f(x) \in D_g \,\}$: $x$ must be acceptable to $f$, and the output $f(x)$ must be acceptable to $g$.
\end{itemize}

\begin{exoresolu}[Domain of a mixed function]
Determine the domain of definition of $f(x) = \dfrac{\sqrt{x-1}}{x-3}$.
\tcblower
\textbf{Solution.} We list the constraints.
\begin{itemize}
\item Radicand: $x - 1 \geq 0 \iff x \geq 1$.
\item Denominator: $x - 3 \neq 0 \iff x \neq 3$.
\end{itemize}
Intersecting: $x \geq 1$ and $x \neq 3$, hence
\[ D_f = [1, 3[ \;\cup\; ]3, +\infty[. \]
Note that $x = 1$ \emph{is} included (the numerator may be zero), while $x = 3$ is excluded.
\end{exoresolu}

\begin{popfigure}
\begin{center}
\begin{tikzpicture}
\begin{axis}[
  width=12cm, height=7.5cm,
  axis lines=middle,
  xlabel={$x$}, ylabel={$y$},
  xmin=-0.5, xmax=8, ymin=-3, ymax=4,
  xtick={1,2,3,4,5,6,7}, ytick={-2,-1,1,2,3},
  domain=1:8, samples=200, clip=false
]
\addplot[popblue, very thick, domain=1:2.55, samples=100] {sqrt(max(0,x-1))/(x-3)};
\addplot[popblue, very thick, domain=3.45:8, samples=100] {sqrt(max(0,x-1))/(x-3)};
\addplot[poporange, dashed, thick] coordinates {(3,-3) (3,4)};
\addplot[popgreen, line width=3pt] coordinates {(1,0.15) (2.9,0.15)};
\addplot[popgreen, line width=3pt] coordinates {(3.1,0.15) (8,0.15)};
\node[popgreen, anchor=south] at (axis cs:2,0.18) {\small $D_f$};
\node[popgreen, anchor=south] at (axis cs:5.5,0.18) {\small $D_f$};
\node[poporange, anchor=west] at (axis cs:3.1,3.6) {\small $x=3$};
\addplot[only marks, mark=*, popblue] coordinates {(1,0)};
\end{axis}
\end{tikzpicture}
\end{center}
\legende{Graph of $f(x)=\dfrac{\sqrt{x-1}}{x-3}$. The domain $[1,3[\,\cup\,]3,+\infty[$ is highlighted in green on the $x$-axis; the dashed line $x=3$ is excluded from the domain.}
\end{popfigure}

\courssub{Limits at the Bounds of the Domain}

Once $D_f$ is known, the next detective question is: \emph{how does $f$ behave as $x$ approaches a bound of $D_f$?} The bounds are of two kinds: \textbf{finite bounds} (like $1$ and $3$ above) and the \textbf{infinite bounds} $-\infty$ and $+\infty$.

\begin{definition}[Finite limit at a point; one-sided limits]
We say that $f(x)$ tends to the real number $\ell$ as $x$ tends to $a$, and we write $\lim_{x \to a} f(x) = \ell$, if $f(x)$ becomes and stays arbitrarily close to $\ell$ whenever $x$ is close enough to $a$ (with $x \neq a$).
\begin{itemize}
\item The \cle{right-hand limit} $\lim_{x \to a^+} f(x)$ uses only values $x > a$.
\item The \cle{left-hand limit} $\lim_{x \to a^-} f(x)$ uses only values $x < a$.
\end{itemize}
The limit $\lim_{x \to a} f(x)$ exists if and only if both one-sided limits exist \textbf{and are equal}.
\end{definition}

\begin{attention}[Compare both sides before concluding]
A classic trap: students compute one side, find $+\infty$, and stop. You must \textbf{always compare the left-hand and right-hand limits}. If $\lim_{x \to a^-} f(x) \neq \lim_{x \to a^+} f(x)$, then $\lim_{x \to a} f(x)$ \textbf{does not exist}. Example: for $f(x) = \dfrac{1}{x-3}$, the left-hand limit is $-\infty$ and the right-hand limit is $+\infty$; the two-sided limit does not exist, even though each side has a perfectly good infinite limit.
\end{attention}

\begin{propriete}[Infinite limits and vertical asymptotes]
If $\lim_{x \to a^+} f(x) = \pm\infty$ or $\lim_{x \to a^-} f(x) = \pm\infty$, then the line $x = a$ is a \cle{vertical asymptote} of the curve of $f$. This typically occurs when the denominator vanishes at $a$ while the numerator does not.
\end{propriete}

\begin{popfigure}
\begin{center}
\begin{tikzpicture}[scale=1.1]
\draw[->, gray] (-0.5,0) -- (6,0) node[right] {\small $x$};
\draw[->, gray] (0,-1.5) -- (0,3.5) node[above] {\small $y$};
\draw[poporange, dashed, thick] (3,-1.5) -- (3,3.5);
\node[poporange, below left] at (3,0) {\small $a$};
\draw[popblue, very thick, domain=3.35:5.6, samples=80] plot (\x, {1/(\x-3)});
\draw[popblue, very thick, domain=0.4:2.65, samples=80] plot (\x, {1/(\x-3)});
\draw[->, popdark, thick] (5.2,0.5) -- (4.1,1.15);
\node[popdark, right, align=center, text width=3.2cm] at (5.3,0.6) {\small right-hand limit $x \to a^{+}$: $+\infty$};
\draw[->, popdark, thick] (0.8,-0.6) -- (1.9,-1.1);
\node[popdark, left, align=center, text width=3.2cm] at (0.7,-0.7) {\small left-hand limit $x \to a^{-}$: $-\infty$};
\end{tikzpicture}
\end{center}
\legende{Left-hand and right-hand limits at a bound $a$: the arrows show the direction of approach along each branch of the curve.}
\end{popfigure}

\begin{exoresolu}[One-sided limits at a finite bound]
For $f(x) = \dfrac{\sqrt{x-1}}{x-3}$, compute the limits at the bounds $x = 1$ and $x = 3$.
\tcblower
\textbf{Solution.}
\textbf{At $x = 1$:} only the right-hand limit makes sense (the domain starts at $1$). As $x \to 1^+$, $\sqrt{x-1} \to 0$ and $x - 3 \to -2 \neq 0$, so
\[ \lim_{x \to 1^+} f(x) = \frac{0}{-2} = 0. \]
The curve starts at the point $(1, 0)$.
\textbf{At $x = 3$:} the numerator tends to $\sqrt{2} > 0$. For the denominator, use a sign table of $x - 3$: it is negative for $x < 3$ and positive for $x > 3$. Hence
\[ \lim_{x \to 3^-} f(x) = \frac{\sqrt{2}}{0^-} = -\infty, \qquad \lim_{x \to 3^+} f(x) = \frac{\sqrt{2}}{0^+} = +\infty. \]
The line $x = 3$ is a vertical asymptote, and the two-sided limit at $3$ does not exist.
\end{exoresolu}

\courssub{Limits at Infinity: Polynomials, Rational Functions, Radicals}

\textbf{Polynomials.} At $\pm\infty$, a polynomial behaves like its \emph{highest-degree term}, because that term eventually dwarfs all the others:
\[ \lim_{x \to +\infty} \bigl(2x^3 - 5x^2 + 7\bigr) = \lim_{x \to +\infty} 2x^3 = +\infty. \]

\begin{propriete}[Rational functions at infinity]
Let $f(x) = \dfrac{P(x)}{Q(x)}$ where $P$ and $Q$ are polynomials. Then at $+\infty$ or $-\infty$, the limit of $f$ equals the limit of the \textbf{ratio of the dominant (highest-degree) terms}:
\[ \lim_{x \to \pm\infty} \frac{a_n x^n + \cdots}{b_m x^m + \cdots} = \lim_{x \to \pm\infty} \frac{a_n x^n}{b_m x^m}. \]
Consequently the limit is: $0$ if $\deg P < \deg Q$; the ratio $\dfrac{a_n}{b_m}$ of leading coefficients if $\deg P = \deg Q$; and $\pm\infty$ (with sign to be determined) if $\deg P > \deg Q$. (\emph{Proof by factorising the dominant term --- see the method below; it is instructive and frequently demanded in examinations.})
\end{propriete}

\textbf{Why factorisation works.} Take $f(x) = \dfrac{3x^2 - x + 1}{2x^2 + 5}$. For $x \neq 0$, factor the dominant term top and bottom:
\[ f(x) = \frac{x^2\left(3 - \dfrac{1}{x} + \dfrac{1}{x^2}\right)}{x^2\left(2 + \dfrac{5}{x^2}\right)} = \frac{3 - \dfrac{1}{x} + \dfrac{1}{x^2}}{2 + \dfrac{5}{x^2}} \xrightarrow[x \to \pm\infty]{} \frac{3}{2}, \]
since $\dfrac{1}{x} \to 0$ and $\dfrac{1}{x^2} \to 0$. This is exactly the mobile money agent's ``ceiling'': a finite limit $\ell$ at infinity means the curve flattens towards the horizontal line $y = \ell$.

\begin{methode}[Removing an indeterminate form]
The \cle{indeterminate forms} are $\dfrac{0}{0}$, $\dfrac{\infty}{\infty}$, $0 \times \infty$, $\infty - \infty$ (and powers such as $1^{\infty}$). They mean ``the rules of limits do not apply directly --- transform the expression''.
\begin{enumerate}
\item \textbf{Form $\dfrac{\infty}{\infty}$ at infinity} (polynomials, rational functions): \emph{factorise the dominant term} in the numerator and in the denominator, simplify, then pass to the limit.
\item \textbf{Form $\dfrac{0}{0}$ at a finite point $a$}: \emph{factorise $(x - a)$} in numerator and denominator and cancel.
\item \textbf{Form $\infty - \infty$ with radicals}: \emph{multiply by the conjugate}: use $\sqrt{A} - \sqrt{B} = \dfrac{A - B}{\sqrt{A} + \sqrt{B}}$.
\end{enumerate}
\end{methode}

\begin{exoresolu}[Conjugate multiplication]
Compute $\lim_{x \to +\infty} \bigl(\sqrt{x^2 + 3x} - x\bigr)$.
\tcblower
\textbf{Solution.} This has the indeterminate form $\infty - \infty$. Multiply by the conjugate:
\begin{align*}
\sqrt{x^2 + 3x} - x &= \frac{\bigl(\sqrt{x^2+3x} - x\bigr)\bigl(\sqrt{x^2+3x} + x\bigr)}{\sqrt{x^2+3x} + x} = \frac{x^2 + 3x - x^2}{\sqrt{x^2+3x} + x} = \frac{3x}{\sqrt{x^2+3x} + x}.
\end{align*}
For $x > 0$, $\sqrt{x^2 + 3x} = x\sqrt{1 + \tfrac{3}{x}}$, so
\[ \frac{3x}{x\sqrt{1 + \tfrac{3}{x}} + x} = \frac{3}{\sqrt{1 + \tfrac{3}{x}} + 1} \xrightarrow[x \to +\infty]{} \frac{3}{2}. \]
Hence $\lim_{x \to +\infty} \bigl(\sqrt{x^2 + 3x} - x\bigr) = \dfrac{3}{2}$.
\end{exoresolu}

\begin{attention}[Radicals and the sign of $x$]
When $x \to -\infty$, $\sqrt{x^2} = |x| = -x$, \textbf{not} $x$. Forgetting this changes signs and ruins the limit. Always write $\sqrt{x^2} = |x|$ first, then discuss the sign of $x$.
\end{attention}

\courssub{Range of Values of a Parameter for Existence of a Function}

Many GCE questions ask: \emph{for which values of the parameter $m$ is the function defined for all real $x$} (or: \emph{on a given interval})? The key tool is the sign of a quadratic, controlled by its discriminant.

\begin{propriete}[Sign of a quadratic and the discriminant]
Let $u(x) = ax^2 + bx + c$ with $a \neq 0$ and discriminant $\Delta = b^2 - 4ac$.
\begin{itemize}
\item If $\Delta < 0$, then $u(x)$ keeps the \textbf{constant sign of $a$} on $\mathbb{R}$.
\item If $\Delta = 0$, then $u(x)$ has the sign of $a$ and vanishes once (at $x = -\tfrac{b}{2a}$).
\item If $\Delta > 0$, then $u(x)$ has the sign of $a$ outside the roots and the opposite sign between them.
\end{itemize}
In particular, $u(x) \geq 0$ for \textbf{all} real $x$ if and only if $a > 0$ and $\Delta \leq 0$.
\end{propriete}

\begin{popfigure}
\begin{center}
\begin{tikzpicture}
\begin{axis}[
  width=12cm, height=5cm,
  axis lines=middle,
  xlabel={$x$}, ylabel={$u(x)$},
  xmin=-6, xmax=6, ymin=-2, ymax=6,
  xtick={-4,-2,0,2,4}, ytick={1,2,3,4,5},
  samples=100
]
\addplot[popblue, very thick, domain=-6:6] {x^2 - 4*x + 4}; % Delta = 0
\addplot[poporange, very thick, domain=-6:6] {x^2 + 2*x + 5}; % Delta < 0
\addplot[popgreen, very thick, domain=-6:6] {x^2 - 2*x - 3}; % Delta > 0
\node[popblue, anchor=south] at (axis cs:2,0.5) {\small $\Delta=0$};
\node[poporange, anchor=south] at (axis cs:-1,4.5) {\small $\Delta<0$};
\node[popgreen, anchor=north] at (axis cs:3,-1.5) {\small $\Delta>0$};
\end{axis}
\end{tikzpicture}
\end{center}
\legende{Sign of a quadratic according to its discriminant $\Delta$.}
\end{popfigure}

\begin{exoresolu}[Parameter condition for existence on $\mathbb{R}$]
Find all values of the real parameter $m$ for which $f(x) = \sqrt{x^2 + mx + 4}$ is defined for every real number $x$.
\tcblower
\textbf{Solution.} We need the radicand $u(x) = x^2 + mx + 4$ to satisfy $u(x) \geq 0$ for all $x \in \mathbb{R}$. Here $a = 1 > 0$, so by the property above it is necessary and sufficient that $\Delta \leq 0$:
\[ \Delta = m^2 - 16 \leq 0 \iff m^2 \leq 16 \iff -4 \leq m \leq 4. \]
Hence $f$ is defined on the whole of $\mathbb{R}$ exactly when $m \in [-4, 4]$. For $|m| > 4$, the radicand becomes negative between the two roots of $u$, and the domain of $f$ is only a union of two unbounded intervals: the function \emph{fails to exist} on a whole central band.

\textbf{Sign table for $u(x)=x^2+mx+4$ when $m=5$ ($\Delta>0$):}
\begin{center}
\begin{tabular}{|c|ccccccccc|}\hline
$x$ & $-\infty$ & & $x_1$ & & $x_2$ & & $+\infty$ \\ \hline
$u(x)$ & $+$ & $\searrow$ & $0$ & $\searrow$ & $-$ & $\nearrow$ & $0$ & $\nearrow$ & $+$ \\ \hline
\end{tabular}
\end{center}
where $x_1, x_2$ are the roots. The function is undefined on $]x_1, x_2[$.
\end{exoresolu}

\begin{exercice}
\begin{enumerate}
\item Determine the domain of definition of $g(x) = \dfrac{\sqrt{4 - x^2}}{x + 1}$ and of $h(x) = \ln\!\left(\dfrac{x-1}{x+2}\right)$.
\item Compute $\lim_{x \to 2^-} \dfrac{x+1}{x^2 - 4}$ and $\lim_{x \to 2^+} \dfrac{x+1}{x^2 - 4}$. Does the two-sided limit exist?
\item Compute $\lim_{x \to -\infty} \dfrac{2x^2 - 3x}{x^2 + 1}$ and $\lim_{x \to +\infty} \bigl(\sqrt{x^2 + x} - x\bigr)$.
\item Find the values of $k$ for which $f(x) = \dfrac{1}{\sqrt{kx^2 - 4x + k}}$ is defined for every real $x$.
\end{enumerate}
\end{exercice}

\begin{corrige}[Exercise 1]
\textbf{1.} For $g$: $4 - x^2 \geq 0 \iff -2 \leq x \leq 2$, and $x \neq -1$. Hence $D_g = [-2, -1[ \,\cup\, ]-1, 2]$. For $h$: we need $\dfrac{x-1}{x+2} > 0$; a sign table gives $D_h = ]-\infty, -2[ \,\cup\, ]1, +\infty[$.

\textbf{2.} $x^2 - 4 = (x-2)(x+2)$. Near $x = 2$ the numerator tends to $3 > 0$. For $x \to 2^-$: $(x-2) < 0$, $(x+2) > 0$, so the denominator is $0^-$ and the limit is $-\infty$. For $x \to 2^+$: denominator $0^+$, limit $+\infty$. The two one-sided limits differ, so the two-sided limit \textbf{does not exist}; $x = 2$ is a vertical asymptote.

\textbf{3.} Ratio of dominant terms: $\dfrac{2x^2}{x^2} = 2$, so the limit is $2$ (note: the sign of $x$ does not matter here since the degrees are even and equal). For the radical: conjugate multiplication gives $\dfrac{x}{\sqrt{x^2+x} + x} = \dfrac{1}{\sqrt{1 + \frac{1}{x}} + 1} \to \dfrac{1}{2}$.

\textbf{4.} We need $kx^2 - 4x + k > 0$ for all $x$ (strict, because of the denominator). Conditions: $k > 0$ and $\Delta = 16 - 4k^2 < 0$, i.e. $k^2 > 4$. Combined with $k > 0$: $k > 2$.
\end{corrige}

\begin{aretenir}[Key points of this section]
\begin{itemize}
\item $D_f$ is found by listing constraints (denominator $\neq 0$; even radicand $\geq 0$; logarithm argument $> 0$), solving each, and intersecting; write the result as a union of intervals.
\item A two-sided limit exists only when the left-hand and right-hand limits exist and are equal.
\item An infinite one-sided limit at $x = a$ signals a vertical asymptote $x = a$.
\item At infinity, keep only the dominant terms; remove indeterminate forms by factorising the dominant term or multiplying by the conjugate.
\item A radicand $ax^2 + bx + c$ stays non-negative on $\mathbb{R}$ exactly when $a > 0$ and $\Delta \leq 0$ --- the standard tool for parameter questions on existence.
\end{itemize}
\end{aretenir}

\coursec{Asymptotes, Intercepts and Infinite Branches}

In the previous section we learnt how to determine the \cle{domain of definition} of a function and how to compute its \cle{limits at the bounds} of that domain. Those limits are not merely numbers: they carry a precise \emph{geometric} meaning. When a limit at a finite point is infinite, the curve climbs or dives without bound near a vertical line; when a limit at infinity is finite, the curve settles down along a horizontal line. This section turns those analytic facts into geometry: we define \cle{asymptotes} and \cle{intercepts}, we classify all possible \cle{infinite branches}, and we learn how to determine the \cle{position of a curve with respect to a line}. These are the tools that will make our curve sketching, at the end of the chapter, fast and reliable.

\courssub{Intercepts of a Curve}

\begin{definition}[Intercepts]
Let $f$ be a function and $\mathcal{C}_f$ its curve in an orthonormal frame $(O;\vec{\imath},\vec{\jmath})$.
\begin{itemize}
\item The \cle{$y$-intercept} is the point where $\mathcal{C}_f$ meets the $y$-axis. If $0$ belongs to the domain of $f$, it is the point $(0,\,f(0))$.
\item The \cle{$x$-intercepts} are the points where $\mathcal{C}_f$ meets the $x$-axis. Their abscissae are the solutions of the equation $f(x)=0$ \emph{that belong to the domain of $f$}.
\end{itemize}
\end{definition}

The intercepts are the first points we plot: they anchor the curve to the axes. Note that a curve has \emph{at most one} $y$-intercept (a function assigns a single image to $0$), but it may have \emph{several} $x$-intercepts, or none at all.

\begin{exemplebox}[Finding intercepts]
Let $f(x)=\dfrac{x^{2}-4}{x+1}$, defined on $\mathbb{R}\setminus\{-1\}$.
\begin{itemize}
\item $y$-intercept: $f(0)=\dfrac{-4}{1}=-4$, so the curve passes through $(0,-4)$.
\item $x$-intercepts: $f(x)=0 \iff x^{2}-4=0 \iff x=2$ or $x=-2$; both values are in the domain. The curve cuts the $x$-axis at $(-2,0)$ and $(2,0)$.
\end{itemize}
\end{exemplebox}

\begin{attention}[Do not solve $f(x)=0$ outside the domain]
A candidate $x$-intercept must belong to the domain of definition. For $g(x)=\dfrac{x^{2}-1}{x-1}$, the equation $g(x)=0$ gives $x=\pm 1$, but $x=1$ is excluded from the domain: only $(-1,0)$ is an intercept. (In fact $g(x)=x+1$ for $x\neq 1$: the curve is a straight line with a ``hole'' at $x=1$.)
\end{attention}

\courssub{Vertical and Horizontal Asymptotes}

Imagine driving along the curve of $f(x)=\dfrac{1}{x-2}$ towards $x=2$. The road does not stop; it becomes a wall that the curve approaches but never touches. That wall is a \cle{vertical asymptote}.

\begin{definition}[Vertical asymptote]
If $\displaystyle\lim_{x\to a^{-}} f(x)=\pm\infty$ or $\displaystyle\lim_{x\to a^{+}} f(x)=\pm\infty$, then the line $x=a$ is a \cle{vertical asymptote} to $\mathcal{C}_f$.
\end{definition}

In practice, vertical asymptotes occur at values excluded from the domain where the function ``blows up'' — typically where the denominator of a rational function vanishes while the numerator does not.

\begin{definition}[Horizontal asymptote]
If $\displaystyle\lim_{x\to+\infty} f(x)=b$ (a finite number), then the line $y=b$ is a \cle{horizontal asymptote} to $\mathcal{C}_f$ at $+\infty$. The same holds at $-\infty$.
\end{definition}

\begin{exemplebox}[Vertical and horizontal asymptotes]
Let $f(x)=\dfrac{2x+1}{x-3}$ on $\mathbb{R}\setminus\{3\}$.
\begin{itemize}
\item $\displaystyle\lim_{x\to 3^{+}}f(x)=+\infty$ and $\displaystyle\lim_{x\to 3^{-}}f(x)=-\infty$: the line $x=3$ is a vertical asymptote.
\item $\displaystyle\lim_{x\to\pm\infty}f(x)=\lim_{x\to\pm\infty}\dfrac{2x}{x}=2$: the line $y=2$ is a horizontal asymptote at both ends.
\end{itemize}
\end{exemplebox}

\begin{attention}[A horizontal asymptote can be crossed!]
A very common mistake at the GCE is to believe that the curve may \emph{never} touch its horizontal asymptote. This is false. The asymptote describes the behaviour \emph{at infinity}, not near the origin. For example, $f(x)=\dfrac{\sin x}{x}$ crosses its horizontal asymptote $y=0$ infinitely many times, yet $\lim_{x\to+\infty}f(x)=0$. A curve may cross a horizontal or oblique asymptote — only a \emph{vertical} asymptote can never be crossed (the function is not defined there).
\end{attention}

\courssub{Oblique Asymptotes}

Some curves, like $f(x)=\dfrac{x^{2}+1}{x-1}$, have no horizontal asymptote because $f(x)\to\pm\infty$ at infinity — yet they do not wander randomly: they straighten up along a slanting line. Such a line is an \cle{oblique asymptote}.

\begin{definition}[Oblique asymptote]
The line $y=mx+c$ (with $m\neq 0$) is an \cle{oblique asymptote} to $\mathcal{C}_f$ at $+\infty$ (resp.\ $-\infty$) if
\[ \lim_{x\to+\infty}\bigl[f(x)-(mx+c)\bigr]=0 \quad\text{(resp.\ at } -\infty\text{)}. \]
\end{definition}

How do we find $m$ and $c$? If $f(x)-(mx+c)\to 0$, then dividing by $x$ (which tends to infinity) gives $\dfrac{f(x)}{x}-m-\dfrac{c}{x}\to 0$, hence $\dfrac{f(x)}{x}\to m$. Once $m$ is known, $c$ is recovered from $f(x)-mx\to c$. This justifies the following method (the derivation above is instructive and may be asked as a justification).

\begin{methode}[Finding an oblique asymptote in two steps]
\begin{enumerate}
\item \textbf{Compute the slope:} $m=\displaystyle\lim_{x\to\pm\infty}\dfrac{f(x)}{x}$. If $m$ is finite and non-zero, continue. If $m=0$, look for a horizontal asymptote instead; if the limit is infinite, the branch is parabolic (see below).
\item \textbf{Compute the intercept:} $c=\displaystyle\lim_{x\to\pm\infty}\bigl[f(x)-mx\bigr]$. If $c$ is finite, the line $y=mx+c$ is an oblique asymptote. If this limit is infinite or fails to exist, there is no oblique asymptote (parabolic branch).
\end{enumerate}
Repeat the study at $+\infty$ and at $-\infty$ separately: the answers may differ.
\end{methode}

\begin{attention}[Compute BOTH $m$ and $c$]
A frequent error is to compute only $m=\lim f(x)/x$, find $m=1$, and declare ``$y=x$ is an asymptote''. This is wrong: $f(x)=x+5+\tfrac{1}{x}$ has $\lim f(x)/x=1$, but its asymptote is $y=x+5$, not $y=x$. The slope alone never determines the line. Always complete step 2.
\end{attention}

\begin{exoresolu}[Complete asymptote study of a rational function]
Let $f(x)=\dfrac{x^{2}+1}{x-1}$ and let $\mathcal{C}_f$ be its curve. Determine all asymptotes of $\mathcal{C}_f$.
\tcblower
\textbf{Domain:} $f$ is defined on $\mathbb{R}\setminus\{1\}$.

\textbf{Vertical asymptote.} $\lim_{x\to 1}(x^2+1)=2>0$ and $x-1\to 0$, so
\[ \lim_{x\to 1^{+}}f(x)=+\infty, \qquad \lim_{x\to 1^{-}}f(x)=-\infty. \]
The line $x=1$ is a vertical asymptote.

\textbf{Behaviour at infinity.} $\lim_{x\to\pm\infty}f(x)=\lim_{x\to\pm\infty}\dfrac{x^{2}}{x}=\pm\infty$: no horizontal asymptote. We search for an oblique one.

\emph{Step 1:} $\dfrac{f(x)}{x}=\dfrac{x^{2}+1}{x(x-1)}=\dfrac{x^{2}+1}{x^{2}-x}\xrightarrow[x\to\pm\infty]{}1$, so $m=1$.

\emph{Step 2:} $f(x)-x=\dfrac{x^{2}+1-x(x-1)}{x-1}=\dfrac{x+1}{x-1}\xrightarrow[x\to\pm\infty]{}1$, so $c=1$.

\textbf{Conclusion:} the line $y=x+1$ is an oblique asymptote at both $+\infty$ and $-\infty$. (Check: polynomial division gives $f(x)=x+1+\dfrac{2}{x-1}$, and $\dfrac{2}{x-1}\to 0$ — a fast alternative for rational functions, which also gives the position of the curve relative to the asymptote for free.)
\end{exoresolu}

\begin{popfigure}
\begin{center}
\begin{tikzpicture}
\begin{axis}[
  width=11cm, height=8.5cm,
  axis lines=middle, axis line style={->},
  xlabel={$x$}, ylabel={$y$},
  xmin=-5, xmax=7, ymin=-8, ymax=10,
  xtick={-4,-2,1,2,4,6}, ytick={-6,-4,-2,2,4,6,8},
  samples=200, domain=-5:7, clip=true,
  legend style={at={(0.02,0.98)}, anchor=north west, font=\small}
]
\addplot[popblue, thick, domain=-5:0.85] {(x^2+1)/(x-1)};
\addplot[popblue, thick, domain=1.15:7] {(x^2+1)/(x-1)};
\addlegendentry{$y=f(x)$}
\addplot[poporange, dashed, thick, domain=-5:7] {x+1};
\addlegendentry{$y=x+1$}
\draw[popgreen, dashed, thick] (axis cs:1,-8) -- (axis cs:1,10);
\node[popgreen, font=\small, rotate=90, anchor=south] at (axis cs:1.15,-6) {$x=1$};
\end{axis}
\end{tikzpicture}
\end{center}
\legende{The curve of $f(x)=\dfrac{x^{2}+1}{x-1}$ approaching its vertical asymptote $x=1$ and its oblique asymptote $y=x+1$.}
\end{popfigure}

\courssub{Infinite Branches: the Complete Classification}

An \cle{infinite branch} is a part of the curve that escapes to infinity. Classifying it means saying \emph{how} it escapes. Suppose $\lim_{x\to+\infty}f(x)=\pm\infty$ (otherwise we have a horizontal asymptote). We examine $\lim_{x\to+\infty}\dfrac{f(x)}{x}$:

\begin{center}
\begin{tabularx}{\linewidth}{|l|X|X|}
\hline
\rowcolor{popblueL} \textbf{Limits} & \textbf{Conclusion} & \textbf{Example} \\
\hline
$\lim f(x)=b$ (finite) & Horizontal asymptote $y=b$ & $f(x)=\dfrac{2x+1}{x-3}$ \\
\hline
$\lim f(x)=\pm\infty$, $\lim \dfrac{f(x)}{x}=m\neq 0$, $\lim[f(x)-mx]=c$ & Oblique asymptote $y=mx+c$ & $f(x)=\dfrac{x^{2}+1}{x-1}$ \\
\hline
$\lim f(x)=\pm\infty$, $\lim \dfrac{f(x)}{x}=\pm\infty$ & \cle{Parabolic branch} in the direction of the $y$-axis & $f(x)=x^{2}$ \\
\hline
$\lim f(x)=\pm\infty$, $\lim \dfrac{f(x)}{x}=0$ & \cle{Parabolic branch} in the direction of the $x$-axis & $f(x)=\sqrt{x}$ \\
\hline
\end{tabularx}
\end{center}

The fourth case of an infinite branch is the \emph{vertical asymptote} itself (infinite limit at a \emph{finite} point). A \cle{parabolic branch} means the curve bends away like a parabola: either it grows faster than any line ($f(x)/x\to\infty$, vertical direction, like $x^2$), or it grows slower than any line ($f(x)/x\to 0$, horizontal direction, like $\sqrt{x}$). In both cases the search for an oblique asymptote \emph{fails}, and that failure is precisely the conclusion to state in the examination.

\begin{popfigure}
\begin{center}
\begin{tikzpicture}[scale=0.85]
% Vertical asymptote
\begin{scope}[shift={(0,0)}]
\draw[->, gray] (-0.3,0) -- (2.6,0); \draw[->, gray] (0,-0.3) -- (0,2.6);
\draw[popgreen, dashed, thick] (1.6,-0.2) -- (1.6,2.5);
\draw[popblue, thick] plot[domain=0.15:1.35, samples=40] (\x, {0.5/(\x-1.6)+1.3});
\draw[popblue, thick] plot[domain=1.85:2.5, samples=20] (\x, {0.5/(\x-1.6)+1.3});
\node[align=center, font=\small] at (1.3,-0.9) {Vertical\\ asymptote};
\end{scope}
% Horizontal asymptote
\begin{scope}[shift={(4.2,0)}]
\draw[->, gray] (-0.3,0) -- (2.6,0); \draw[->, gray] (0,-0.3) -- (0,2.6);
\draw[poporange, dashed, thick] (-0.2,1.4) -- (2.5,1.4);
\draw[popblue, thick] plot[domain=0.1:2.5, samples=40] (\x, {1.4+0.9*exp(-1.5*\x)});
\node[align=center, font=\small] at (1.3,-0.9) {Horizontal\\ asymptote};
\end{scope}
% Oblique asymptote
\begin{scope}[shift={(8.4,0)}]
\draw[->, gray] (-0.3,0) -- (2.6,0); \draw[->, gray] (0,-0.3) -- (0,2.6);
\draw[poporange, dashed, thick] (-0.1,0.15) -- (2.5,2.45);
\draw[popblue, thick] plot[domain=0.1:2.4, samples=40] (\x, {0.9*\x+0.25+0.5*exp(-2*\x)});
\node[align=center, font=\small] at (1.3,-0.9) {Oblique\\ asymptote};
\end{scope}
% Parabolic branch
\begin{scope}[shift={(12.6,0)}]
\draw[->, gray] (-0.3,0) -- (2.6,0); \draw[->, gray] (0,-0.3) -- (0,2.6);
\draw[popblue, thick] plot[domain=0:1.55, samples=40] (\x, {\x*\x});
\node[align=center, font=\small] at (1.3,-0.9) {Parabolic\\ branch};
\end{scope}
\end{tikzpicture}
\end{center}
\legende{The four types of infinite branches: vertical asymptote, horizontal asymptote, oblique asymptote, parabolic branch.}
\end{popfigure}

\courssub{Position of a Curve with Respect to a Line}

Knowing the asymptote is only half the story: the examiner will also ask \emph{on which side} of it the curve lies. The idea is beautifully simple. Two points with the same abscissa $x$, one on the curve $(x,f(x))$ and one on the line $(x,g(x))$, are compared by the sign of the \cle{difference function} $d(x)=f(x)-g(x)$.

\begin{propriete}[Position of $\mathcal{C}_f$ relative to a line]
Let $\mathcal{C}_f$ be the curve of $f$ and $\Delta$ the line $y=g(x)=mx+c$, and let $d(x)=f(x)-g(x)$ on a common interval.
\begin{itemize}
\item If $d(x)>0$, then $\mathcal{C}_f$ is \textbf{above} $\Delta$.
\item If $d(x)<0$, then $\mathcal{C}_f$ is \textbf{below} $\Delta$.
\item If $d(x)=0$, the curve and the line \textbf{meet} at $(x,f(x))$.
\end{itemize}
This property is applied directly; its justification is the definition of the order on the ordinates of points with the same abscissa.
\end{propriete}

\begin{propriete}[Crossing an asymptote]
A curve crosses its horizontal or oblique asymptote $y=mx+c$ exactly at the points whose abscissae solve $f(x)=mx+c$, i.e.\ $d(x)=0$. If the sign of $d$ changes there, the curve passes from one side of the asymptote to the other.
\end{propriete}

\begin{methode}[Studying the position of a curve relative to a line]
\begin{enumerate}
\item Form the difference $d(x)=f(x)-g(x)$ and reduce it to a single fraction.
\item \textbf{Factorise} numerator and denominator as much as possible.
\item Build a \textbf{sign table} of $d(x)$ on the domain of $f$ (do not forget forbidden values).
\item Conclude in words: ``above on \dots'', ``below on \dots'', ``intersects the line at \dots''.
\end{enumerate}
\end{methode}

\begin{exoresolu}[Position of a curve relative to its oblique asymptote]
Let $f(x)=\dfrac{x^{2}+1}{x-1}$ with oblique asymptote $\Delta:\ y=x+1$. Study the position of $\mathcal{C}_f$ relative to $\Delta$.
\tcblower
We compute the difference:
\[ d(x)=f(x)-(x+1)=\frac{x^{2}+1}{x-1}-(x+1)=\frac{x^{2}+1-(x+1)(x-1)}{x-1}=\frac{x^{2}+1-x^{2}+1}{x-1}=\frac{2}{x-1}. \]
The numerator $2$ is positive, so $d(x)$ has the sign of $x-1$:
\begin{center}
\begin{tabular}{|c|ccccc|}
\hline
$x$ & $-\infty$ & & $1$ & & $+\infty$ \\
\hline
$d(x)=\dfrac{2}{x-1}$ & & $-$ & $\|$ & $+$ & \\
\hline
\end{tabular}
\end{center}
\textbf{Conclusion:}
\begin{itemize}
\item On $\left]-\infty,1\right[$, $d(x)<0$: $\mathcal{C}_f$ is \textbf{below} $\Delta$.
\item On $\left]1,+\infty\right[$, $d(x)>0$: $\mathcal{C}_f$ is \textbf{above} $\Delta$.
\item $d(x)=0$ has no solution: the curve \textbf{never crosses} its oblique asymptote.
\end{itemize}
This matches the figure above: the left branch creeps up under the line $y=x+1$, the right branch descends onto it from above.
\end{exoresolu}

\begin{exoresolu}[Position relative to a given line, with crossing]
Let $f(x)=\dfrac{x^{2}-3x+3}{x-2}$ on $\mathbb{R}\setminus\{2\}$. Show that $\Delta:\ y=x-1$ is an oblique asymptote and study the position of $\mathcal{C}_f$ relative to $\Delta$.
\tcblower
\textbf{Asymptote.} $\dfrac{f(x)}{x}=\dfrac{x^{2}-3x+3}{x^{2}-2x}\to 1=m$; then
\[ f(x)-x=\frac{x^{2}-3x+3-x(x-2)}{x-2}=\frac{-x+3}{x-2}\to -1=c. \]
So $y=x-1$ is an oblique asymptote at $\pm\infty$.

\textbf{Position.} $d(x)=f(x)-(x-1)=\dfrac{x^{2}-3x+3-(x-1)(x-2)}{x-2}=\dfrac{1}{x-2}$.

\begin{center}
\begin{tabular}{|c|ccccc|}
\hline
$x$ & $-\infty$ & & $2$ & & $+\infty$ \\
\hline
$d(x)$ & & $-$ & $\|$ & $+$ & \\
\hline
\end{tabular}
\end{center}
$\mathcal{C}_f$ is below $\Delta$ on $\left]-\infty,2\right[$ and above on $\left]2,+\infty\right[$; since $d$ never vanishes, there is no intersection.

\textbf{A case with crossing.} Take $h(x)=\dfrac{x^{2}+4x+2}{x+1}$ on $\mathbb{R}\setminus\{-1\}$. Polynomial division gives $h(x)=x+2+\dfrac{x}{x+1}$, so $y=x+2$ is the oblique asymptote and $d(x)=\dfrac{x}{x+1}$, which vanishes at $x=0$: the curve \textbf{crosses} its asymptote at the point $(0,2)$, passing from above (for $-1<x<0$, where $d>0$) to below (for $x>0$, where $d<0$).
\end{exoresolu}

\begin{attention}[Sign table must respect the domain]
When making the sign table of $d(x)$, always include the forbidden values of $f$ (double bar $\|$). The position of the curve can change across a vertical asymptote even when $d$ does not vanish there, exactly as in the two examples above.
\end{attention}

\begin{aretenir}[Key points of the section]
\begin{itemize}
\item Intercepts: compute $f(0)$ and solve $f(x)=0$ \emph{within the domain}.
\item $\lim_{x\to a}f(x)=\pm\infty \Rightarrow$ vertical asymptote $x=a$; $\lim_{x\to\pm\infty}f(x)=b \Rightarrow$ horizontal asymptote $y=b$.
\item Oblique asymptote: $m=\lim \dfrac{f(x)}{x}$, then $c=\lim[f(x)-mx]$; \textbf{both} must be finite, with $m\neq 0$.
\item If the oblique search fails ($f(x)/x\to\infty$ or $\to 0$), the branch is \textbf{parabolic} (vertical or horizontal direction).
\item Position relative to a line: sign of $d(x)=f(x)-g(x)$; the curve crosses its horizontal/oblique asymptote where $d(x)=0$.
\item A horizontal or oblique asymptote \emph{may} be crossed; a vertical one never.
\end{itemize}
\end{aretenir}

\begin{exercice}[Asymptotes and position]
Let $f(x)=\dfrac{x^{2}+2x+2}{x+1}$ and $\mathcal{C}_f$ its curve.
\begin{enumerate}
\item Give the domain of $f$ and show that $x=-1$ is a vertical asymptote.
\item Show that the line $y=x+1$ is an oblique asymptote at $+\infty$ and at $-\infty$.
\item Study the position of $\mathcal{C}_f$ relative to the line $y=x+1$.
\item Find the intercepts of $\mathcal{C}_f$ with the axes.
\end{enumerate}
\end{exercice}

\begin{corrige}[Exercise 1]
\textbf{1.} Domain: $\mathbb{R}\setminus\{-1\}$. $\lim_{x\to -1}(x^2+2x+2)=1>0$ and $x+1\to 0$, so $\lim_{x\to -1^{+}}f=+\infty$, $\lim_{x\to -1^{-}}f=-\infty$: $x=-1$ is a vertical asymptote.

\textbf{2.} $\dfrac{f(x)}{x}=\dfrac{x^{2}+2x+2}{x^{2}+x}\to 1=m$. Then $f(x)-x=\dfrac{x^{2}+2x+2-x^{2}-x}{x+1}=\dfrac{x+2}{x+1}\to 1=c$. Hence $y=x+1$ is an oblique asymptote at both ends.

\textbf{3.} $d(x)=f(x)-(x+1)=\dfrac{x^{2}+2x+2-(x+1)^{2}}{x+1}=\dfrac{1}{x+1}$. So $d(x)<0$ on $\left]-\infty,-1\right[$ (curve below the line) and $d(x)>0$ on $\left]-1,+\infty\right[$ (curve above the line); $d$ never vanishes, so no crossing.

\textbf{4.} $f(0)=2$: $y$-intercept $(0,2)$. $f(x)=0 \iff x^{2}+2x+2=0$, discriminant $4-8=-4<0$: no $x$-intercept.
\end{corrige}

\begin{savaistu}[Asymptotes in real life]
Asymptotes are not only exam material. The time needed to transmit data over a network as more users connect, the concentration of a medicine in the blood after an injection, or the cost per unit in a factory in Douala as production grows — all these quantities approach horizontal asymptotes. Recognising an asymptote means understanding the long-term behaviour of a real system.
\end{savaistu}

\coursec{Stationary Points and Variation of a Function}

In the previous section we learned how to describe the \emph{global} behaviour of a curve: where it meets the axes, how it behaves at infinity, and which lines it approaches. We now turn to the \emph{local} behaviour: where does the curve rise, where does it fall, and where does it turn around? The tool that answers all these questions is the \cle{derivative}, and the document that summarises all the answers is the \cle{table of variation}. Mastering this table is essential: at the GCE Advanced Level, almost every curve-sketching question requires it, and it is also the key to counting the solutions of an equation $f(x)=k$.

\courssub{Derivative and Monotonicity}

\textbf{Intuition.}\par
Imagine driving from Bamenda to Bafoussam along a hilly road, and think of the road profile as the graph of a function $f$. Where the car climbs, the slope is positive; where it descends, the slope is negative. The slope of the road at a point is exactly the slope of the tangent line to the graph, that is, the number $f'(x)$. It is therefore natural that the \emph{sign} of $f'$ should tell us whether $f$ is increasing or decreasing.

\begin{propriete}[Derivative and monotonicity]
Let $f$ be differentiable on an interval $I$.
\begin{itemize}
\item If $f'(x)>0$ for all $x\in I$ (except possibly at isolated points where $f'=0$), then $f$ is \cle{strictly increasing} on $I$.
\item If $f'(x)<0$ for all $x\in I$ (except possibly at isolated points where $f'=0$), then $f$ is \cle{strictly decreasing} on $I$.
\item If $f'(x)=0$ for all $x\in I$, then $f$ is \cle{constant} on $I$.
\end{itemize}
\end{propriete}

\textbf{Why this is true.}\par
The rigorous proof uses the Mean Value Theorem (which is part of the Further Mathematics syllabus). The idea is simple: if $f'(x)>0$ on $I$, then for any $a<b$ in $I$ the chord joining $(a,f(a))$ to $(b,f(b))$ has gradient $\dfrac{f(b)-f(a)}{b-a}$. The Mean Value Theorem guarantees a point $c\in(a,b)$ where the tangent is parallel to this chord: $f'(c)=\dfrac{f(b)-f(a)}{b-a}>0$. Since $b-a>0$, we get $f(b)-f(a)>0$, i.e.\ $f(b)>f(a)$: the function increases.

\begin{attention}[Sign of $f'$ versus sign of $f$]
Do not confuse the sign of $f'(x)$ with the sign of $f(x)$. The sign of $f$ tells you whether the curve is \emph{above or below} the $x$-axis; the sign of $f'$ tells you whether the curve is \emph{rising or falling}. A function can be negative and increasing (for example $f(x)=x$ for $x<0$).
\end{attention}

\courssub{Stationary Points and Their Nature}

\textbf{Intuition.}\par
At the very top of a hill and at the very bottom of a valley, the road is momentarily flat: the tangent to the graph is horizontal, so $f'(x)=0$. Such points deserve a name.

\begin{definition}[Stationary point]
A \cle{stationary point} of a differentiable function $f$ is a point $a$ of its domain such that
\[ f'(a)=0. \]
Geometrically, the tangent to the curve at $(a,f(a))$ is \textbf{horizontal}.
\end{definition}

A stationary point is of one of three types:
\begin{itemize}
\item a \cle{local maximum}: $f(a)\geq f(x)$ for all $x$ near $a$ (the curve turns from rising to falling);
\item a \cle{local minimum}: $f(a)\leq f(x)$ for all $x$ near $a$ (the curve turns from falling to rising);
\item a \cle{horizontal point of inflexion}: the tangent is horizontal but the curve does not turn around; it keeps rising (or keeps falling), as for $f(x)=x^3$ at $x=0$.
\end{itemize}

\begin{attention}[$f'(a)=0$ does not guarantee an extremum]
The most frequent mistake at the GCE is to write ``$f'(a)=0$, therefore $f$ has a maximum (or minimum) at $a$''. This is \textbf{false} in general: for $f(x)=x^3$ we have $f'(0)=0$, yet $f$ is strictly increasing on $\mathbb{R}$ and has no extremum at $0$. The condition $f'(a)=0$ is \emph{necessary} but \emph{not sufficient}. You must always check the \textbf{sign change} of $f'$ around $a$ (first derivative test) or use the second derivative test.
\end{attention}

\begin{propriete}[First derivative test]
Let $f$ be differentiable on an open interval containing $a$, with $f'(a)=0$.
\begin{itemize}
\item If $f'$ changes sign from $+$ to $-$ at $a$, then $f$ has a \cle{local maximum} at $a$.
\item If $f'$ changes sign from $-$ to $+$ at $a$, then $f$ has a \cle{local minimum} at $a$.
\item If $f'$ keeps the same sign on both sides of $a$, then $a$ is a \cle{horizontal point of inflexion}.
\end{itemize}
\emph{Justification.} If $f'>0$ just before $a$ and $f'<0$ just after, then $f$ is increasing before $a$ and decreasing after $a$, so $f(x)\leq f(a)$ on each side: $f(a)$ is a local maximum. The other cases are similar. This justification is examinable as a short reasoning question.
\end{propriete}

\begin{propriete}[Second derivative test]
Let $f$ be twice differentiable near $a$, with $f'(a)=0$.
\begin{itemize}
\item If $f''(a)>0$, then $f$ has a \cle{local minimum} at $a$.
\item If $f''(a)<0$, then $f$ has a \cle{local maximum} at $a$.
\item If $f''(a)=0$, the test is \textbf{inconclusive}: use the first derivative test.
\end{itemize}
\emph{Idea of the proof.} If $f''(a)>0$, then $f'$ is increasing through $a$; since $f'(a)=0$, $f'$ must be negative just before $a$ and positive just after, which is exactly the minimum case of the first derivative test.
\end{propriete}

\begin{attention}[The inconclusive case]
$f''(a)=0$ does \emph{not} mean ``no extremum''. Compare $f(x)=x^4$ and $g(x)=x^3$ at $a=0$: both have $f'(0)=f''(0)=0$, but $x^4$ has a minimum at $0$ while $x^3$ has a horizontal inflexion. When $f''(a)=0$, go back to the sign of $f'$.
\end{attention}

\begin{exemplebox}[Classifying stationary points]
Let $f(x)=x^3-3x$. Then $f'(x)=3x^2-3=3(x-1)(x+1)$, so the stationary points are $x=-1$ and $x=1$. Also $f''(x)=6x$.
\begin{itemize}
\item At $x=-1$: $f''(-1)=-6<0$, so $f$ has a \textbf{local maximum} $f(-1)=2$.
\item At $x=1$: $f''(1)=6>0$, so $f$ has a \textbf{local minimum} $f(1)=-2$.
\end{itemize}
\end{exemplebox}

\begin{popfigure}
\begin{center}
\begin{tikzpicture}
\begin{axis}[
  axis lines=middle, width=12cm, height=8cm,
  xlabel={$x$}, ylabel={$y$},
  xmin=-2.8, xmax=2.8, ymin=-3.2, ymax=3.2,
  xtick={-2,-1,1,2}, ytick={-2,2},
  domain=-2.3:2.3, samples=120, clip=false]
\addplot[popblue, very thick] {x^3-3*x};
\addplot[poporange, thick, dashed, domain=-2:0] {2};
\addplot[poporange, thick, dashed, domain=0:2] {-2};
\addplot[only marks, mark=*, mark size=2pt, popdark] coordinates {(-1,2) (1,-2)};
\node[popdark, anchor=south east] at (axis cs:-1,2) {max $(-1,\,2)$};
\node[popdark, anchor=north west] at (axis cs:1,-2) {min $(1,\,-2)$};
\end{axis}
\end{tikzpicture}
\end{center}
\legende{The curve of $f(x)=x^3-3x$ with horizontal tangents (dashed) at its stationary points.}
\end{popfigure}

\courssub{Sign Table of $f'(x)$ and Construction of the Table of Variation}

To exploit the link between the sign of $f'$ and the monotonicity of $f$, we first study the sign of $f'$ on the whole domain. The strategy is always the same: \textbf{factorise} $f'(x)$ as much as possible, find its zeros (the \cle{critical values}), and use the sign of each factor.

\begin{methode}[Building a table of variation in four steps]
\begin{enumerate}
\item \textbf{Domain.} Determine the domain of definition $D_f$ and compute the limits of $f$ at the bounds of $D_f$ (open bounds, vertical asymptotes, $\pm\infty$).
\item \textbf{Derivative.} Compute $f'(x)$ and factorise it completely; solve $f'(x)=0$ to find the critical values, and note any value where $f'$ does not exist.
\item \textbf{Sign.} Build the sign table of $f'(x)$ over $D_f$, marking every critical value and every excluded value with the appropriate symbol ($0$ or double bar).
\item \textbf{Limits and extrema.} Draw the arrows of variation ($\nearrow$ where $f'>0$, $\searrow$ where $f'<0$) and write the value of $f$ at each remarkable point: limits at the bounds and exact values $f(a)$ at the stationary points.
\end{enumerate}
\end{methode}

\begin{exoresolu}[Complete study of $f(x)=x^3-3x$]
Study the variations of $f(x)=x^3-3x$ on $\mathbb{R}$ and draw its table of variation.
\tcblower
\textbf{Step 1 — Domain and limits.} $f$ is a polynomial, so $D_f=\mathbb{R}$. As $x\to+\infty$, $x^3$ dominates, so $\lim_{+\infty}f=+\infty$; as $x\to-\infty$, $\lim_{-\infty}f=-\infty$.

\textbf{Step 2 — Derivative.} $f'(x)=3x^2-3=3(x-1)(x+1)$; critical values $x=-1$ and $x=1$.

\textbf{Step 3 — Sign of $f'$.} The quadratic $3(x-1)(x+1)$ is positive outside its roots and negative between them: $f'>0$ on $(-\infty,-1)\cup(1,+\infty)$ and $f'<0$ on $(-1,1)$.

\textbf{Step 4 — Values.} $f(-1)=-1+3=2$ (local maximum), $f(1)=1-3=-2$ (local minimum). Hence the table:

\begin{center}
\begin{tabular}{|c|ccccccc|}
\hline
$x$ & $-\infty$ & & $-1$ & & $1$ & & $+\infty$ \\
\hline
$f'(x)$ & & $+$ & $0$ & $-$ & $0$ & $+$ & \\
\hline
$f(x)$ & $-\infty$ & $\nearrow$ & $2$ & $\searrow$ & $-2$ & $\nearrow$ & $+\infty$ \\
\hline
\end{tabular}
\end{center}

The curve rises from $-\infty$ to the peak $(-1,2)$, falls to the valley $(1,-2)$, then rises to $+\infty$.
\end{exoresolu}

\begin{popfigure}
\begin{center}
\begin{tikzpicture}[scale=1.0]
\draw[popline] (0,0) rectangle (10.5,3.4);
\draw[popline] (1.7,0) -- (1.7,3.4);
\draw[popline] (0,2.3) -- (10.5,2.3);
\draw[popline] (0,1.2) -- (10.5,1.2);
\node at (0.85,2.85) {$x$};
\node at (0.85,1.75) {$f'(x)$};
\node at (0.85,0.6) {$f(x)$};
\node at (2.6,2.85) {$-\infty$};
\node at (4.7,2.85) {$-1$};
\node at (6.8,2.85) {$1$};
\node at (9.6,2.85) {$+\infty$};
\node at (3.6,1.75) {$+$};
\node at (4.7,1.75) {$0$};
\node at (5.75,1.75) {$-$};
\node at (6.8,1.75) {$0$};
\node at (8.2,1.75) {$+$};
\node at (2.6,0.25) {$-\infty$};
\node at (4.7,1.0) {$2$};
\node at (6.8,0.25) {$-2$};
\node at (9.6,1.0) {$+\infty$};
\draw[->, popblue, thick] (2.9,0.35) -- (4.4,0.9);
\draw[->, popblue, thick] (5.0,0.9) -- (6.5,0.35);
\draw[->, popblue, thick] (7.1,0.35) -- (9.3,0.9);
\end{tikzpicture}
\end{center}
\legende{Template of the table of variation of $f(x)=x^3-3x$.}
\end{popfigure}

\begin{attention}[Values where $f$ is not defined]
Never write a finite value of $f$ in the table at a point where $f$ is \emph{not defined}. For example, for $g(x)=\dfrac{x^2}{x-1}$, the value $x=1$ is excluded from the domain: in the table you must place a \textbf{double vertical bar} at $x=1$ in every row and write the \emph{limits} $\lim_{x\to1^-}g(x)=-\infty$ and $\lim_{x\to1^+}g(x)=+\infty$ on either side — never $g(1)$. A value written under a double bar is a serious error that examiners penalise immediately.
\end{attention}

\courssub{Reading Information from the Table of Variation}

A correct table of variation is a goldmine of information. From it we read directly:
\begin{itemize}
\item the \cle{extrema} of $f$ (local maxima and minima, with their values);
\item the \cle{range} of $f$ on each monotonic branch;
\item the \textbf{number of solutions} of the equation $f(x)=k$ for any constant $k$.
\end{itemize}

\begin{propriete}[Counting the solutions of $f(x)=k$]
On each interval where $f$ is continuous and strictly monotonic, the equation $f(x)=k$ has \textbf{exactly one} solution if $k$ lies strictly between the values (or limits) of $f$ at the ends of that interval, and \textbf{no} solution otherwise. This follows from the Intermediate Value Theorem (existence) together with strict monotonicity (uniqueness). The total number of solutions is the sum over all monotonic branches, and it is read directly from the table of variation.
\end{propriete}

\begin{exemplebox}[Counting solutions from the table]
For $f(x)=x^3-3x$, the table above shows three monotonic branches:
\begin{itemize}
\item on $(-\infty,-1]$, $f$ goes from $-\infty$ to $2$: one solution of $f(x)=k$ iff $k\leq 2$;
\item on $[-1,1]$, $f$ goes from $2$ down to $-2$: one solution iff $-2\leq k\leq 2$;
\item on $[1,+\infty)$, $f$ goes from $-2$ to $+\infty$: one solution iff $k\geq -2$.
\end{itemize}
Counting carefully (and not double-counting the shared endpoints): if $|k|<2$ the equation has \textbf{three} solutions; if $k=\pm2$ it has \textbf{two} solutions (one of them a double root at the extremum); if $|k|>2$ it has \textbf{one} solution. This is exactly the kind of discussion the GCE Board asks for.
\end{exemplebox}

\courssub{Interpretation of the Sign of $f(x)$: Position of the Curve Relative to the $x$-Axis}

The table of variation also tells us the \emph{sign} of $f$ itself, which locates the curve relative to the $x$-axis. For a point $(x,y)$ in the plane, the sign of $y-f(x)$ tells us whether the point lies above ($y-f(x)>0$) or below ($y-f(x)<0$) the curve $y=f(x)$. In particular, on the curve itself $y=f(x)$ so $y-f(x)=0$.
\begin{itemize}
\item if $f(x)>0$ at $x$, the curve lies \textbf{above} the $x$-axis at that abscissa;
\item if $f(x)<0$, the curve lies \textbf{below} the $x$-axis;
\item the zeros of $f$ are the abscissae of the \textbf{intercepts} with the $x$-axis.
\end{itemize}
More generally, comparing the curve with a line $y=mx+c$ amounts to studying the sign of the difference $g(x)=f(x)-(mx+c)$: the curve is above the line where $g(x)>0$ and below it where $g(x)<0$. This idea will be developed fully when we study the position of a curve with respect to a line.

\begin{exoresolu}[Sign of $f$ from the table of variation]
Using the table of variation of $f(x)=x^3-3x$, determine the sign of $f(x)$ on $\mathbb{R}$.
\tcblower
Factorise: $f(x)=x(x^2-3)=x(x-\sqrt{3})(x+\sqrt{3})$, so the zeros are $-\sqrt{3},\,0,\,\sqrt{3}$. From the table: on $(-\infty,-1]$, $f$ increases from $-\infty$ to $2$, so it crosses zero exactly once, at $x=-\sqrt{3}$; hence $f<0$ on $(-\infty,-\sqrt{3})$ and $f>0$ on $(-\sqrt{3},-1]$. On $[-1,1]$, $f$ decreases from $2$ to $-2$, crossing zero once, at $x=0$: $f>0$ on $[-1,0)$ and $f<0$ on $(0,1]$. On $[1,+\infty)$, $f$ increases from $-2$ to $+\infty$, crossing zero once, at $x=\sqrt{3}$: $f<0$ on $[1,\sqrt{3})$ and $f>0$ on $(\sqrt{3},+\infty)$. Conclusion: the curve is above the $x$-axis on $(-\sqrt{3},0)\cup(\sqrt{3},+\infty)$ and below it on $(-\infty,-\sqrt{3})\cup(0,\sqrt{3})$.
\end{exoresolu}

\begin{exercice}[Practice]
Let $f(x)=x^3-3x^2+2$ defined on $\mathbb{R}$.
\begin{enumerate}
\item Compute $f'(x)$, factorise it, and find the stationary points.
\item Use the second derivative test to determine their nature.
\item Draw the complete table of variation of $f$.
\item Deduce the number of solutions of $f(x)=k$ according to the values of $k$.
\end{enumerate}
\end{exercice}

\begin{corrige}[Exercise 1]
1. $f'(x)=3x^2-6x=3x(x-2)$; stationary points $x=0$ and $x=2$.

2. $f''(x)=6x-6$. $f''(0)=-6<0$: local maximum $f(0)=2$. $f''(2)=6>0$: local minimum $f(2)=8-12+2=-2$.

3. Limits: $\lim_{-\infty}f=-\infty$, $\lim_{+\infty}f=+\infty$. Sign of $f'$: positive on $(-\infty,0)\cup(2,+\infty)$, negative on $(0,2)$.
\begin{center}
\begin{tabular}{|c|ccccccc|}
\hline
$x$ & $-\infty$ & & $0$ & & $2$ & & $+\infty$ \\
\hline
$f'(x)$ & & $+$ & $0$ & $-$ & $0$ & $+$ & \\
\hline
$f(x)$ & $-\infty$ & $\nearrow$ & $2$ & $\searrow$ & $-2$ & $\nearrow$ & $+\infty$ \\
\hline
\end{tabular}
\end{center}

4. By the Intermediate Value Theorem on each branch: if $k>2$ or $k<-2$, one solution; if $k=2$ or $k=-2$, two solutions; if $-2<k<2$, three solutions.
\end{corrige}

\begin{aretenir}[Key points of this section]
\begin{itemize}
\item $f$ increases where $f'>0$ and decreases where $f'<0$; the sign of $f'$ governs the \emph{shape}, the sign of $f$ governs the \emph{position relative to the $x$-axis}.
\item A stationary point satisfies $f'(a)=0$ (horizontal tangent); it is a maximum, a minimum, or a horizontal inflexion — decide with the sign change of $f'$ or with $f''(a)$.
\item The table of variation is built in four steps: domain and limits, derivative, sign of $f'$, arrows with exact values.
\item The number of solutions of $f(x)=k$ is read from the table by applying the Intermediate Value Theorem on each monotonic branch.
\item Never place a value of $f$ at a point excluded from the domain: use a double bar and limits.
\end{itemize}
\end{aretenir}

With the variations of $f$ now under control, one refinement remains before we can sketch curves with full confidence: the way the curve \emph{bends}. This is the question of concavity and points of inflexion, which we study in the next section.

\coursec{Concavity, Points of Inflexion and Symmetry}

In the previous section we learnt how the \emph{first} derivative $f'$ reveals where a function is increasing or decreasing and where its stationary points lie. But two curves can both be increasing on an interval and yet look completely different: one may bend upwards like a cup, the other downwards like a cap. To distinguish these shapes we need the \emph{second} derivative $f''$. In this section we study the \cle{concavity} of a curve, the \cle{points of inflexion} where the bending changes direction, and the \cle{symmetries} (centre and axis) that allow us to halve our work when sketching a curve. By the end, you will have the complete study plan of a function.

\courssub{Concavity of a Curve}

\textbf{Intuition.} Imagine driving along the graph of $f$ from left to right. If the road bends so that you keep turning \emph{left} (anti-clockwise), the curve is \cle{concave up}; if you keep turning \emph{right} (clockwise), it is \cle{concave down}. Equivalently: a concave-up curve \emph{holds water} like a bowl; a concave-down curve \emph{sheds water} like a roof.

Geometrically, concavity is detected by the position of the curve relative to its \textbf{tangent lines}:
\begin{itemize}
\item if near a point the curve lies \textbf{above} each of its tangents, it is concave up;
\item if near a point the curve lies \textbf{below} each of its tangents, it is concave down.
\end{itemize}

Why does the second derivative decide this? The slope of the tangent is $f'(x)$. If $f''(x) > 0$ on an interval, then $f'$ is \emph{increasing}: the tangent turns anti-clockwise as $x$ grows, so the curve bends upwards and stays above its tangents. If $f''(x) < 0$, then $f'$ is decreasing and the curve bends downwards, staying below its tangents.

\begin{definition}[Concavity]
Let $f$ be twice differentiable on an interval $I$.
\begin{itemize}
\item $f$ is \cle{concave up} (convex) on $I$ if $f''(x) > 0$ for all $x \in I$; the curve lies above each of its tangents on $I$.
\item $f$ is \cle{concave down} (concave) on $I$ if $f''(x) < 0$ for all $x \in I$; the curve lies below each of its tangents on $I$.
\end{itemize}
\end{definition}

\begin{popfigure}
\begin{center}
\begin{tikzpicture}[scale=1.0, line cap=round]
% Concave up arc: y = 0.5 x^2
\begin{scope}[shift={(-3.2,0)}]
\draw[thick, popblue, domain=-1.4:1.4, samples=40] plot (\x, {0.5*\x*\x});
\draw[poporange, thick, domain=-1.9:1.9] plot (\x, {0});
\draw[popgreen, thick, domain=-2.0:0.6] plot (\x, {-\x - 0.5});
\fill[popdark] (0,0) circle (1.5pt);
\fill[popdark] (-1,0.5) circle (1.5pt);
\node[popblue, align=center] at (0,-1.3) {Concave up: $f''>0$};
\node[popmuted, align=center, text width=3.4cm] at (0,-2.1) {curve above its tangents};
\end{scope}
% Concave down arc: y = 1.5 - 0.5 x^2
\begin{scope}[shift={(3.2,0)}]
\draw[thick, poppurple, domain=-1.4:1.4, samples=40] plot (\x, {1.5 - 0.5*\x*\x});
\draw[poporange, thick, domain=-1.9:1.9] plot (\x, {1.5});
\draw[popgreen, thick, domain=-0.6:2.0] plot (\x, {2 - \x});
\fill[popdark] (0,1.5) circle (1.5pt);
\fill[popdark] (1,1.0) circle (1.5pt);
\node[poppurple, align=center] at (0,-1.3) {Concave down: $f''<0$};
\node[popmuted, align=center, text width=3.4cm] at (0,-2.1) {curve below its tangents};
\end{scope}
\end{tikzpicture}
\end{center}
\legende{Concavity seen through tangent lines: concave up (left, tangents at $x=0$ and $x=-1$), concave down (right, tangents at $x=0$ and $x=1$).}
\end{popfigure}

\begin{exemplebox}[Testing concavity]
Let $f(x) = x^3 - 3x$. Then $f'(x) = 3x^2 - 3$ and $f''(x) = 6x$.
\begin{itemize}
\item For $x < 0$, $f''(x) < 0$: the curve is concave down on $\left]-\infty, 0\right[$.
\item For $x > 0$, $f''(x) > 0$: the curve is concave up on $\left]0, +\infty\right[$.
\end{itemize}
\end{exemplebox}

\courssub{Points of Inflexion}

\textbf{Intuition.} When you drive over the top of a hill and the road changes from bending right to bending left, there is exactly one instant when the steering wheel passes through the neutral position. On a curve, that instant is a \cle{point of inflexion}: the point where the concavity changes.

\begin{definition}[Point of inflexion]
A point $M(a, f(a))$ on the curve of $f$ is a \cle{point of inflexion} if the concavity of $f$ changes at $a$: the curve crosses its tangent at $M$, passing from one side of the tangent to the other.
\end{definition}

\begin{propriete}[Condition for a point of inflexion]
Let $f$ be twice differentiable on an open interval containing $a$.
\begin{itemize}
\item \textbf{Necessary condition:} if $(a, f(a))$ is a point of inflexion, then $f''(a) = 0$.
\item \textbf{Sufficient condition:} if $f''(a) = 0$ \emph{and} $f''$ \textbf{changes sign} at $a$ (i.e.\ the concavity changes), then $(a, f(a))$ is a point of inflexion.
\end{itemize}
(The sign change of $f''$ is the examinable justification; always exhibit it, e.g.\ in a sign table.)
\end{propriete}

\begin{attention}[$f''(a) = 0$ is not enough]
A very common mistake at the GCE is to write ``$f''(a)=0$, hence inflexion point''. This is \textbf{false} without a sign change. Counter-example: $f(x) = x^4$ has $f''(x) = 12x^2$, so $f''(0) = 0$, yet $f''(x) > 0$ on both sides of $0$: the curve stays concave up and $(0,0)$ is \emph{not} an inflexion point (it is a minimum). \textbf{Always verify that $f''$ changes sign.}
\end{attention}

\textbf{Stationary versus ordinary inflexion points.} An inflexion point $(a,f(a))$ is called:
\begin{itemize}
\item a \cle{stationary inflexion point} if in addition $f'(a) = 0$: the tangent there is \emph{horizontal} (example: $f(x)=x^3$ at the origin);
\item an \cle{ordinary inflexion point} if $f'(a) \neq 0$: the tangent there is oblique (example: $f(x) = x^3 - 3x$ at the origin, where $f'(0) = -3$).
\end{itemize}
To classify, compute $f'(a)$ after confirming the sign change of $f''$.

\begin{exoresolu}[Inflexion points of a cubic]
\textbf{Statement.} Let $f(x) = x^3 - 3x$ defined on $\mathbb{R}$. Determine the concavity of $f$ and find its points of inflexion, specifying their nature.
\tcblower
\textbf{Solution.} We compute $f'(x) = 3x^2 - 3$ and $f''(x) = 6x$.
$f''(x) = 0 \iff x = 0$, and $f''(x) = 6x$ changes sign at $0$:
\begin{center}
\begin{tabular}{|c|ccccc|}
\hline
$x$ & $-\infty$ & & $0$ & & $+\infty$ \\ \hline
$f''(x)$ & & $-$ & $0$ & $+$ & \\ \hline
Concavity & & down & infl. & up & \\ \hline
\end{tabular}
\end{center}
Since $f''$ changes sign at $0$, the point $(0, f(0)) = (0,0)$ is a point of inflexion. Moreover $f'(0) = -3 \neq 0$, so it is an \textbf{ordinary inflexion point}; the tangent there has equation $y = -3x$, and indeed the curve crosses this line at the origin since $x^3 - 3x - (-3x) = x^3$ changes sign at $0$.
\end{exoresolu}

\courssub{Centre of Symmetry of a Curve}

\textbf{Intuition.} Some curves look exactly the same after a half-turn about a special point $I$. The graph of $y = x^3$ rotated through $180^{\circ}$ about the origin falls back on itself. Such a point $I$ is called a \cle{centre of symmetry}. Detecting it is precious: half of the curve determines the other half.

\begin{definition}[Centre of symmetry]
The point $I(a, b)$ is a \cle{centre of symmetry} of the curve $\mathcal{C}_f$ if for every $x$ such that $a+x$ and $a-x$ are in the domain,
\[ f(a+x) + f(a-x) = 2b. \]
Equivalently, with $M(a+x, f(a+x))$ and $M'(a-x, f(a-x))$, the point $I$ is the midpoint of $[MM']$.
\end{definition}

\begin{methode}[Proving that $I(a,b)$ is a centre of symmetry]
\begin{enumerate}
\item Check the domain is symmetric about $a$: $a+x \in D_f \iff a-x \in D_f$.
\item Compute $f(a+x)$ and $f(a-x)$ separately.
\item Add them and simplify.
\item If $f(a+x) + f(a-x) = 2b$ for all admissible $x$, conclude that $I(a,b)$ is a centre of symmetry.
\end{enumerate}
To \emph{find} a candidate $(a,b)$, look at the domain (e.g.\ an excluded value $a$ for a rational function) and compute $b = f(a)$ or use limits and asymptotes.
\end{methode}

\begin{exoresolu}[Centre of symmetry of a hyperbola]
\textbf{Statement.} Let $f(x) = \dfrac{2x+1}{x-1}$ on $D_f = \mathbb{R} \setminus \{1\}$. Show that $I(1, 2)$ is a centre of symmetry of $\mathcal{C}_f$.
\tcblower
\textbf{Solution.} The domain is symmetric about $1$ since $1+x \in D_f \iff x \neq 0 \iff 1-x \in D_f$. For $x \neq 0$:
\begin{align*}
f(1+x) &= \frac{2(1+x)+1}{(1+x)-1} = \frac{2x+3}{x} = 2 + \frac{3}{x}, \\
f(1-x) &= \frac{2(1-x)+1}{(1-x)-1} = \frac{3-2x}{-x} = 2 - \frac{3}{x}.
\end{align*}
Hence
\[ f(1+x) + f(1-x) = \left(2 + \tfrac{3}{x}\right) + \left(2 - \tfrac{3}{x}\right) = 4 = 2 \times 2. \]
Therefore $I(1,2)$ is a centre of symmetry of $\mathcal{C}_f$. (Notice $a=1$ is the excluded value and $b=2$ is the horizontal asymptote: for a homographic function the centre of symmetry is the intersection of the two asymptotes.)
\end{exoresolu}

\courssub{Axis of Symmetry of a Curve}

\textbf{Intuition.} Other curves are their own mirror image in a vertical line, like the parabola $y = x^2$ in the $y$-axis. Such a line is an \cle{axis of symmetry}.

\begin{definition}[Axis of symmetry]
The vertical line $x = a$ is an \cle{axis of symmetry} of $\mathcal{C}_f$ if for every $x$ such that $a+x$ and $a-x$ are in the domain,
\[ f(a+x) = f(a-x). \]
The points $M(a+x, f(a+x))$ and $M'(a-x, f(a-x))$ are then mirror images in the line $x = a$.
\end{definition}

\begin{methode}[Proving that the line $x = a$ is an axis of symmetry]
\begin{enumerate}
\item Check the domain is symmetric about $a$.
\item Compute $f(a+x)$ and $f(a-x)$.
\item Show the two expressions are \emph{equal} for all admissible $x$.
\item Conclude that the line $x = a$ is an axis of symmetry.
\end{enumerate}
\end{methode}

\begin{exemplebox}[Axis of symmetry of a parabola]
Let $f(x) = x^2 - 4x + 5$ on $\mathbb{R}$. We test the line $x = 2$:
\[ f(2+x) = (2+x)^2 - 4(2+x) + 5 = x^2 + 1, \qquad f(2-x) = (2-x)^2 - 4(2-x) + 5 = x^2 + 1. \]
Since $f(2+x) = f(2-x)$ for all $x$, the line $x = 2$ is an axis of symmetry. (This is the vertical line through the vertex $(2,1)$.)
\end{exemplebox}

\begin{attention}[Centre versus axis: do not confuse them]
A \textbf{centre of symmetry} is a \emph{point}: the test is $f(a+x) + f(a-x) = 2b$ (a \emph{sum}). An \textbf{axis of symmetry} is a \emph{line} $x = a$: the test is $f(a+x) = f(a-x)$ (an \emph{equality} of values). Mixing the two tests is a classic source of lost marks. Also note: a curve can have both, one, or neither.
\end{attention}

\courssub{Even and Odd Functions: the Special Cases}

Two symmetries occur so often that they have their own names.

\begin{definition}[Even and odd functions]
Let $f$ have a domain symmetric about $0$.
\begin{itemize}
\item $f$ is \cle{even} if $f(-x) = f(x)$ for all $x \in D_f$: the \textbf{$y$-axis} (line $x=0$) is an axis of symmetry.
\item $f$ is \cle{odd} if $f(-x) = -f(x)$ for all $x \in D_f$: the \textbf{origin} $O(0,0)$ is a centre of symmetry.
\end{itemize}
\end{definition}

Examples: $x^2$, $\cos x$, $|x|$ are even; $x^3$, $\sin x$, $\dfrac{1}{x}$ are odd. The function $f(x) = x^3 - 3x$ is odd since $f(-x) = -x^3 + 3x = -f(x)$: its curve has the origin as centre of symmetry.

\begin{popfigure}
\begin{center}
\begin{tikzpicture}
\begin{axis}[
  width=11cm, height=8cm,
  axis lines=middle,
  xlabel={$x$}, ylabel={$y$},
  xmin=-3, xmax=3, ymin=-4, ymax=4,
  xtick={-2,-1,1,2}, ytick={-3,-2,-1,1,2,3},
  domain=-2.15:2.15, samples=120,
  legend style={at={(0.02,0.98)}, anchor=north west, font=\small}
]
\addplot[thick, popblue] {x^3 - 3*x};
\addlegendentry{$f(x)=x^3-3x$}
\addplot[dashed, poporange] {-3*x};
\addlegendentry{tangent $y=-3x$ at $O$}
\addplot[only marks, mark=*, popdark, mark size=1.6pt] coordinates {(0,0) (1,-2) (-1,2)};
\node[popdark, anchor=west, font=\small] at (axis cs:0.1,0.7) {$O$: inflexion, centre of symmetry};
\node[popmuted, anchor=south, font=\small] at (axis cs:1,-2) {min $(1,-2)$};
\node[popmuted, anchor=north, font=\small] at (axis cs:-1,2) {max $(-1,2)$};
\end{axis}
\end{tikzpicture}
\end{center}
\legende{The curve of $f(x)=x^3-3x$: the origin is both a point of inflexion (the curve crosses its tangent $y=-3x$) and a centre of symmetry ($f$ is odd).}
\end{popfigure}

\begin{savaistu}[Why symmetry saves you half the work]
If $I(a,b)$ is a centre of symmetry or $x=a$ is an axis of symmetry, it is enough to study $f$ on \emph{half} the domain (say $x \geq a$): variations, limits and asymptotes on the other half follow by symmetry. For even or odd functions, study on $[0, +\infty[$ only. GCE examiners reward candidates who state explicitly: ``we restrict the study to $\ldots$ and complete by symmetry''.
\end{savaistu}

\courssub{Assembling Everything: the Complete Study Plan}

Concavity, inflexion and symmetry are the final bricks of the full study of a function. Here is the master plan you should follow in every curve-sketching problem.

\begin{methode}[Complete study plan of a function]
\begin{enumerate}
\item \textbf{Domain} $D_f$; reduce it using \textbf{parity / symmetry} (test $f(-x)$, or $f(a \pm x)$).
\item \textbf{Limits} at the bounds of $D_f$; deduce \textbf{asymptotes} (vertical, horizontal, oblique) and study infinite branches.
\item \textbf{Intercepts} with the axes.
\item \textbf{First derivative} $f'$: sign, stationary points, \textbf{table of variation}.
\item \textbf{Second derivative} $f''$: sign, \textbf{concavity}, \textbf{points of inflexion} (justify by the sign change of $f''$).
\item \textbf{Position} of the curve relative to remarkable lines (asymptotes, tangents at inflexion points).
\item \textbf{Sketch} the curve, using symmetry to complete it.
\end{enumerate}
\end{methode}

\begin{exoresolu}[Concavity and symmetry of a rational function]
\textbf{Statement.} Let $g(x) = \dfrac{x^2}{x-1}$ on $\mathbb{R}\setminus\{1\}$. Study the concavity of $g$ and show that $I(1,2)$ is a centre of symmetry.
\tcblower
\textbf{Solution.} \textbf{Concavity.} Write $g(x) = x + 1 + \dfrac{1}{x-1}$ (polynomial division). Then
\[ g'(x) = 1 - \frac{1}{(x-1)^2}, \qquad g''(x) = \frac{2}{(x-1)^3}. \]
$g''(x)$ has the sign of $x-1$: negative on $\left]-\infty,1\right[$ (concave down), positive on $\left]1,+\infty\right[$ (concave up). Since $g$ is not defined at $1$, there is \textbf{no point of inflexion} — the concavity changes across the vertical asymptote $x=1$, not across a point of the curve.

\textbf{Centre of symmetry.} For $x \neq 0$:
\begin{align*}
g(1+x) &= (1+x) + 1 + \frac{1}{x} = 2 + x + \frac{1}{x}, \\
g(1-x) &= (1-x) + 1 - \frac{1}{x} = 2 - x - \frac{1}{x}.
\end{align*}
Thus $g(1+x) + g(1-x) = 4 = 2 \times 2$, so $I(1,2)$ is a centre of symmetry. Note that $I$ is the intersection of the vertical asymptote $x=1$ and the oblique asymptote $y = x+1$.
\end{exoresolu}

\begin{exercice}
Let $f(x) = \dfrac{x^3}{3} - x + 1$ on $\mathbb{R}$.
\begin{enumerate}
\item Compute $f''$ and study the concavity of $f$.
\item Show that $(0,1)$ is a point of inflexion and state whether it is stationary or ordinary.
\item Prove that $I(0,1)$ is a centre of symmetry of $\mathcal{C}_f$.
\end{enumerate}
\end{exercice}

\begin{corrige}[Exercise 1]
\textbf{1.} $f'(x) = x^2 - 1$ and $f''(x) = 2x$. Hence $f''(x) < 0$ on $\left]-\infty,0\right[$ (concave down) and $f''(x) > 0$ on $\left]0,+\infty\right[$ (concave up).

\textbf{2.} $f''(0) = 0$ and $f''$ changes sign at $0$, so $(0, f(0)) = (0,1)$ is a point of inflexion. Since $f'(0) = -1 \neq 0$, it is an \emph{ordinary} inflexion point (tangent of slope $-1$).

\textbf{3.} Compute:
\[ f(x) + f(-x) = \left(\tfrac{x^3}{3} - x + 1\right) + \left(-\tfrac{x^3}{3} + x + 1\right) = 2 = 2 \times 1, \]
so $f(0+x) + f(0-x) = 2 \times 1$ for all $x$: $I(0,1)$ is a centre of symmetry.
\end{corrige}

\begin{exercice}
Let $h(x) = \dfrac{x^2 - 2x + 2}{x - 1}$ on $\mathbb{R}\setminus\{1\}$.
\begin{enumerate}
\item Show that $h(x) = x - 1 + \dfrac{1}{x-1}$.
\item Deduce that $I(1,0)$ is a centre of symmetry of $\mathcal{C}_h$.
\item Compute $h''$ and study the concavity of $h$. Does $h$ have a point of inflexion?
\end{enumerate}
\end{exercice}

\begin{corrige}[Exercise 2]
\textbf{1.} $x - 1 + \dfrac{1}{x-1} = \dfrac{(x-1)^2 + 1}{x-1} = \dfrac{x^2 - 2x + 2}{x-1} = h(x)$.

\textbf{2.} For $x \neq 0$: $h(1+x) = x + \dfrac{1}{x}$ and $h(1-x) = -x - \dfrac{1}{x}$, so $h(1+x) + h(1-x) = 0 = 2 \times 0$: $I(1,0)$ is a centre of symmetry.

\textbf{3.} $h'(x) = 1 - \dfrac{1}{(x-1)^2}$, $h''(x) = \dfrac{2}{(x-1)^3}$. Thus $h$ is concave down on $\left]-\infty,1\right[$ and concave up on $\left]1,+\infty\right[$. There is \textbf{no} point of inflexion because $h$ is not defined at $1$ (the concavity changes across the vertical asymptote).
\end{corrige}

\begin{aretenir}[Key points of this section]
\begin{itemize}
\item $f'' > 0$: concave up, curve above its tangents; $f'' < 0$: concave down, curve below its tangents.
\item Inflexion point at $a$: requires $f''(a) = 0$ \textbf{and} a \textbf{sign change} of $f''$ at $a$. It is \emph{stationary} if also $f'(a)=0$, \emph{ordinary} otherwise.
\item Centre of symmetry $I(a,b)$: prove $f(a+x) + f(a-x) = 2b$.
\item Axis of symmetry $x = a$: prove $f(a+x) = f(a-x)$.
\item Even function $\Rightarrow$ $y$-axis symmetry; odd function $\Rightarrow$ origin symmetry.
\item Symmetry halves the study domain; concavity completes the table of variation and refines the sketch.
\end{itemize}
\end{aretenir}

\coursec{Algebraic Solutions of Absolute-Value Inequalities}

In the previous sections we learned to read the shape of a curve: its concavity, its points of inflexion and its possible symmetries. We now change our point of view. Instead of asking \emph{``what does the graph look like?''}, we ask \emph{``for which values of $x$ does a given condition hold?''} When the condition involves an absolute value, the question becomes an \cle{absolute-value inequality}, and the GCE examiner expects you to solve it \emph{algebraically}, with complete rigour, and to present the answer as a union of intervals. This section gives you every tool you need.

\courssub{Reminder: the Absolute Value and its Properties}

Imagine a trader walking along the Douala--Yaound\'e highway represented as a number line. Whether she is $5$~km to the east or $5$~km to the west of a toll gate, her \emph{distance} from the gate is the same: $5$~km. The absolute value is precisely this idea of distance, stripped of direction.

\begin{definition}[Absolute value]
For any real number $x$, the \cle{absolute value} (or \cle{modulus}) of $x$ is
\[
|x| = \begin{cases} x & \text{if } x \geq 0, \\ -x & \text{if } x < 0. \end{cases}
\]
Geometrically, $|x|$ is the distance from $x$ to $0$ on the number line, and $|a - b|$ is the distance between the numbers $a$ and $b$.
\end{definition}

\begin{propriete}[Main properties of the absolute value]
For all real numbers $x$ and $y$:
\begin{itemize}
\item $|x| \geq 0$, and $|x| = 0 \iff x = 0$;
\item $|-x| = |x|$ and $|x|^2 = x^2$;
\item $|xy| = |x|\,|y|$ and, for $y \neq 0$, $\left|\dfrac{x}{y}\right| = \dfrac{|x|}{|y|}$;
\item $|x + y| \leq |x| + |y|$ (the \cle{triangle inequality});
\item $\sqrt{x^2} = |x|$ (not $x$!).
\end{itemize}
These properties are not usually examined as proofs, but they are used constantly; the last one is a favourite GCE trap.
\end{propriete}

From the ``distance'' picture, three basic facts follow immediately for $a > 0$:
\[
|x| = a \iff x = a \ \text{or}\ x = -a; \qquad
|x| < a \iff -a < x < a; \qquad
|x| > a \iff x < -a \ \text{or}\ x > a.
\]
In interval notation: $|x| < a$ gives $x \in \left]-a\,;\,a\right[$, while $|x| > a$ gives $x \in \left]-\infty\,;\,-a\right[ \cup \left]a\,;\,+\infty\right[$. Notice the \emph{union}: the solution of a ``greater than'' modulus inequality is \textbf{two separate pieces}, never one interval.

\courssub{The Key Equivalences}

Everything in this section rests on one idea: replace the modulus by an equivalent condition \emph{without} modulus.

\begin{propriete}[Fundamental equivalences]
Let $A$ be any real expression and $b > 0$. Then
\[
|A| < b \iff -b < A < b, \qquad\qquad
|A| > b \iff A < -b \ \ \text{or}\ \ A > b.
\]
The same holds with $\leq$ and $\geq$ in place of $<$ and $>$.
\end{propriete}

\textbf{Justification.} $|A|$ is the distance from $A$ to $0$. Saying this distance is less than $b$ means $A$ lies strictly between $-b$ and $b$, i.e.\ $-b < A < b$. Saying the distance is greater than $b$ means $A$ lies outside that segment, on either side: $A < -b$ or $A > b$. $\blacksquare$

\begin{attention}[A union, not a single interval]
A very common mistake is to write the solution of $|A| > b$ as ``$-b > A > b$'' or as one interval. This is meaningless: no number is simultaneously less than $-b$ and greater than $b$. The word is \textbf{or}, and the answer is a \textbf{union} of two intervals.
\end{attention}

\courssub{Solving $|ax + b| < c$ and $|ax + b| > c$}

\begin{methode}[One-modulus inequalities]
\begin{enumerate}
\item Check the right-hand side $c$. If $c \leq 0$, then $|A| < c$ has \textbf{no solution} and $|A| > c$ (with $c<0$) holds for \textbf{all} $x$ in the domain.
\item If $c > 0$, apply the fundamental equivalence to remove the modulus.
\item Solve the resulting double inequality (or the two separate inequalities) for $x$, isolating $x$ in the middle.
\item Write the solution set in interval notation and, if useful, draw it on a number line.
\end{enumerate}
\end{methode}

\begin{exoresolu}[Solving $|2x - 3| < 5$]
Solve in $\mathbb{R}$ the inequality $|2x - 3| < 5$ and illustrate the solution on a number line.
\tcblower
\textbf{Solution.} Since $5 > 0$, the fundamental equivalence gives
\[
|2x - 3| < 5 \iff -5 < 2x - 3 < 5.
\]
Add $3$ to all three parts: $-2 < 2x < 8$. Divide by $2$ (positive, so the inequalities keep their direction):
\[
-1 < x < 4.
\]
Hence $S = \left]-1\,;\,4\right[$. \textbf{Check:} $x = 0$ gives $|{-3}| = 3 < 5$ $\checkmark$; $x = 5$ gives $|7| = 7 \not< 5$ $\checkmark$ (correctly excluded).
\end{exoresolu}

\begin{popfigure}
\begin{center}
\begin{tikzpicture}[scale=1.1]
\draw[->, thick, popink] (-3.5,0) -- (5.5,0);
\foreach \x in {-3,-2,-1,0,1,2,3,4,5}
  \draw[popink] (\x,0.08) -- (\x,-0.08) node[below, font=\small] {$\x$};
\draw[line width=2.5pt, poporange] (-1,0) -- (4,0);
\draw[fill=white, draw=poporange, line width=1.2pt] (-1,0) circle (0.09);
\draw[fill=white, draw=poporange, line width=1.2pt] (4,0) circle (0.09);
\node[poporange, font=\small, align=center] at (1.5,0.45) {$|2x-3|<5$};
\end{tikzpicture}
\end{center}
\legende{Solution set of $|2x-3|<5$: the open interval $\left]-1\,;\,4\right[$ (open dots: endpoints excluded).}
\end{popfigure}

\begin{exemplebox}[A ``greater than'' case]
Solve $|3x + 1| \geq 4$. Since $4 > 0$:
\[
3x + 1 \leq -4 \quad\text{or}\quad 3x + 1 \geq 4
\iff x \leq -\tfrac{5}{3} \quad\text{or}\quad x \geq 1.
\]
So $S = \left]-\infty\,;\,-\tfrac{5}{3}\right] \cup \left[1\,;\,+\infty\right[$. Note the closed brackets because the inequality is non-strict.
\end{exemplebox}

\courssub{Inequalities with Two Moduli: Squaring Both Sides}

How do we solve $|x - 3| < |2x + 1|$? Neither fundamental equivalence applies directly. The rescue comes from a simple observation: \textbf{both sides of $|A| < |B|$ are non-negative}, and the squaring function $t \mapsto t^2$ is strictly increasing on $\left[0\,;\,+\infty\right[$. Therefore squaring preserves the inequality.

\begin{propriete}[Squaring a two-modulus inequality]
For any real expressions $A$ and $B$,
\[
|A| < |B| \iff A^2 < B^2 \iff (A - B)(A + B) < 0,
\]
because $|A|^2 = A^2$ and both sides are non-negative. The same holds with $\leq$, $>$, $\geq$. This derivation \emph{is} examinable: quote it before using it.
\end{propriete}

\begin{methode}[Solving $|A| < |B|$]
\begin{enumerate}
\item Square both sides (legal, since both are $\geq 0$): $A^2 < B^2$.
\item Move everything to one side and factorise as a difference of squares: $(A-B)(A+B) < 0$.
\item Solve the resulting polynomial inequality using a sign table.
\item Write the solution in interval notation and verify with a test value.
\end{enumerate}
\end{methode}

\begin{exoresolu}[Solving $|x - 3| < |2x + 1|$]
Solve in $\mathbb{R}$: $|x - 3| < |2x + 1|$.
\tcblower
\textbf{Solution.} Both sides are non-negative, so we may square:
\[
(x - 3)^2 < (2x + 1)^2
\iff (x-3)^2 - (2x+1)^2 < 0.
\]
Factorise as a difference of two squares with $a = x-3$, $b = 2x+1$:
\[
\big[(x-3) - (2x+1)\big]\big[(x-3) + (2x+1)\big] < 0
\iff (-x - 4)(3x - 2) < 0.
\]
Multiplying by $-1$ (reversing the inequality): $(x + 4)(3x - 2) > 0$. The roots are $x = -4$ and $x = \tfrac{2}{3}$. The product is positive outside the roots:
\[
S = \left]-\infty\,;\,-4\right[ \cup \left]\tfrac{2}{3}\,;\,+\infty\right[.
\]
\textbf{Check:} $x = 0$: $|{-3}| = 3 > |1| = 1$, so $0 \notin S$ $\checkmark$; $x = 1$: $|{-2}| = 2 < |3| = 3$, so $1 \in S$ $\checkmark$.
\end{exoresolu}

\begin{attention}[Never square when a side may be negative]
Squaring is valid \textbf{only} when both sides are known to be non-negative. From $-3 < 2$ one cannot deduce $9 < 4$! So for an inequality like $|x - 1| < 2x + 3$, you must \emph{first} impose the condition $2x + 3 \geq 0$ (a modulus is never less than a negative number), and only then square. Forgetting this condition introduces false solutions — a classic GCE trap.
\end{attention}

\courssub{The Critical Value (Case) Method}

When moduli are \emph{added} rather than compared, squaring does not help. We then use the definition of $|x|$ itself: each modulus changes its formula at the point where its contents vanish. These points are called \cle{critical values}; they split the number line into cases.

\begin{methode}[Critical value method]
\begin{enumerate}
\item Find the zero of \textbf{each} modulus: these are the critical values.
\item Arrange them in increasing order; they divide $\mathbb{R}$ into intervals (cases).
\item On each interval, replace every modulus by its correct expression ($+(\cdots)$ or $-(\cdots)$) using a sign table.
\item Solve the resulting ordinary inequality \emph{on that interval}, and intersect the answer with the interval.
\item Take the \textbf{union} of the partial solutions over all cases.
\end{enumerate}
\end{methode}

\begin{exoresolu}[Solving $|x - 1| + |x + 2| \leq 4$]
Solve in $\mathbb{R}$: $|x - 1| + |x + 2| \leq 4$.
\tcblower
\textbf{Solution.} The critical values are $x = 1$ (from $|x-1|$) and $x = -2$ (from $|x+2|$). Sign table:
\begin{center}
\begin{tabular}{|c|ccccccc|}
\hline
$x$ & $-\infty$ & & $-2$ & & $1$ & & $+\infty$ \\
\hline
$x - 1$ & & $-$ & & $-$ & $0$ & & $+$ \\
\hline
$x + 2$ & & $-$ & $0$ & $+$ & & & $+$ \\
\hline
\end{tabular}
\end{center}
\textbf{Case 1:} $x < -2$. Then $|x-1| = -(x-1)$ and $|x+2| = -(x+2)$:
\[
-(x-1) - (x+2) \leq 4 \iff -2x - 1 \leq 4 \iff x \geq -\tfrac{5}{2}.
\]
Intersecting with $x < -2$: $\left[-\tfrac{5}{2}\,;\,-2\right[$.

\textbf{Case 2:} $-2 \leq x < 1$. Then $|x-1| = -(x-1)$, $|x+2| = x+2$:
\[
-(x-1) + (x+2) \leq 4 \iff 3 \leq 4,
\]
which is always true. The whole interval $\left[-2\,;\,1\right[$ is a solution.

\textbf{Case 3:} $x \geq 1$. Then $|x-1| = x-1$, $|x+2| = x+2$:
\[
(x-1) + (x+2) \leq 4 \iff 2x + 1 \leq 4 \iff x \leq \tfrac{3}{2}.
\]
Intersecting with $x \geq 1$: $\left[1\,;\,\tfrac{3}{2}\right]$.

\textbf{Union:} $S = \left[-\tfrac{5}{2}\,;\,-2\right[ \cup \left[-2\,;\,1\right[ \cup \left[1\,;\,\tfrac{3}{2}\right] = \left[-\tfrac{5}{2}\,;\,\tfrac{3}{2}\right]$.

\textbf{Check:} $x = 0$: $|{-1}| + |2| = 3 \leq 4$ $\checkmark$; $x = 2$: $1 + 4 = 5 \not\leq 4$ $\checkmark$; endpoint $x = -\tfrac{5}{2}$: $\tfrac{7}{2} + \tfrac{1}{2} = 4 \leq 4$ $\checkmark$ (included, as the inequality is non-strict).
\end{exoresolu}

\begin{popfigure}
\begin{center}
\begin{tikzpicture}[scale=1.05]
\draw[->, thick, popink] (-4.5,0) -- (3.5,0);
\foreach \x in {-4,-3,-2,-1,0,1,2,3}
  \draw[popink] (\x,0.08) -- (\x,-0.08) node[below, font=\small] {$\x$};
\draw[fill=white, draw=popblue, line width=1.2pt] (-2,0) circle (0.09);
\draw[fill=white, draw=popblue, line width=1.2pt] (1,0) circle (0.09);
\node[popblue, font=\small, above] at (-2,0.15) {$-2$};
\node[popblue, font=\small, above] at (1,0.15) {$1$};
\node[font=\small, poppurple, align=center] at (-3.2,0.55) {Case 1: $x<-2$};
\node[font=\small, poppurple, align=center] at (-0.5,0.55) {Case 2: $-2\leq x<1$};
\node[font=\small, poppurple, align=center] at (2.2,0.55) {Case 3: $x\geq 1$};
\draw[<->, poppurple, dashed] (-4.3,-0.7) -- (-2.1,-0.7);
\draw[<->, poppurple, dashed] (-1.9,-0.7) -- (0.9,-0.7);
\draw[<->, poppurple, dashed] (1.1,-0.7) -- (3.3,-0.7);
\draw[line width=2.5pt, poporange] (-2.5,0) -- (1.5,0);
\draw[fill=poporange, draw=poporange] (-2.5,0) circle (0.09);
\draw[fill=poporange, draw=poporange] (1.5,0) circle (0.09);
\node[poporange, font=\small, below] at (-2.5,-0.25) {$-\tfrac{5}{2}$};
\node[poporange, font=\small, below] at (1.5,-0.25) {$\tfrac{3}{2}$};
\end{tikzpicture}
\end{center}
\legende{The critical values $-2$ and $1$ split the line into three cases; the union of the partial solutions is the shaded interval $\left[-\tfrac{5}{2}\,;\,\tfrac{3}{2}\right]$.}
\end{popfigure}

\begin{attention}[Intersect before you unite]
In each case, the solution of the simplified inequality is only valid \emph{on that case's interval}. A frequent error is to keep $x \geq -\tfrac{5}{2}$ from Case 1 without intersecting with $x < -2$. Always intersect within the case, then take the union across cases.
\end{attention}

\courssub{Modulus Combined with Rational Expressions}

GCE questions often mix a modulus with a fraction, e.g.\ $\left|\dfrac{x-2}{x+1}\right| < 2$. Two safe strategies exist.

\textbf{Strategy A (quotient of moduli).} Use $\left|\dfrac{A}{B}\right| = \dfrac{|A|}{|B|}$ with $B \neq 0$:
\[
\left|\frac{x-2}{x+1}\right| < 2 \iff \frac{|x-2|}{|x+1|} < 2 \iff |x-2| < 2|x+1| \quad (x \neq -1).
\]
Both sides are non-negative, so we square: $(x-2)^2 < 4(x+1)^2$, i.e.\ $(x-2)^2 - 4(x+1)^2 < 0$, which factorises as $\big[(x-2) - 2(x+1)\big]\big[(x-2) + 2(x+1)\big] = (-x-4)(3x) < 0$, or $x(x+4) > 0$. Hence $x < -4$ or $x > 0$, and $x = -1$ is already excluded. So $S = \left]-\infty\,;\,-4\right[ \cup \left]0\,;\,+\infty\right[$.

\textbf{Strategy B (sign tables per case).} If the inequality has the form $\dfrac{|A|}{B} < c$ with $B$ not in a modulus, split at the zero of $A$ \emph{and} at the zero of $B$, remove the modulus per case, and solve each rational inequality with a sign table, remembering the excluded value.

\begin{exemplebox}[Verification as an exam technique]
Whatever method you use, finish by \textbf{testing sample values}: one inside your proposed solution set, one outside, and each endpoint. This takes less than a minute and catches sign slips, forgotten intersections and wrongly included endpoints. In Strategy A above: $x = 1$ gives $\frac{1}{2} < 2$ $\checkmark$ (in $S$); $x = -2$ gives $4 \not< 2$ $\checkmark$ (not in $S$); $x = -4$ gives $2 \not< 2$ $\checkmark$ (endpoint correctly excluded).
\end{exemplebox}

\begin{aretenir}[Key points of the section]
\begin{itemize}
\item $|A| < b \iff -b < A < b$; \quad $|A| > b \iff A < -b$ \textbf{or} $A > b$ (a \emph{union}).
\item $|A| < |B| \iff A^2 < B^2$: squaring is legal because both sides are $\geq 0$.
\item Never square an inequality when one side may be negative; impose the sign condition first.
\item For sums of moduli, use the critical value method: zeros of each modulus $\to$ cases $\to$ solve per case $\to$ intersect, then unite.
\item Always present the answer in interval notation and verify with test values.
\end{itemize}
\end{aretenir}

\begin{exercice}[Practice]
Solve in $\mathbb{R}$ and give each answer in interval notation:
\begin{enumerate}
\item $|4x - 5| \leq 7$;
\item $|2x + 3| > 1$;
\item $|x + 1| \geq |3x - 5|$;
\item $|x| + |x - 2| < 6$;
\item $\left|\dfrac{2x - 1}{x + 2}\right| \leq 3$.
\end{enumerate}
\end{exercice}

\begin{corrige}[Exercises 1--5]
\textbf{1.} $-7 \leq 4x - 5 \leq 7 \iff -2 \leq 4x \leq 12 \iff -\tfrac12 \leq x \leq 3$; $S = \left[-\tfrac12\,;\,3\right]$.

\textbf{2.} $2x + 3 < -1$ or $2x + 3 > 1 \iff x < -2$ or $x > -1$; $S = \left]-\infty\,;\,-2\right[ \cup \left]-1\,;\,+\infty\right[$.

\textbf{3.} Square: $(x+1)^2 \geq (3x-5)^2 \iff (x+1)^2 - (3x-5)^2 \geq 0 \iff (-2x+6)(4x-4) \geq 0 \iff (x-3)(x-1) \leq 0$; $S = \left[1\,;\,3\right]$.

\textbf{4.} Critical values $0$ and $2$. Case $x<0$: $-x - (x-2) < 6 \iff x > -2$, giving $\left]-2\,;\,0\right[$. Case $0 \leq x < 2$: $x - (x-2) = 2 < 6$ always true, giving $\left[0\,;\,2\right[$. Case $x \geq 2$: $2x - 2 < 6 \iff x < 4$, giving $\left[2\,;\,4\right[$. Union: $S = \left]-2\,;\,4\right[$.

\textbf{5.} $|2x-1| \leq 3|x+2|$, $x \neq -2$. Square: $(2x-1)^2 \leq 9(x+2)^2 \iff (2x-1)^2 - 9(x+2)^2 \leq 0 \iff (-x-7)(5x+5) \leq 0 \iff (x+7)(x+1) \geq 0$; hence $x \leq -7$ or $x \geq -1$ ($-2$ not concerned). $S = \left]-\infty\,;\,-7\right] \cup \left[-1\,;\,+\infty\right[$.
\end{corrige}

With these algebraic techniques mastered, we are ready for the final section, where the same inequalities are solved \emph{graphically} and where all the skills of the chapter come together in complete curve sketching.

\coursec{Graphical Solutions of Inequalities and Complete Curve Sketching}

In the previous section we solved absolute-value inequalities \emph{algebraically}, by squaring, splitting into cases and using the double-inequality trick. Algebra is powerful, but it has a price: heavy computation, and no picture of what is really happening. In this final section we take the opposite route. We learn to \emph{read} solutions directly from a graph, and then we assemble everything studied in this chapter --- domain, limits, asymptotes, variation, concavity, symmetry --- into one complete, examinable skill: the \cle{full study and sketching of a function}, the flagship question of GCE Further Mathematics Paper 2.

\courssub{The Graphical Meaning of an Inequality}

\begin{experience}[One picture is worth a page of algebra]
Draw on the same axes the parabola $y = x^{2}$ and the line $y = x + 2$. They meet where $x^{2} = x + 2$, i.e. at $x = -1$ and $x = 2$. Now ask: \emph{where is the parabola above the line?} Looking at the picture, the parabola dominates outside the interval $[-1, 2]$. You have just solved $x^{2} > x + 2$ graphically: $x < -1$ or $x > 2$. No sign table was needed --- only the two intersection abscissae and a glance at the picture.
\end{experience}

\begin{propriete}[Graphical meaning of $f(x) > g(x)$]
Let $f$ and $g$ be two functions with curves $\mathcal{C}_{f}$ and $\mathcal{C}_{g}$ drawn on the same axes.
\begin{itemize}
\item $f(x) = g(x)$ at the \cle{abscissae of the intersection points} of the two curves;
\item $f(x) > g(x)$ exactly on the set of $x$-values where $\mathcal{C}_{f}$ lies \cle{strictly above} $\mathcal{C}_{g}$;
\item $f(x) < g(x)$ exactly where $\mathcal{C}_{f}$ lies strictly below $\mathcal{C}_{g}$.
\end{itemize}
The solution set is always a set of values of $x$: it is read on the \textbf{$x$-axis}, never on the $y$-axis.
\end{propriete}

\begin{methode}[Solving $f(x) > g(x)$ graphically]
\begin{enumerate}
\item Draw (or use the given) curves $\mathcal{C}_{f}$ and $\mathcal{C}_{g}$ on the same axes.
\item Find the abscissae of the intersection points, either by reading them off the graph or by solving $f(x) = g(x)$ algebraically.
\item On each interval determined by these abscissae, decide which curve is on top.
\item \cle{Project} the relevant portions of the curves onto the $x$-axis and write the solution set as a union of intervals.
\end{enumerate}
\end{methode}

\begin{attention}[Project onto the $x$-axis!]
The most frequent error in graphical questions is to read the answer on the $y$-axis. An inequality in $x$ has a solution set made of \textbf{values of $x$}. After locating the intersection \emph{points}, drop vertical lines down to the $x$-axis and read the intervals there. Writing ``$y > 2$'' as a final answer to ``solve $f(x) > 2$'' costs marks at the GCE.
\end{attention}

\courssub{Graphical Solution of Absolute-Value Inequalities}

\courssub{The curve $y = |f(x)|$}

To solve $|f(x)| < k$ graphically we first need the graph of $y = |f(x)|$. Recall that
\[
|f(x)| = \begin{cases} f(x) & \text{if } f(x) \geq 0, \\ -f(x) & \text{if } f(x) < 0. \end{cases}
\]
So wherever $f$ is non-negative, the two curves coincide; wherever $f$ is negative, taking the modulus changes the sign of the ordinate, which geometrically is a \cle{reflection in the $x$-axis}.

\begin{methode}[Sketching $y = |f(x)|$ from $y = f(x)$]
\begin{enumerate}
\item Sketch $y = f(x)$ and mark its $x$-intercepts.
\item \textbf{Keep unchanged} every part of the curve lying on or above the $x$-axis.
\item \textbf{Reflect in the $x$-axis} every part lying below the $x$-axis (it flips up).
\item The resulting curve never goes below the $x$-axis; at the intercepts it typically forms a sharp ``corner'' (a cusp-like point where the curve is not differentiable).
\end{enumerate}
\end{methode}

\begin{popfigure}
\begin{center}
\begin{tikzpicture}
\begin{axis}[width=6.6cm, height=5.6cm, axis lines=middle, xlabel={$x$}, ylabel={$y$},
xmin=-2.5, xmax=4.5, ymin=-4.5, ymax=5.5, xtick={-1,1,3}, ytick={-4,2,4},
title={$y=f(x)=x^{2}-2x-3$}, title style={font=\small}, domain=-2:4, samples=100]
\addplot[popblue, thick] {x^2-2*x-3};
\addplot[popmuted, dashed, domain=-2:4] {0};
\end{axis}
\end{tikzpicture}\hfill
\begin{tikzpicture}
\begin{axis}[width=6.6cm, height=5.6cm, axis lines=middle, xlabel={$x$}, ylabel={$y$},
xmin=-2.5, xmax=4.5, ymin=-4.5, ymax=5.5, xtick={-1,1,3}, ytick={2,4},
title={$y=|f(x)|$}, title style={font=\small}, domain=-2:4, samples=200]
\addplot[poporange, thick] {abs(x^2-2*x-3)};
\addplot[popblue, dashed, domain=-1:3] {-(x^2-2*x-3)};
\end{axis}
\end{tikzpicture}
\end{center}
\legende{Left: the parabola $y = f(x)$. Right: $y = |f(x)|$; the dashed arc shows the original negative part, which has been reflected upwards in the $x$-axis.}
\end{popfigure}

\courssub{Solving $|f(x)| < k$ and $|f(x)| > k$ graphically}

\begin{propriete}[Modulus equations and horizontal lines]
Let $k > 0$. The solutions of $|f(x)| = k$ are exactly the abscissae of the intersection points of the curve $y = f(x)$ (or equivalently $y = |f(x)|$) with the horizontal lines $y = k$ and $y = -k$. Consequently:
\begin{itemize}
\item $|f(x)| < k$ holds where the curve $y = |f(x)|$ lies \textbf{below} the line $y = k$;
\item $|f(x)| > k$ holds where the curve $y = |f(x)|$ lies \textbf{above} the line $y = k$.
\end{itemize}
(The proof is immediate from $|f(x)| < k \iff -k < f(x) < k$, seen in the previous section; it is not examinable but must be understood.)
\end{propriete}

\begin{exoresolu}[Solving $|x^{2} - 2x - 3| < 2$ graphically]
Using the graph of $y = |x^{2} - 2x - 3|$, solve the inequality $|x^{2} - 2x - 3| < 2$.
\tcblower
\textbf{Step 1: find the intersections with $y = 2$.} We solve $|x^{2} - 2x - 3| = 2$, i.e.
\[
x^{2} - 2x - 3 = 2 \quad\text{or}\quad x^{2} - 2x - 3 = -2.
\]
The first gives $x^{2} - 2x - 5 = 0$, so $x = 1 \pm \sqrt{6}$. The second gives $x^{2} - 2x - 1 = 0$, so $x = 1 \pm \sqrt{2}$. Numerically: $1 - \sqrt{6} \approx -1.45$, $1 - \sqrt{2} \approx -0.41$, $1 + \sqrt{2} \approx 2.41$, $1 + \sqrt{6} \approx 3.45$.

\textbf{Step 2: read the graph.} The modulus curve lies \emph{below} the line $y = 2$ between the first and second intersections, and again between the third and fourth.

\textbf{Step 3: project onto the $x$-axis.}
\[
S = \left(1 - \sqrt{6},\; 1 - \sqrt{2}\right) \cup \left(1 + \sqrt{2},\; 1 + \sqrt{6}\right).
\]
This agrees with the algebraic solution obtained by the double-inequality method --- a good cross-check.
\end{exoresolu}

\begin{popfigure}
\begin{center}
\begin{tikzpicture}
\begin{axis}[width=11.5cm, height=7cm, axis lines=middle, xlabel={$x$}, ylabel={$y$},
xmin=-2.6, xmax=4.6, ymin=-1.2, ymax=5.6,
xtick={-1.449,-0.414,1,2.414,3.449}, xticklabels={$1{-}\sqrt{6}$,$1{-}\sqrt{2}$,$1$,$1{+}\sqrt{2}$,$1{+}\sqrt{6}$},
xticklabel style={font=\scriptsize}, ytick={2,4}, domain=-2.2:4.2, samples=200]
\addplot[poporange, very thick] {abs(x^2-2*x-3)};
\addplot[poppurple, thick, domain=-2.6:4.6] {2};
\addplot[popgreen, very thick] coordinates {(-1.449,0) (-0.414,0)};
\addplot[popgreen, very thick] coordinates {(2.414,0) (3.449,0)};
\addplot[only marks, mark=*, popdark, mark size=1.6pt] coordinates {(-1.449,2) (-0.414,2) (2.414,2) (3.449,2)};
\node[popgreen, font=\small] at (axis cs:-0.93,-0.75) {solution};
\node[popgreen, font=\small] at (axis cs:2.93,-0.75) {solution};
\node[poppurple, font=\small] at (axis cs:4.35,2.35) {$y=2$};
\end{axis}
\end{tikzpicture}
\end{center}
\legende{Graphical solution of $|x^{2}-2x-3|<2$: the curve $y=|f(x)|$ is below $y=2$ exactly above the thick green segments of the $x$-axis.}
\end{popfigure}

\courssub{Comparing two modulus curves: $|f(x)| < |g(x)|$}

The same principle applies with two curves: sketch $y = |f(x)|$ and $y = |g(x)|$ on the same axes, find the abscissae of their intersection points (solve $|f(x)| = |g(x)|$, i.e. $f(x) = g(x)$ or $f(x) = -g(x)$), then read where the first modulus curve is below the second.

\begin{exemplebox}[A quick comparison]
Solve graphically $|x - 1| < |x + 1|$. The two V-shaped curves meet where $x - 1 = -(x + 1)$, i.e. $x = 0$ (the case $x-1 = x+1$ is impossible). For $x > 0$ the line $y = |x - 1|$ lies below $y = |x + 1|$. Hence $S = (0, +\infty)$. Geometrically this says: the points closer to $1$ than to $-1$ are exactly the points of the right half-line --- a fact you also met with perpendicular bisectors in pure geometry.
\end{exemplebox}

\courssub{The curve $y = f(|x|)$ (a useful GCE extra)}

\begin{propriete}[The graph of $y = f(|x|)$]
For $x \geq 0$, $f(|x|) = f(x)$, so the graph coincides with the right half of $\mathcal{C}_{f}$. For $x < 0$, $f(|x|) = f(-x)$, which is the mirror image of the right half in the $y$-axis. Hence $y = f(|x|)$ is always an \cle{even} function: keep the part of $\mathcal{C}_{f}$ with $x \geq 0$, discard the left half, and reflect the kept part in the $y$-axis.
\end{propriete}

\begin{attention}[Do not confuse $|f(x)|$ with $f(|x|)$]
These are two different transformations. $y = |f(x)|$ reflects the \emph{negative ordinates} in the $x$-axis (vertical flip). $y = f(|x|)$ reflects the \emph{right half-graph} in the $y$-axis (horizontal flip) and produces an even curve. A frequent mistake is to reflect the left half of $\mathcal{C}_{f}$ instead of the right half: remember that $f(|x|)$ at $x = -3$ uses the value of $f$ at $+3$, so it is the \textbf{right} half that is copied.
\end{attention}

\courssub{The Complete Curve-Sketching Algorithm}

Every lesson of this chapter now comes together. A full study of a function, as demanded in Paper 2, follows a fixed checklist. Mastering the \emph{order} of the steps is what separates a confident candidate from a lost one.

\begin{methode}[The ten-step complete study of a function]
\begin{enumerate}
\item \textbf{Domain:} find the domain of definition $D_{f}$ (denominators, square roots, logarithms).
\item \textbf{Symmetry:} test $f(-x)$; if $f$ is even or odd, study on half the domain and deduce the rest (axis or centre of symmetry). Also look for other centres/axes via $f(2a - x) = f(x)$ or $f(a+t) + f(a-t) = 2b$.
\item \textbf{Limits at the bounds} of $D_{f}$, including $\pm\infty$.
\item \textbf{Asymptotes:} vertical (infinite limits), horizontal (finite limits at infinity), oblique $y = ax + b$ with $a = \lim f(x)/x$, $b = \lim [f(x) - ax]$.
\item \textbf{Derivative:} compute $f'$, solve $f'(x) = 0$, find stationary points.
\item \textbf{Table of variation:} sign of $f'$, arrows, remarkable values.
\item \textbf{Concavity:} sign of $f''$; points of inflexion where $f''$ changes sign.
\item \textbf{Intercepts:} $y$-intercept $f(0)$; $x$-intercepts from $f(x) = 0$ (exact or approximate).
\item \textbf{Position of the curve} relative to its asymptotes (sign of $f(x) - (ax+b)$) and infinite branches.
\item \textbf{Plot:} draw asymptotes first, then remarkable points, then the curve respecting variation, concavity and position.
\end{enumerate}
\end{methode}

\begin{popfigure}
\begin{center}
\begin{tikzpicture}[node distance=0.32cm,
box/.style={draw=popblue, fill=popblueL, rounded corners=2pt, text width=3.3cm, align=center, font=\small, inner sep=3pt},
boxb/.style={draw=poppurple, fill=white, rounded corners=2pt, text width=3.3cm, align=center, font=\small, inner sep=3pt},
arr/.style={->, thick, popdark}]
\node[box] (n1) {1. Domain $D_f$};
\node[boxb, below=of n1] (n2) {2. Symmetry};
\node[box, below=of n2] (n3) {3. Limits at bounds};
\node[boxb, below=of n3] (n4) {4. Asymptotes};
\node[box, below=of n4] (n5) {5. Derivative $f'$};
\node[boxb, right=1.1cm of n5] (n6) {6. Table of variation};
\node[box, above=of n6] (n7) {7. Concavity, inflexion};
\node[boxb, above=of n7] (n8) {8. Intercepts};
\node[box, above=of n8] (n9) {9. Position vs asymptotes};
\node[boxb, above=of n9] (n10) {10. Plot the curve};
\draw[arr] (n1) -- (n2); \draw[arr] (n2) -- (n3); \draw[arr] (n3) -- (n4);
\draw[arr] (n4) -- (n5); \draw[arr] (n5) -- (n6); \draw[arr] (n6) -- (n7);
\draw[arr] (n7) -- (n8); \draw[arr] (n8) -- (n9); \draw[arr] (n9) -- (n10);
\end{tikzpicture}
\end{center}
\legende{The ten-step curve-sketching algorithm: follow the arrows and no GCE question will surprise you.}
\end{popfigure}

\courssub{Full Worked Study: a Rational Function with an Oblique Asymptote}

\begin{exoresolu}[Complete study of $f(x) = \dfrac{x^{2} - 3x + 3}{x - 2}$]
Let $f(x) = \dfrac{x^{2} - 3x + 3}{x - 2}$ and let $\mathcal{C}$ be its curve.
\begin{enumerate}
\item Study $f$ completely (domain, limits, asymptotes, variation, position).
\item Sketch $\mathcal{C}$.
\item Deduce the number of solutions of $f(x) = m$ according to the values of the real number $m$.
\item Solve graphically $|f(x)| < 1$.
\end{enumerate}
\tcblower
\textbf{1. Domain.} $D_{f} = \mathbb{R} \setminus \{2\}$. No symmetry (the domain is not symmetric about $0$).

\textbf{2. Limits.} $\displaystyle\lim_{x \to 2^{-}} f(x) = -\infty$ and $\lim_{x \to 2^{+}} f(x) = +\infty$ (numerator at $x=2$ equals $1 > 0$). At infinity, $f(x) \sim x$, so $\lim_{x \to \pm\infty} f(x) = \pm\infty$.

\textbf{3. Asymptotes.} The line $x = 2$ is a vertical asymptote. Polynomial division gives
\[
f(x) = x - 1 + \frac{1}{x - 2},
\]
so $\lim_{x \to \pm\infty} [f(x) - (x - 1)] = 0$: the line $y = x - 1$ is an \cle{oblique asymptote}. Moreover $f(x) - (x-1) = \dfrac{1}{x-2}$: the curve is \textbf{above} the asymptote for $x > 2$ and \textbf{below} it for $x < 2$.

\textbf{4. Derivative and variation.}
\[
f'(x) = 1 - \frac{1}{(x-2)^{2}} = \frac{(x-2)^{2} - 1}{(x-2)^{2}} = \frac{(x-1)(x-3)}{(x-2)^{2}}.
\]
$f'(x) = 0$ at $x = 1$ and $x = 3$; $f(1) = -1$ (local maximum) and $f(3) = 3$ (local minimum). Note the amusing fact $f(3) > f(1)$: the two branches are separate.
\begin{center}
\begin{tabular}{|c|ccccccccc|}
\hline
$x$ & $-\infty$ & & $1$ & & $2$ & & $3$ & & $+\infty$ \\ \hline
$f'(x)$ & & $+$ & $0$ & $-$ & $\|$ & $-$ & $0$ & $+$ & \\ \hline
$f(x)$ & $-\infty$ & $\nearrow$ & $-1$ & $\searrow$ & $\|$ & $\searrow$ & $3$ & $\nearrow$ & $+\infty$ \\ \hline
\end{tabular}
\end{center}
(On $(3, +\infty)$, $f' > 0$ and $f$ rises from $3$ to $+\infty$; on $(-\infty,1)$, $f$ rises from $-\infty$ to $-1$.)

\textbf{5. Concavity.} $f''(x) = \dfrac{2}{(x-2)^{3}}$: concave down on $(-\infty, 2)$, concave up on $(2, +\infty)$; no inflexion point since $f''$ never vanishes (the concavity changes across the asymptote, not at a point of the curve).

\textbf{6. Intercepts.} $f(0) = -\tfrac{3}{2}$. The numerator $x^{2} - 3x + 3$ has discriminant $9 - 12 = -3 < 0$: \textbf{no $x$-intercept}.

\textbf{7. Sketch.} Draw the asymptotes $x = 2$ and $y = x - 1$, plot $(1, -1)$, $(3, 3)$, $(0, -1.5)$, then the two branches respecting the table and the position relative to the oblique asymptote.

\textbf{8. Number of solutions of $f(x) = m$.} Reading horizontal lines $y = m$ across the sketch:
\begin{itemize}
\item if $-1 < m < 3$: no solution;
\item if $m = -1$ or $m = 3$: exactly one solution;
\item if $m < -1$ or $m > 3$: exactly two solutions.
\end{itemize}

\textbf{9. Solving $|f(x)| < 1$.} Reflect the negative part of $\mathcal{C}$ in the $x$-axis to get $y = |f(x)|$, then cut with $y = 1$: $|f(x)| = 1 \iff f(x) = 1$ or $f(x) = -1$. From $f(x) = 1$: $x^{2} - 3x + 3 = x - 2$, i.e. $x^{2} - 4x + 5 = 0$, no real root. From $f(x) = -1$: $x^{2} - 3x + 3 = -(x - 2)$, i.e. $x^{2} - 2x + 1 = 0$, so $x = 1$ (double). The curve $y = |f(x)|$ only \emph{touches} $y = 1$ at $x = 1$ and stays below it nowhere else... checking the variation: on $(-\infty, 2)$, $|f|$ ranges over $[1, +\infty)$ except near $x=1$ where $|f(1)| = 1$; hence $|f(x)| < 1$ has \textbf{no solution}: $S = \varnothing$. The graph confirms it instantly.
\end{exoresolu}

\begin{attention}[Two classic traps in full studies]
\begin{itemize}
\item Forgetting the \textbf{position of the curve relative to the oblique asymptote}: the sign of $f(x) - (ax + b)$ is a one-line computation and is explicitly demanded by many mark schemes.
\item Concluding ``no inflexion point'' too quickly, or claiming one at $x = 2$ here: an inflexion point must be a point \emph{of the curve}; $f$ is not even defined at $x = 2$.
\end{itemize}
\end{attention}

\begin{savaistu}[Why Paper 2 loves this question]
A complete function study tests, in a single item, limits, differentiation, asymptotes, inequalities and graph reading --- nearly the whole chapter. Cameroonian candidates who rehearse the ten-step checklist on three or four functions before the examination routinely turn this long question into their highest-scoring one. Treat the checklist as a ritual: same order, same table, same care with the axes, every time.
\end{savaistu}

\begin{exercice}[Practice]
\begin{enumerate}
\item Using the graph of $y = |x^{2} - 2x - 3|$ drawn in this section, solve $|x^{2} - 2x - 3| > 2$.
\item From the graph of $y = x^{3} - 3x$, sketch $y = |x^{3} - 3x|$ and $y = |x|^{3} - 3|x|$, stating precisely which parts are kept or reflected.
\item Carry out the complete study of $g(x) = \dfrac{x^{2} + x - 2}{x - 1}$ (oblique asymptote $y = x + 2$), then deduce the number of solutions of $g(x) = m$ and solve graphically $|g(x)| < 2$.
\end{enumerate}
\end{exercice}

\begin{corrige}[Exercise 1]
\begin{enumerate}
\item The curve is above $y = 2$ outside the intervals found in the worked example: $S = (-\infty, 1-\sqrt{6}) \cup (1-\sqrt{2}, 1+\sqrt{2}) \cup (1+\sqrt{6}, +\infty)$.
\item For $y = |x^{3} - 3x|$: keep the parts with $x \in [-\sqrt{3}, 0] \cup [\sqrt{3}, +\infty)$, reflect the rest in the $x$-axis. For $y = |x|^{3} - 3|x| = f(|x|)$: keep the right half ($x \geq 0$) of the cubic and reflect it in the $y$-axis; the result is even, with minima at $x = \pm 1$.
\item $D_{g} = \mathbb{R}\setminus\{1\}$; $g(x) = x + 2$ exactly (since $x^{2} + x - 2 = (x-1)(x+2)$)! The curve is the line $y = x + 2$ with the single point $(1, 3)$ removed --- a ``hole'', not an asymptote. $g(x) = m$ has one solution for $m \neq 3$, none for $m = 3$. $|g(x)| < 2 \iff -2 < x + 2 < 2$ with $x \neq 1$: $S = (-4, 0)$.
\end{enumerate}
\end{corrige}

\begin{aretenir}[Key points of the section]
\begin{itemize}
\item $f(x) > g(x)$ where $\mathcal{C}_{f}$ is above $\mathcal{C}_{g}$; solutions are read on the \textbf{$x$-axis}.
\item $y = |f(x)|$: keep the upper part, reflect the lower part in the $x$-axis. $y = f(|x|)$: keep the right half, reflect it in the $y$-axis (even function).
\item $|f(x)| = k$ at the intersections with $y = k$ and $y = -k$; then read above/below.
\item The ten-step checklist (domain, symmetry, limits, asymptotes, $f'$, variation, $f''$, intercepts, position, plot) is the backbone of every Paper 2 curve-sketching question.
\item A finished sketch is a tool: use it to count roots, solve inequalities and discuss positions.
\end{itemize}
\end{aretenir}

\newpage

\coursec{Integration activity}

\courssub{Aims of the activity}

This integration activity closes the chapter by mobilising, in one authentic optimisation problem, every skill developed in Lessons 66 to 76. Working in teams, you will:

\begin{itemize}
\item determine the \cle{domain of definition} of a rational-type function and compute its \cle{limits at the bounds} of that domain;
\item prove the existence of an \cle{oblique asymptote} and study the \cle{position of the curve} relative to it (infinite branches);
\item compute a derivative, discuss \cle{stationary points}, and build a complete \cle{table of variation};
\item study \cle{concavity} through the sign of the second derivative and decide whether a \cle{point of inflexion} exists;
\item solve an \cle{absolute-value inequality} \emph{both} algebraically and graphically, and interpret the solution set in context;
\item produce a complete, well-annotated \cle{curve sketch} and defend it before the class.
\end{itemize}

\begin{aretenir}[Skills being assessed]
Domain $\to$ limits $\to$ asymptotes $\to$ derivative and variation $\to$ concavity $\to$ sketch $\to$ inequality solving. This is exactly the chain of reasoning expected in the GCE Advanced Level Further Mathematics examination: a full function study is never a list of disconnected calculations, but a single coherent argument in which each step feeds the next.
\end{aretenir}

\courssub{Situation}

\textbf{Optimising a Water Tank's Profit Curve in Buea.}\par
Madam Eposi runs a water-vending business near Molyko, Buea. Every day she fills her tank from the Cameroon Water Utilities network and resells water by the jerrycan to students of the University of Buea. After several months of record keeping, her elder brother, a Further Mathematics student, models her \cle{daily profit} $P(x)$, measured in \emph{thousands of FCFA}, as a function of the quantity $x$ she sells, measured in \emph{hundreds of litres}:
\[
P(x)=\frac{-x^{2}+10x+7}{x+1},\qquad x\ge 0.
\]
The model reflects two realities of the business: moderate sales bring good margins, but beyond a certain level the extra costs (hiring an extra carrier, buying fuel for the pump during power cuts, discounts to attract hostels) eat into the profit, which eventually falls.

Madam Eposi wants answers to three practical questions:
\begin{itemize}
\item What quantity of water maximises the daily profit, and what is that maximum profit?
\item For large sales, the profit decreases almost along a straight line: which line, and does the profit curve lie above or below it?
\item For peace of mind, she accepts any sales level for which the profit stays \emph{within 2 thousand FCFA of the maximum}. What is the corresponding range of daily sales?
\end{itemize}

Your team is hired as the consulting mathematicians. You must produce a poster containing the full study of the function $P$ and a clear, justified answer to each question.

\courssub{Tasks}

\begin{enumerate}
\item \textbf{Domain and bounds.} Justify that the natural domain of $P$ is $\mathbb{R}\setminus\{-1\}$ and state the domain of the \emph{model}. Compute $P(0)$ and $\displaystyle\lim_{x\to+\infty}P(x)$. Interpret each result for the business.
\item \textbf{Oblique asymptote and position.} Show that for all $x\ge 0$,
\[
P(x)=-x+11-\frac{4}{x+1}.
\]
Deduce that the line $(\Delta):\ y=-x+11$ is an oblique asymptote to the curve $(\mathcal{C})$ of $P$, and determine the position of $(\mathcal{C})$ relative to $(\Delta)$.
\item \textbf{Variation and maximum.} Compute $P'(x)$ and show that $P'(x)=\dfrac{4}{(x+1)^{2}}-1$. Solve $P'(x)=0$ on $[0,+\infty[$, build the table of variation of $P$, and deduce the maximum profit and the quantity sold that achieves it.
\item \textbf{Concavity.} Compute $P''(x)$, study its sign on $[0,+\infty[$, state the concavity of $(\mathcal{C})$ and decide whether $(\mathcal{C})$ has a point of inflexion.
\item \textbf{Intercepts and sketch.} Find the intercepts of $(\mathcal{C})$ with the axes (give surd form where necessary). Sketch $(\mathcal{C})$, the asymptote $(\Delta)$, the maximum point and the intercepts on the same axes.
\item \textbf{The tolerance band.} Let $P_{\max}$ denote the maximum profit. Solve \emph{algebraically} the inequality
\[
\left|P(x)-P_{\max}\right|\le 2 \qquad (x\ge 0),
\]
then verify your solution set \emph{graphically} by drawing the horizontal lines $y=P_{\max}-2$ and $y=P_{\max}+2$ on your sketch. Interpret the solution set for Madam Eposi.
\item \textbf{Poster and critique.} Present your sketch and solution set on a poster. Each team checks another team's poster against the marking scheme: correct domain, correct asymptote with position, correct table of variation, correct concavity, correct solution set.
\end{enumerate}

\courssub{Model answer}

\textbf{1. Domain and bounds.}\par
$P$ is a rational function (quotient of two polynomials); it is defined unless its denominator vanishes, i.e. unless $x+1=0$, that is $x=-1$. The \cle{natural domain} is therefore $\mathbb{R}\setminus\{-1\}$. Since a quantity of water sold cannot be negative, the \emph{domain of the model} is $D=[0,+\infty[$. Note that the forbidden value $x=-1$ lies outside the model's domain, so no vertical asymptote is relevant to the business study.

At the left bound of $D$:
\[
P(0)=\frac{-0+0+7}{0+1}=7.
\]
At the right bound, factorise the highest powers of $x$ (standard technique for limits of rational functions at infinity): for $x>0$,
\[
P(x)=\frac{x^{2}\left(-1+\dfrac{10}{x}+\dfrac{7}{x^{2}}\right)}{x\left(1+\dfrac{1}{x}\right)}
=x\cdot\frac{-1+\dfrac{10}{x}+\dfrac{7}{x^{2}}}{1+\dfrac{1}{x}}.
\]
As $x\to+\infty$, the bracketed fraction tends to $\dfrac{-1}{1}=-1$, while the factor $x$ tends to $+\infty$. Hence
\[
\lim_{x\to+\infty}P(x)=-\infty .
\]
\emph{Interpretation:} with no sales, a fixed income (for instance renting out empty cans) yields $7$ thousand FCFA, i.e. \SI{7000}{FCFA}; excessively large sales drive the profit down without bound because of the extra costs. The curve therefore has an \cle{infinite branch} at $+\infty$, whose nature is determined in Task 2.

\textbf{2. Oblique asymptote and position.}\par
We perform the polynomial division of $-x^{2}+10x+7$ by $x+1$. Since
\[
(x+1)(-x+11)=-x^{2}+11x-x+11=-x^{2}+10x+11,
\]
we obtain
\[
-x^{2}+10x+7=(x+1)(-x+11)-4.
\]
Dividing both sides by $x+1$ (valid for $x\neq -1$, in particular for all $x\ge 0$):
\[
P(x)=-x+11-\frac{4}{x+1}.
\]
Now examine the gap between the curve and the line $(\Delta):\ y=-x+11$:
\[
P(x)-(-x+11)=-\frac{4}{x+1}\xrightarrow[x\to+\infty]{}0.
\]
By definition, $(\Delta)$ is an \cle{oblique asymptote} of $(\mathcal{C})$ at $+\infty$: this is the ``straight line'' Madam Eposi observed.

For the \cle{position of the curve} relative to the asymptote, study the sign of the gap: for all $x\ge 0$, $x+1>0$, so
\[
P(x)-(-x+11)=-\frac{4}{x+1}<0.
\]
Therefore $(\mathcal{C})$ lies strictly \textbf{below} $(\Delta)$ on the whole domain: the profit approaches the line from underneath, never reaching it.

\textbf{3. Variation and maximum.}\par
Differentiate the divided form $P(x)=-x+11-4(x+1)^{-1}$ term by term:
\[
P'(x)=-1-4\cdot\bigl(-(x+1)^{-2}\bigr)=-1+\frac{4}{(x+1)^{2}}=\frac{4-(x+1)^{2}}{(x+1)^{2}}.
\]
Factorising the numerator as a difference of two squares, $4-(x+1)^{2}=\bigl(2-(x+1)\bigr)\bigl(2+(x+1)\bigr)$:
\[
P'(x)=\frac{(1-x)(x+3)}{(x+1)^{2}}.
\]
On $[0,+\infty[$, $(x+1)^{2}>0$ and $x+3>0$, so $P'(x)$ has the sign of $1-x$. The equation $P'(x)=0$ reduces to $1-x=0$, i.e. $x=1$ (the root $x=-3$ is outside the domain). Thus $P'(x)>0$ on $[0,1[$, $P'(1)=0$, and $P'(x)<0$ on $]1,+\infty[$. The only \cle{stationary point} is at $x=1$, and
\[
P(1)=-1+11-\frac{4}{2}=8.
\]

\begin{center}
\begin{tabular}{|c|ccccc|}
\hline
$x$ & $0$ & & $1$ & & $+\infty$ \\
\hline
$P'(x)$ & & $+$ & $0$ & $-$ & \\
\hline
$P(x)$ & $7$ & $\nearrow$ & $8$ & $\searrow$ & $-\infty$ \\
\hline
\end{tabular}
\end{center}

Since $P'$ changes sign from $+$ to $-$ at $x=1$, the stationary point is a \cle{maximum}. \emph{Interpretation:} the maximum daily profit is $P_{\max}=8$ thousand FCFA, i.e. \SI{8000}{FCFA}, obtained by selling $x=1$ hundred litres, that is \SI{100}{litres} per day.

\textbf{4. Concavity.}\par
Differentiate $P'(x)=-1+4(x+1)^{-2}$:
\[
P''(x)=4\cdot\bigl(-2(x+1)^{-3}\bigr)=-\frac{8}{(x+1)^{3}}.
\]
For all $x\ge 0$, $(x+1)^{3}>0$, so $P''(x)<0$: the curve is \cle{concave down} on the whole domain (its tangents all lie above it). Since $P''$ never vanishes and never changes sign on $[0,+\infty[$, $(\mathcal{C})$ has \textbf{no point of inflexion}. This is consistent with Task 3: a function that is concave down around a stationary point has a maximum there.

\textbf{5. Intercepts and sketch.}\par
$y$-intercept: $P(0)=7$, giving the point $(0,7)$.

$x$-intercepts: solve $P(x)=0$, i.e. $-x^{2}+10x+7=0$, equivalently $x^{2}-10x-7=0$. The discriminant is $\Delta=100+28=128=(8\sqrt{2})^{2}$, so
\[
x=\frac{10\pm 8\sqrt{2}}{2}=5\pm 4\sqrt{2}.
\]
Only $x=5+4\sqrt{2}\approx 10.66$ lies in the domain $[0,+\infty[$ (the other root $5-4\sqrt{2}\approx -0.66$ is rejected). \emph{Interpretation:} the profit becomes negative beyond about \SI{1066}{litres} sold per day.

\begin{popfigure}
\begin{tikzpicture}
\begin{axis}[
  width=0.86\linewidth, height=7.2cm,
  axis lines=middle, xlabel={$x$ (hundred litres)}, ylabel={$P(x)$},
  xmin=-0.5, xmax=11.5, ymin=-6, ymax=13,
  xtick={0,1,2,4.236,5,10.66}, xticklabels={$0$,$1$,$2$,$2{+}\sqrt{5}$,$5$,$5{+}4\sqrt{2}$},
  ytick={6,7,8,10}, grid=both, grid style={popline, very thin},
  legend style={font=\small, at={(0.98,0.55)}, anchor=east}
]
\addplot[popblue, very thick, domain=0:11.2, samples=120] {-x+11-4/(x+1)};
\addlegendentry{$y=P(x)$}
\addplot[poporange, dashed, thick, domain=0:11.2] {-x+11};
\addlegendentry{asymptote $y=-x+11$}
\addplot[poppurple, dashed, domain=0:11.2] {6};
\addplot[poppurple, dashed, domain=0:11.2] {10};
\addplot[only marks, mark=*, popdark] coordinates {(1,8) (0,7) (10.66,0) (4.236,6)};
\node[popdark, anchor=south west, font=\small] at (axis cs:1,8) {max $(1,8)$};
\end{axis}
\end{tikzpicture}
\legende{The profit curve, its oblique asymptote, and the tolerance band $6\le y\le 10$.}
\end{popfigure}

\textbf{6. The tolerance band.}\par
\emph{Algebraic solution.} With $P_{\max}=8$, the key equivalence for absolute values, $|u|\le a\iff -a\le u\le a$, gives
\[
|P(x)-8|\le 2\iff -2\le P(x)-8\le 2\iff 6\le P(x)\le 10.
\]
The table of variation shows $P(x)\le 8<10$ for every $x\ge 0$, so the right-hand inequality $P(x)\le 10$ is automatically satisfied. Only $P(x)\ge 6$ remains:
\[
-x+11-\frac{4}{x+1}\ge 6
\iff 5-x\ge \frac{4}{x+1}.
\]
Since $x\ge 0$, we have $x+1>0$, so multiplying both sides by $x+1$ preserves the inequality:
\[
(5-x)(x+1)\ge 4
\iff -x^{2}+4x+5\ge 4
\iff x^{2}-4x-1\le 0.
\]
The quadratic $x^{2}-4x-1$ has discriminant $16+4=20$ and roots
\[
x=\frac{4\pm\sqrt{20}}{2}=2\pm\sqrt{5}.
\]
A quadratic with positive leading coefficient is negative \emph{between} its roots, so
\[
x\in[\,2-\sqrt{5},\;2+\sqrt{5}\,].
\]
Finally, intersect with the domain $x\ge 0$; since $2-\sqrt{5}\approx -0.24<0$:
\[
\boxed{\,S=[\,0,\;2+\sqrt{5}\,]\approx[\,0,\;4.24\,]\,}
\]
\emph{Graphical verification.} The horizontal lines $y=P_{\max}-2=6$ and $y=P_{\max}+2=10$ delimit a band. On the sketch, the curve enters the band at $x=0$ (since $P(0)=7$) and stays inside it until it crosses the line $y=6$ at $x=2+\sqrt{5}\approx 4.24$; beyond that it leaves the band from below. The graphical reading matches the algebraic solution set exactly.

\emph{Interpretation:} Madam Eposi's profit stays within \SI{2000}{FCFA} of her maximum as long as daily sales do not exceed about $424$ litres; beyond that level, the profit drops more than \SI{2000}{FCFA} below the best possible value.

\begin{attention}[Common errors observed during the critique session]
\begin{itemize}
\item Forgetting to intersect the algebraic solution $[2-\sqrt{5},\,2+\sqrt{5}]$ with the domain $x\ge 0$ of the model.
\item Claiming an inflexion point ``because there is a maximum'': a maximum comes from $P'$ changing sign, an inflexion point from $P''$ changing sign. Here $P''(x)<0$ throughout the domain.
\item Writing the asymptote as $y=-x+11$ ``by looking at the leading terms'' without computing $P(x)-(-x+11)$ and showing it tends to $0$: this justification is examinable and carries marks.
\item Multiplying an inequality by $x+1$ without first noting that $x+1>0$ on the domain, so the direction of the inequality is preserved.
\item Solving $|P(x)-8|\le 2$ as only $P(x)\ge 6$ without first checking that $P(x)\le 10$ holds automatically.
\end{itemize}
\end{attention}

\begin{savaistu}[Why a tolerance band?]
In quality control and business management, decisions are rarely based on an exact optimum: one accepts a \emph{tolerance interval} around it. Writing ``within $a$ units of the target $m$'' is precisely the inequality $|x-m|\le a$, which is why absolute-value inequalities appear in economics, engineering and the GCE examination alike.
\end{savaistu}

\newpage

\coursec{Methods and worked examples}

\coursec{Further Curve Sketching and Inequalities}

\courssub{Domain of Definition and Limits at the Bounds}

\begin{definition}[Domain of Definition]
The \cle{domain of definition} $D_f$ of a function $f$ is the set of all real numbers $x$ for which the expression $f(x)$ makes sense (i.e. can be computed without violating the rules of real arithmetic).
\end{definition}

\begin{propriete}[Common Restrictions on $D_f$]
For real-valued functions of a real variable:
\begin{itemize}
\item \textbf{Denominator:} $q(x) \neq 0$ for $f(x) = \frac{p(x)}{q(x)}$.
\item \textbf{Even root:} $u(x) \geq 0$ for $f(x) = \sqrt[n]{u(x)}$ with $n$ even.
\item \textbf{Logarithm:} $u(x) > 0$ for $f(x) = \ln u(x)$ or $\log_a u(x)$.
\item \textbf{Tangent:} $u(x) \neq \frac{\pi}{2} + k\pi,\ k \in \mathbb{Z}$ for $f(x) = \tan u(x)$.
\item \textbf{Arcsine/Arccosine:} $-1 \leq u(x) \leq 1$ for $f(x) = \arcsin u(x)$ or $\arccos u(x)$.
\end{itemize}
$D_f$ is the \textbf{intersection} of all individual conditions, expressed as a union of intervals.
\end{propriete}

\begin{methode}[Finding the Domain and Limits at its Bounds]
\begin{enumerate}
\item \textbf{Identify the constraints.} Ask: what could make $f(x)$ impossible to compute? (See the list above.)
\item \textbf{Solve each constraint} as an equation or inequality in $x$, then \textbf{intersect} all the conditions. Write $D_f$ as a union of intervals.
\item \textbf{Compute the limits at every bound} of $D_f$:
\begin{itemize}
\item At $\pm\infty$: compare degrees for rational functions; factorise the dominant term.
\item At each excluded finite value $a$: compute $\lim_{x\to a^-} f(x)$ and $\lim_{x\to a^+} f(x)$ separately (sign of the denominator on each side).
\end{itemize}
\item \textbf{Interpret:} an infinite limit at a finite bound $x=a$ announces a \cle{vertical asymptote}; a finite limit at $\pm\infty$ announces a \cle{horizontal asymptote}.
\end{enumerate}
\textbf{Reflex:} never write $\lim_{x\to a}$ when $a$ is excluded from the domain — always use one-sided limits.
\end{methode}

\begin{exemplebox}[Domain and limits of a rational function with a radical]
Let $f(x) = \dfrac{\sqrt{x+1}}{x-2}$.
\begin{enumerate}
\item Determine the domain of definition $D_f$ of $f$.
\item Compute the limits of $f$ at the bounds of $D_f$ and interpret the results graphically.
\end{enumerate}
\tcblower
\textbf{1. Domain of definition.}

Two constraints act simultaneously:
\begin{itemize}
\item The square root requires $x + 1 \geq 0$, i.e. $x \geq -1$.
\item The denominator requires $x - 2 \neq 0$, i.e. $x \neq 2$.
\end{itemize}
Both conditions must hold, so
\[ D_f = [-1, 2) \cup (2, +\infty). \]

\textbf{2. Limits at the bounds.}

\emph{At $x = -1$ (included bound):} we simply evaluate:
\[ \lim_{x \to -1^+} f(x) = \frac{\sqrt{0}}{-1-2} = 0. \]
The curve starts at the point $(-1, 0)$.

\emph{At $x = 2$ (excluded bound):} the numerator tends to $\sqrt{3} > 0$; the denominator tends to $0$ with a sign depending on the side:
\begin{itemize}
\item If $x \to 2^-$, then $x - 2 < 0$, so $f(x) \to \dfrac{\sqrt{3}}{0^-} = -\infty$.
\item If $x \to 2^+$, then $x - 2 > 0$, so $f(x) \to \dfrac{\sqrt{3}}{0^+} = +\infty$.
\end{itemize}
\textbf{Interpretation:} the line $x = 2$ is a \cle{vertical asymptote}.

\emph{At $+\infty$:} for $x > 0$, write $\sqrt{x+1} = \sqrt{x}\sqrt{1 + \tfrac{1}{x}}$, so
\[ f(x) = \frac{\sqrt{x}\,\sqrt{1+\frac1x}}{x\left(1-\frac2x\right)} = \frac{1}{\sqrt{x}}\cdot\frac{\sqrt{1+\frac1x}}{1-\frac2x} \xrightarrow[x\to+\infty]{} 0 \times \frac{1}{1} = 0. \]
\textbf{Interpretation:} the line $y = 0$ (the $x$-axis) is a \cle{horizontal asymptote} at $+\infty$.
\end{exemplebox}

\begin{aretenir}[Key Points: Domain and Limits]
\begin{itemize}
\item $D_f$ is an intersection of conditions; write it as a union of intervals.
\item At a finite bound $a \notin D_f$, always compute \textbf{one-sided limits} $\lim_{x\to a^-}$ and $\lim_{x\to a^+}$.
\item $\lim_{x\to a^\pm} f(x) = \pm\infty \implies$ vertical asymptote $x = a$.
\item $\lim_{x\to\pm\infty} f(x) = b \in \mathbb{R} \implies$ horizontal asymptote $y = b$.
\end{itemize}
\end{aretenir}

\courssub{Asymptotes, Intercepts and Infinite Branches}

\begin{propriete}[Types of Asymptotes]
Let $(\mathcal{C})$ be the curve of $f$.
\begin{itemize}
\item \textbf{Vertical asymptote} $x = a$: $\lim_{x\to a^\pm} f(x) = \pm\infty$ (or $\mp\infty$).
\item \textbf{Horizontal asymptote} $y = b$: $\lim_{x\to\pm\infty} f(x) = b$ (finite).
\item \textbf{Oblique asymptote} $y = ax + b$ ($a \neq 0$): $\lim_{x\to\pm\infty} \big[f(x) - (ax+b)\big] = 0$.
\end{itemize}
For a rational function $f(x) = \frac{P(x)}{Q(x)}$:
\begin{itemize}
\item $\deg P < \deg Q \implies y = 0$ horizontal asymptote.
\item $\deg P = \deg Q \implies y = \frac{\text{leading coeff of }P}{\text{leading coeff of }Q}$ horizontal asymptote.
\item $\deg P = \deg Q + 1 \implies$ oblique asymptote found by polynomial long division.
\item $\deg P \ge \deg Q + 2 \implies$ no oblique asymptote; \cle{parabolic branch} (study $\lim \frac{f(x)}{x}$).
\end{itemize}
\end{propriete}

\begin{methode}[Asymptotes, Intercepts and Position Relative to Asymptotes]
\begin{enumerate}
\item \textbf{Vertical asymptotes:} at each excluded value $x = a$, if $\lim_{x\to a^{\pm}} f(x) = \pm\infty$, then $x = a$ is a vertical asymptote.
\item \textbf{Horizontal asymptotes:} if $\lim_{x\to\pm\infty} f(x) = b$ (finite), then $y = b$ is a horizontal asymptote at that end.
\item \textbf{Oblique asymptotes} (rational function with $\deg(\text{num}) = \deg(\text{den}) + 1$): perform \textbf{polynomial long division} to write $f(x) = ax + b + \dfrac{r(x)}{q(x)}$ with $\deg r < \deg q$. Since $\dfrac{r(x)}{q(x)} \to 0$, the line $y = ax + b$ is an oblique asymptote.
\item \textbf{Position of the curve relative to the line:} study the \textbf{sign of the remainder} $f(x) - (ax+b) = \dfrac{r(x)}{q(x)}$. Positive $\Rightarrow$ curve above; negative $\Rightarrow$ curve below.
\item \textbf{Parabolic branches:} if $\lim_{x\to\pm\infty} f(x) = \pm\infty$ and no oblique asymptote, compute $\lim_{x\to\pm\infty} \dfrac{f(x)}{x}$:
\begin{itemize}
\item If $0$: parabolic branch in the direction of $(Ox)$.
\item If $\pm\infty$: parabolic branch in the direction of $(Oy)$.
\end{itemize}
\item \textbf{Intercepts:} $y$-intercept $= f(0)$ (if $0 \in D_f$); $x$-intercepts: solve $f(x) = 0$ (numerator $=0$ for a rational function, provided the denominator is non-zero).
\end{enumerate}
\end{methode}

\begin{exemplebox}[Oblique asymptote and relative position]
Let $f(x) = \dfrac{x^2 - 3x + 5}{x - 1}$ and let $(\mathcal{C})$ be its curve.
\begin{enumerate}
\item Show that the line $(\Delta): y = x - 2$ is an asymptote to $(\mathcal{C})$ at $+\infty$ and at $-\infty$.
\item Study the position of $(\mathcal{C})$ relative to $(\Delta)$.
\item Give all asymptotes and the intercepts of $(\mathcal{C})$.
\end{enumerate}
\tcblower
\textbf{1. Oblique asymptote.} Divide $x^2 - 3x + 5$ by $x - 1$:
\[ x^2 - 3x + 5 = (x-1)(x-2) + 3. \]
(Check: $(x-1)(x-2) = x^2 - 3x + 2$, and $x^2 - 3x + 2 + 3 = x^2 - 3x + 5$. \checkmark) Hence
\[ f(x) = x - 2 + \frac{3}{x-1}. \]
Since $\lim_{x\to\pm\infty} \dfrac{3}{x-1} = 0$, we get $\lim_{x\to\pm\infty} \big[f(x) - (x-2)\big] = 0$: the line $y = x - 2$ is an \cle{oblique asymptote} at both $+\infty$ and $-\infty$.

\textbf{2. Relative position.} The sign of the gap is the sign of
\[ f(x) - (x-2) = \frac{3}{x-1}. \]
\begin{itemize}
\item If $x > 1$: $\dfrac{3}{x-1} > 0$, so $(\mathcal{C})$ is \textbf{above} $(\Delta)$.
\item If $x < 1$: $\dfrac{3}{x-1} < 0$, so $(\mathcal{C})$ is \textbf{below} $(\Delta)$.
\end{itemize}

\textbf{3. All asymptotes and intercepts.}

$D_f = \mathbb{R}\setminus\{1\}$. At $x = 1$: the numerator equals $1 - 3 + 5 = 3 > 0$, and
\[ \lim_{x\to 1^-} f(x) = -\infty, \qquad \lim_{x\to 1^+} f(x) = +\infty, \]
so $x = 1$ is a \cle{vertical asymptote}.

$y$-intercept: $f(0) = \dfrac{5}{-1} = -5$; the curve passes through $(0, -5)$.

$x$-intercepts: solve $x^2 - 3x + 5 = 0$; $\Delta = 9 - 20 = -11 < 0$, so there is \textbf{no} $x$-intercept (the curve never crosses the $x$-axis; indeed $x^2 - 3x + 5 > 0$ for all $x$).
\end{exemplebox}

\begin{attention}[Common Mistake: Oblique Asymptote]
Do \textbf{not} confuse the quotient of the division with the asymptote equation. If $f(x) = \frac{P(x)}{Q(x)}$ and $P(x) = Q(x)(ax+b) + R(x)$, then the oblique asymptote is $y = ax + b$, \textbf{not} $y = \frac{P(x)}{Q(x)}$. Always verify that $\deg R < \deg Q$.
\end{attention}

\courssub{Stationary Points and Variation of a Function}

\begin{definition}[Stationary (Critical) Point]
A point $x_0 \in D_f$ is a \cle{stationary point} (or \cle{critical point}) of $f$ if $f'(x_0) = 0$ or $f'(x_0)$ does not exist (while $f$ is continuous at $x_0$).
\end{definition}

\begin{propriete}[First Derivative Test]
Let $x_0$ be a stationary point of $f$ and suppose $f'$ changes sign at $x_0$.
\begin{itemize}
\item $f'$ changes from $+$ to $-$: $f$ has a \cle{local maximum} at $x_0$.
\item $f'$ changes from $-$ to $+$: $f$ has a \cle{local minimum} at $x_0$.
\item $f'$ does \textbf{not} change sign: $x_0$ is a \cle{stationary point of inflexion} (horizontal tangent crossing the curve).
\end{itemize}
\end{propriete}

\begin{methode}[Stationary Points and Table of Variation]
\begin{enumerate}
\item \textbf{Differentiate} $f$ on each interval of $D_f$ (quotient rule for rational functions: $\left(\tfrac{u}{v}\right)' = \tfrac{u'v - uv'}{v^2}$).
\item \textbf{Solve $f'(x) = 0$}: the solutions, together with points where $f'$ fails to exist (but $f$ is continuous), are the \cle{stationary points}.
\item \textbf{Sign of $f'$:} for a rational function, the denominator is a square ($>0$ on $D_f$), so the sign of $f'$ is the sign of its numerator. Build a sign table.
\item \textbf{Table of variation:} one column per remarkable value (bounds of $D_f$, roots of $f'$) \emph{and} one column between consecutive values; place the limits at the bounds, the extrema values $f(x_i)$, and arrows $\nearrow$ / $\searrow$.
\item \textbf{Classify:} $f'$ changes $+ \to -$: local maximum; $- \to +$: local minimum; no sign change: stationary point of inflexion.
\end{enumerate}
\textbf{Reflex:} compute each extremum value $f(x_i)$ exactly — it goes into the table and onto the graph.
\end{methode}

\begin{exemplebox}[Full variation study of a rational function]
Let $f(x) = \dfrac{x^2}{x - 1}$.
\begin{enumerate}
\item Compute $f'(x)$ and solve $f'(x) = 0$.
\item Draw the complete table of variation of $f$ on its domain.
\end{enumerate}
\tcblower
$D_f = \mathbb{R}\setminus\{1\}$.

\textbf{1. Derivative.} With $u = x^2$, $v = x - 1$:
\[ f'(x) = \frac{2x(x-1) - x^2 \cdot 1}{(x-1)^2} = \frac{2x^2 - 2x - x^2}{(x-1)^2} = \frac{x^2 - 2x}{(x-1)^2} = \frac{x(x-2)}{(x-1)^2}. \]
Since $(x-1)^2 > 0$ on $D_f$, $f'(x)$ has the sign of $x(x-2)$:
\[ f'(x) = 0 \iff x = 0 \ \text{or}\ x = 2. \]
Sign of $x(x-2)$: positive on $(-\infty, 0)$, negative on $(0, 2)$ (excluding $x=1$ where $f'$ is undefined), positive on $(2, +\infty)$.

Extremum values: $f(0) = 0$ and $f(2) = \dfrac{4}{1} = 4$.

Limits at the bounds: $\lim_{x\to-\infty} f(x) = -\infty$, $\lim_{x\to+\infty} f(x) = +\infty$ (degree of numerator exceeds degree of denominator by one); at $x = 1$, numerator $= 1 > 0$: $\lim_{x\to1^-} f = -\infty$, $\lim_{x\to1^+} f = +\infty$.

\textbf{2. Table of variation.}
\begin{center}
\begin{tabular}{|c|ccccccccc|}
\hline
$x$ & $-\infty$ & & $0$ & & $1$ & & $2$ & & $+\infty$ \\
\hline
$f'(x)$ & & $+$ & $0$ & $-$ & $\|$ & $-$ & $0$ & $+$ & \\
\hline
$f(x)$ & $-\infty$ & $\nearrow$ & $0$ & $\searrow$ & $-\infty\,\|\,+\infty$ & $\searrow$ & $4$ & $\nearrow$ & $+\infty$ \\
\hline
\end{tabular}
\end{center}
\textbf{Conclusion:} $f$ has a \cle{local maximum} $0$ at $x = 0$ and a \cle{local minimum} $4$ at $x = 2$. Note the maximum value is \emph{smaller} than the minimum value — perfectly possible since they lie on different branches separated by the asymptote $x = 1$.
\end{exemplebox}

\begin{savaistu}[Range of a Function]
The \cle{range} (or image set) $f(D_f)$ is the set of all values $f(x)$ for $x \in D_f$. It can be read directly from the table of variation: it is the union of all intervals covered by $f(x)$ as $x$ runs through $D_f$. For the example above, the range is $(-\infty, 0] \cup [4, +\infty)$. This addresses Lesson 68: "Range of values for the existence or non-existence of a function".
\end{savaistu}

\courssub{Concavity, Points of Inflexion and Symmetry}

\begin{propriete}[Concavity and Second Derivative]
Let $f$ be twice differentiable on an interval $I$.
\begin{itemize}
\item If $f''(x) > 0$ on $I$, $f$ is \cle{convex} (concave up, $\cup$-shaped) on $I$.
\item If $f''(x) < 0$ on $I$, $f$ is \cle{concave} (concave down, $\cap$-shaped) on $I$.
\end{itemize}
\end{propriete}

\begin{definition}[Point of Inflexion]
A point $I(x_0, f(x_0))$ on the curve of $f$ is a \cle{point of inflexion} if the concavity changes at $x_0$ (i.e. $f''$ changes sign). At such a point, the tangent line crosses the curve.
\end{definition}

\begin{propriete}[Centre and Axis of Symmetry]
\begin{itemize}
\item \textbf{Centre of symmetry} $\Omega(a, b)$: the curve is invariant under a half-turn about $\Omega$. Equivalently, for all $h$ such that $a \pm h \in D_f$,
\[ f(a + h) + f(a - h) = 2b. \]
\item \textbf{Axis of symmetry} $x = a$: the curve is invariant under reflection across the vertical line $x = a$. Equivalently, for all admissible $h$,
\[ f(a + h) = f(a - h). \]
For an even function, $a = 0$: $f(-x) = f(x)$.
\end{itemize}
\textbf{Note:} For a rational function $f(x) = \frac{ax+b}{cx+d}$ ($c \neq 0$), the centre of symmetry is the intersection of its asymptotes: $\Omega\left(-\frac{d}{c},\ \frac{a}{c}\right)$.
\end{propriete}

\begin{methode}[Concavity, Inflexion, Centre and Axis of Symmetry]
\begin{enumerate}
\item \textbf{Concavity:} compute $f''$. Where $f'' > 0$, $f$ is \cle{convex} ($\cup$); where $f'' < 0$, $f$ is \cle{concave} ($\cap$).
\item \textbf{Point of inflexion:} a point where $f''$ \textbf{changes sign} (in particular $f'' = 0$ or $f''$ undefined there, \emph{and} the sign changes). The tangent crosses the curve at such a point.
\item \textbf{Centre of symmetry} $\Omega(a, b)$: show that for all $x$ with $a \pm h \in D_f$,
\[ f(a + h) + f(a - h) = 2b. \]
\item \textbf{Axis of symmetry} $x = a$: show that $f(a + h) = f(a - h)$ for all admissible $h$.
\item \textbf{Exploit symmetry:} study the function on half the domain, then complete the curve by symmetry.
\end{enumerate}
\end{methode}

\begin{exemplebox}[Inflexion point and centre of symmetry]
Let $f(x) = \dfrac{2x+1}{x-1}$ and $(\mathcal{C})$ its curve.
\begin{enumerate}
\item Show that the point $\Omega(1, 2)$ is a centre of symmetry of $(\mathcal{C})$.
\item Let $g(x) = x^3 - 3x^2 + 2x + 1$. Determine the concavity of $g$ and its point of inflexion.
\end{enumerate}
\tcblower
\textbf{1. Centre of symmetry.} $D_f = \mathbb{R}\setminus\{1\}$, which is symmetric about $1$. For $h \neq 0$:
\begin{align*}
f(1+h) + f(1-h) &= \frac{2(1+h)+1}{(1+h)-1} + \frac{2(1-h)+1}{(1-h)-1} \\
&= \frac{2h+3}{h} + \frac{3-2h}{-h} = \frac{2h+3}{h} + \frac{2h-3}{h} = \frac{4h}{h} = 4 = 2 \times 2.
\end{align*}
Thus $f(1+h) + f(1-h) = 2 \times 2$ for all $h \neq 0$: $\Omega(1, 2)$ is a \cle{centre of symmetry} of $(\mathcal{C})$. Note $\Omega$ is the intersection of the asymptotes $x = 1$ and $y = 2$ — a general fact for hyperbolas $y = \dfrac{ax+b}{cx+d}$.

\textbf{2. Concavity of $g$.} Compute:
\[ g'(x) = 3x^2 - 6x + 2, \qquad g''(x) = 6x - 6 = 6(x-1). \]
\begin{itemize}
\item $g''(x) < 0$ for $x < 1$: $g$ is \textbf{concave} ($\cap$) on $(-\infty, 1)$;
\item $g''(x) > 0$ for $x > 1$: $g$ is \textbf{convex} ($\cup$) on $(1, +\infty)$.
\end{itemize}
Since $g''$ vanishes \emph{and changes sign} at $x = 1$, the point $I(1, g(1))$ is a \cle{point of inflexion}, with $g(1) = 1 - 3 + 2 + 1 = 1$: $I(1, 1)$. (For a cubic, the inflexion point is always a centre of symmetry of the curve.)
\end{exemplebox}

\begin{attention}[Inflexion Point Trap]
$f''(x_0) = 0$ is \textbf{not sufficient} for an inflexion point. You must check that $f''$ \textbf{changes sign} at $x_0$. Example: $f(x) = x^4$, $f''(0) = 0$ but $f''(x) = 12x^2 \geq 0$ everywhere — no inflexion at $0$ (it is a minimum).
\end{attention}

\courssub{Algebraic Solutions of Absolute-Value Inequalities}

\begin{propriete}[Fundamental Equivalences for $|u|$]
For $a > 0$:
\[ |u| < a \iff -a < u < a; \qquad |u| > a \iff u < -a \ \text{or}\ u > a. \]
For $a \geq 0$:
\[ |u| \leq a \iff -a \leq u \leq a; \qquad |u| \geq a \iff u \leq -a \ \text{or}\ u \geq a. \]
\end{propriete}

\begin{propriete}[Two Absolute Values]
For any expressions $u, v$:
\[ |u| \leq |v| \iff u^2 \leq v^2 \iff (u-v)(u+v) \leq 0. \]
(Squaring is valid because both sides are non-negative.)
\end{propriete}

\begin{methode}[Solving Absolute-Value Inequalities Algebraically]
\begin{enumerate}
\item \textbf{Simple case $|u| \lessgtr a$:} use the fundamental equivalences directly.
\item \textbf{Two absolute values $|u| \lessgtr |v|$:} square both sides: $|u| \lessgtr |v| \iff u^2 \lessgtr v^2 \iff (u-v)(u+v) \lessgtr 0$.
\item \textbf{Several absolute values with different signs (e.g. $|u| + |v| < k$):} use the \textbf{critical-point method}:
\begin{enumerate}
\item Find where each expression inside $|.|$ vanishes (critical points).
\item Split $\mathbb{R}$ into intervals delimited by these points.
\item On each interval, remove the absolute value bars with the correct sign (replace $|E|$ by $E$ or $-E$).
\item Solve the resulting inequality \textbf{on that interval}.
\item \textbf{Intersect} the solution with the interval.
\item Take the \textbf{union} of all interval solutions.
\end{enumerate}
\item \textbf{Always} give the solution set in interval notation and check a sample value per interval.
\end{enumerate}
\end{methode}

\begin{attention}[Trap: Squaring Inequalities]
Never square an inequality $|A| < B$ unless you know $B \geq 0$; squaring $|A| < B$ when $B$ could be negative is invalid. Also, never forget to \textbf{intersect} each case's solution with the interval defining that case.
\end{attention}

\begin{exemplebox}[GCE-style absolute-value inequalities]
Solve in $\mathbb{R}$:
\begin{enumerate}
\item $|2x - 3| \leq 5$;
\item $|x - 1| > |2x + 1|$;
\item $|x - 2| + |x + 1| < 5$.
\end{enumerate}
\tcblower
\textbf{1.} $|2x-3| \leq 5 \iff -5 \leq 2x - 3 \leq 5 \iff -2 \leq 2x \leq 8 \iff -1 \leq x \leq 4$.
\[ S = [-1, 4]. \]

\textbf{2.} Both sides are non-negative; square:
\[ |x-1| > |2x+1| \iff (x-1)^2 > (2x+1)^2 \iff (x-1)^2 - (2x+1)^2 > 0. \]
Factorise as a difference of squares:
\[ \big[(x-1) - (2x+1)\big]\big[(x-1) + (2x+1)\big] = (-x-2)(3x) > 0 \iff -3x(x+2) > 0 \iff x(x+2) < 0. \]
The roots are $-2$ and $0$; the product is negative between the roots:
\[ S = (-2, 0). \]
Check $x = -1$: $|-2| = 2 > |-1| = 1$. \checkmark

\textbf{3.} Critical points: $x - 2 = 0$ at $x = 2$; $x + 1 = 0$ at $x = -1$. Three cases.

\emph{Case $x < -1$:} $|x-2| = 2-x$, $|x+1| = -x-1$. Inequality: $(2-x) + (-x-1) < 5 \iff 1 - 2x < 5 \iff x > -2$. Intersect with $x < -1$: $(-2, -1)$.

\emph{Case $-1 \leq x \leq 2$:} $|x-2| = 2-x$, $|x+1| = x+1$. Inequality: $(2-x)+(x+1) = 3 < 5$: always true. Solution on this case: $[-1, 2]$.

\emph{Case $x > 2$:} $|x-2| = x-2$, $|x+1| = x+1$. Inequality: $2x - 1 < 5 \iff x < 3$. Intersect with $x > 2$: $(2, 3)$.

Union of the three cases:
\[ S = (-2, -1) \cup [-1, 2] \cup (2, 3) = (-2, 3). \]
Check $x = 0$: $|{-2}| + |1| = 3 < 5$ \checkmark; check $x = 4$: $2 + 5 = 7 \not< 5$ \checkmark.
\end{exemplebox}

\courssub{Graphical Solutions of Inequalities}

\begin{methode}[Solving Inequalities Graphically]
To solve $f(x) < g(x)$, $f(x) \leq g(x)$, $|f(x)| < k$, etc., graphically:
\begin{enumerate}
\item \textbf{Sketch both sides} as separate curves on the same axes: $y = f(x)$ and $y = g(x)$ (for $|f(x)|$: reflect the part of the curve below the $x$-axis upwards).
\item \textbf{Find the intersection points} (solve $f(x) = g(x)$ algebraically if needed) — their abscissas are the boundary values of the solution.
\item \textbf{Read off:} $f(x) < g(x)$ holds exactly where the curve of $f$ lies \textbf{strictly below} the curve of $g$.
\item \textbf{Write the solution} as interval(s) of $x$-values; use open bounds for strict inequalities, closed bounds for non-strict ones (provided the point is in the domain).
\end{enumerate}
\textbf{Reflex:} a graphical solution must be supported by a clear, labelled sketch — in the GCE, marks are awarded for the diagram.
\end{methode}

\begin{exemplebox}[Graphical solution of $|x^2 - 2x| < 3 - x$]
The curve of $f(x) = |x^2 - 2x|$ and the line $y = 3 - x$ are given.
\begin{enumerate}
\item Determine algebraically the abscissas of the intersection points.
\item Hence solve graphically $|x^2 - 2x| < 3 - x$.
\end{enumerate}
\tcblower
\textbf{1. Intersections.} Solve $|x^2 - 2x| = 3 - x$ (requires $3 - x \geq 0$, i.e. $x \leq 3$).

\emph{Case A:} $x^2 - 2x \geq 0$ (i.e. $x \leq 0$ or $x \geq 2$). Then
\[ x^2 - 2x = 3 - x \iff x^2 - x - 3 = 0 \iff x = \frac{1 \pm \sqrt{13}}{2}. \]
$\dfrac{1+\sqrt{13}}{2} \approx 2.30$ (valid: $\geq 2$ and $\leq 3$ \checkmark); $\dfrac{1-\sqrt{13}}{2} \approx -1.30$ (valid: $\leq 0$ \checkmark).

\emph{Case B:} $x^2 - 2x < 0$ (i.e. $0 < x < 2$). Then
\[ -(x^2 - 2x) = 3 - x \iff x^2 - 3x + 3 = 0, \quad \Delta = 9 - 12 < 0: \text{ no solution.} \]

Intersection abscissas: $x_1 = \dfrac{1-\sqrt{13}}{2} \approx -1.30$ and $x_2 = \dfrac{1+\sqrt{13}}{2} \approx 2.30$.

\textbf{2. Graphical reading.} Sketch: $y = x^2 - 2x$ is a parabola through $(0,0)$ and $(2,0)$ with vertex $(1,-1)$; taking the absolute value reflects the dip between $0$ and $2$ into a ``bump'' peaking at $(1, 1)$. The line $y = 3 - x$ falls through $(0,3)$ and $(3,0)$.

\begin{popfigure}
\begin{center}
\begin{tikzpicture}[scale=0.85]
\draw[->, popline] (-2.6,0) -- (4,0) node[right]{$x$};
\draw[->, popline] (0,-1.4) -- (0,4.4) node[above]{$y$};
\draw[domain=-1.9:3.6, smooth, thick, popblue] plot (\x, {abs(\x*\x - 2*\x)});
\draw[domain=-2.2:3.6, smooth, thick, poporange] plot (\x, {3 - \x});
\fill[popdark] (-1.303,0) circle (1.5pt) node[below, font=\small]{$x_1$};
\fill[popdark] (2.303,0) circle (1.5pt) node[below, font=\small]{$x_2$};
\draw[dashed, popmuted] (-1.303,0) -- (-1.303,4.303);
\draw[dashed, popmuted] (2.303,0) -- (2.303,0.697);
\node[popblue, font=\small] at (3.4,3.4) {$y=|x^2-2x|$};
\node[poporange, font=\small] at (-1.9,4.2) {$y=3-x$};
\end{tikzpicture}
\end{center}
\legende{The curve $y = |x^2-2x|$ lies strictly below the line between the two intersection points.}
\end{popfigure}

The curve of $f$ is strictly below the line for $x$ between the two intersection abscissas:
\[ S = \left( \frac{1-\sqrt{13}}{2},\ \frac{1+\sqrt{13}}{2} \right). \]
The bounds are open because the inequality is strict (at $x_1, x_2$ the two graphs meet).
\end{exemplebox}

\courssub{Complete Curve Sketching (Synthesis)}

\begin{methode}[The Seven-Step Plan for Sketching a Curve]
\begin{enumerate}
\item \textbf{Domain} $D_f$ (and symmetry: even, odd, centre/axis of symmetry — to reduce the study).
\item \textbf{Limits at the bounds} of $D_f$; deduce vertical/horizontal asymptotes.
\item \textbf{Derivative} $f'$: sign table, stationary points, extrema values.
\item \textbf{Second derivative} $f''$: concavity and points of inflexion (if required).
\item \textbf{Table of variation} summarising everything (limits, extrema, arrows).
\item \textbf{Infinite branches:} oblique asymptote by division, or parabolic branch via $\lim f(x)/x$; position of curve relative to asymptote by the sign of the difference.
\item \textbf{Sketch:} draw asymptotes (dashed), plot intercepts and extrema, then draw each branch respecting the table, concavity and relative positions.
\end{enumerate}
\end{methode}

\begin{exemplebox}[Full study and sketch — GCE Advanced Level style]
Let $f(x) = \dfrac{x^2 + x - 2}{x - 2}$ and $(\mathcal{C})$ its representative curve.
\begin{enumerate}
\item Determine $D_f$ and the limits at its bounds; deduce the asymptotes parallel to the axes.
\item Show that the line $y = x + 3$ is an oblique asymptote and study the position of $(\mathcal{C})$ relative to it.
\item Study the variations of $f$ and draw the table of variation.
\item Sketch $(\mathcal{C})$.
\end{enumerate}
\tcblower
\textbf{1. Domain, limits, asymptotes.} $D_f = \mathbb{R} \setminus \{2\}$.

At $x = 2$: numerator $= 4 + 2 - 2 = 4 > 0$; denominator $\to 0^{\pm}$:
\[ \lim_{x\to 2^-} f(x) = -\infty, \qquad \lim_{x\to 2^+} f(x) = +\infty \quad\Rightarrow\quad x = 2 \ \text{vertical asymptote.} \]

At $\pm\infty$: $\deg(\text{num}) = \deg(\text{den}) + 1$, so $\lim_{x\to\pm\infty} f(x) = \pm\infty$ (leading terms $\tfrac{x^2}{x} = x$): no horizontal asymptote; there is an infinite branch to analyse.

\textbf{2. Oblique asymptote and position.} Divide:
\[ x^2 + x - 2 = (x-2)(x+3) + 4 \quad\big(\text{check: } (x-2)(x+3) = x^2 + x - 6;\ +4 \text{ gives } x^2+x-2\ \checkmark\big) \]
\[ f(x) = x + 3 + \frac{4}{x-2}, \qquad \lim_{x\to\pm\infty} \frac{4}{x-2} = 0, \]
so $y = x + 3$ is an \cle{oblique asymptote} at both ends.

Position: $f(x) - (x+3) = \dfrac{4}{x-2}$. For $x > 2$: positive, $(\mathcal{C})$ \textbf{above} the line; for $x < 2$: negative, $(\mathcal{C})$ \textbf{below} the line.

\textbf{3. Variations.} Quotient rule:
\[ f'(x) = \frac{(2x+1)(x-2) - (x^2+x-2)}{(x-2)^2} = \frac{2x^2 - 3x - 2 - x^2 - x + 2}{(x-2)^2} = \frac{x^2 - 4x}{(x-2)^2} = \frac{x(x-4)}{(x-2)^2}. \]
$f'(x) = 0$ at $x = 0$ and $x = 4$; sign of $f'$ = sign of $x(x-4)$: $+$ on $(-\infty, 0)$, $-$ on $(0, 2)$ and $(2, 4)$, $+$ on $(4, +\infty)$.

Extrema: $f(0) = \dfrac{-2}{-2} = 1$ (local max); $f(4) = \dfrac{16+4-2}{2} = 9$ (local min).

Intercepts: $f(0) = 1$; $f(x) = 0 \iff x^2 + x - 2 = 0 \iff (x+2)(x-1) = 0 \iff x = -2$ or $x = 1$.

\begin{center}
\begin{tabular}{|c|ccccccccc|}
\hline
$x$ & $-\infty$ & & $0$ & & $2$ & & $4$ & & $+\infty$ \\
\hline
$f'(x)$ & & $+$ & $0$ & $-$ & $\|$ & $-$ & $0$ & $+$ & \\
\hline
$f(x)$ & $-\infty$ & $\nearrow$ & $1$ & $\searrow$ & $-\infty\,\|\,+\infty$ & $\searrow$ & $9$ & $\nearrow$ & $+\infty$ \\
\hline
\end{tabular}
\end{center}

\textbf{4. Sketch.} Draw the dashed asymptotes $x = 2$ and $y = x + 3$; plot $(-2, 0)$, $(1, 0)$, $(0, 1)$ (maximum), $(4, 9)$ (minimum).

\begin{popfigure}
\begin{center}
\begin{tikzpicture}[scale=0.5]
\draw[->, popline] (-6.5,0) -- (7.5,0) node[right]{$x$};
\draw[->, popline] (0,-7.5) -- (0,11.5) node[above]{$y$};
\draw[dashed, popmuted] (2,-7.5) -- (2,11.5) node[below left, font=\small]{$x=2$};
\draw[dashed, popmuted, domain=-6:7] plot (\x, {\x+3}) node[pos=0.06, above, sloped, font=\small]{$y=x+3$};
\draw[domain=-6:1.55, smooth, thick, popblue, samples=80] plot (\x, {(\x*\x+\x-2)/(\x-2)});
\draw[domain=2.42:7, smooth, thick, popblue, samples=80] plot (\x, {(\x*\x+\x-2)/(\x-2)});
\fill[popdark] (0,1) circle (2.5pt) node[above right, font=\small]{$(0,1)$ max};
\fill[popdark] (4,9) circle (2.5pt) node[above right, font=\small]{$(4,9)$ min};
\fill[poporange] (-2,0) circle (2.5pt) node[below, font=\small]{$-2$};
\fill[poporange] (1,0) circle (2.5pt) node[below, font=\small]{$1$};
\end{tikzpicture}
\end{center}
\legende{The curve $(\mathcal{C})$: two branches, asymptotes $x=2$ and $y=x+3$, maximum at $(0,1)$, minimum at $(4,9)$.}
\end{popfigure}

The left branch rises from the asymptote $y = x+3$ (below it) to the maximum $(0,1)$, then falls through $(1,0)$ down to $-\infty$ along $x = 2$. The right branch descends from $+\infty$ along $x = 2$ to the minimum $(4, 9)$, then rises towards the asymptote $y = x+3$, staying above it — exactly as the table of variation and the position study predict.
\end{exemplebox}

\begin{exercice}[Practice Exercises]
\begin{enumerate}
\item Let $f(x) = \dfrac{x^2 - 4}{x - 1}$. Find $D_f$, all asymptotes, intercepts, stationary points, table of variation, and sketch the curve.
\item Solve algebraically: $|3x - 2| \geq |x + 4|$.
\item Solve graphically: $|x^2 - 4| < x + 2$. (Sketch the curves and give the solution set.)
\item Let $f(x) = \dfrac{x}{x^2 + 1}$. Show that the origin is a centre of symmetry. Find the range of $f$.
\item A curve has equation $y = \dfrac{x^2 + 2x + 3}{x + 1}$. Find its asymptotes, stationary points, and sketch it.
\end{enumerate}
\end{exercice}

\begin{corrige}[Exercise 1]
$f(x) = \frac{x^2-4}{x-1}$, $D_f = \mathbb{R}\setminus\{1\}$.
Vertical asymptote $x=1$ ($\lim_{x\to1^-}=-\infty$, $\lim_{x\to1^+}=+\infty$).
Division: $x^2-4 = (x-1)(x+1) - 3$, so $f(x) = x+1 - \frac{3}{x-1}$. Oblique asymptote $y=x+1$.
Position: $f(x)-(x+1) = -\frac{3}{x-1}$. Above for $x<1$, below for $x>1$.
$f'(x) = \frac{x^2-2x+4}{(x-1)^2} > 0$ on $D_f$ (numerator $\Delta = -12 < 0$). No stationary points; $f$ strictly increasing on each branch.
Intercepts: $y$-int $f(0)=4$; $x$-ints $x=\pm2$.
Table: $(-\infty,1)$: $-\infty \nearrow +\infty$; $(1,+\infty)$: $-\infty \nearrow +\infty$.
Sketch: two increasing branches, asymptotes $x=1$, $y=x+1$, passing through $(-2,0)$, $(0,4)$, $(2,0)$.
\end{corrige}

\begin{corrige}[Exercise 2]
$|3x-2| \geq |x+4| \iff (3x-2)^2 \geq (x+4)^2 \iff 9x^2-12x+4 \geq x^2+8x+16 \iff 8x^2-20x-12 \geq 0 \iff 2x^2-5x-3 \geq 0 \iff (2x+1)(x-3) \geq 0$.
Roots: $-\frac12$, $3$. Solution: $(-\infty, -\frac12] \cup [3, +\infty)$.
\end{corrige}

\begin{corrige}[Exercise 3]
$|x^2-4| < x+2$ requires $x+2 > 0 \implies x > -2$.
Case $x^2-4 \geq 0$ ($x \leq -2$ or $x \geq 2$): on $x > -2$, this means $x \geq 2$. Then $x^2-4 < x+2 \iff x^2-x-6 < 0 \iff (x-3)(x+2) < 0 \iff -2 < x < 3$. Intersect with $x \geq 2$: $[2, 3)$.
Case $x^2-4 < 0$ ($-2 < x < 2$): $-(x^2-4) < x+2 \iff -x^2+4 < x+2 \iff x^2+x-2 > 0 \iff (x+2)(x-1) > 0 \iff x < -2$ or $x > 1$. Intersect with $-2 < x < 2$: $(1, 2)$.
Union: $(1, 2) \cup [2, 3) = (1, 3)$.
Graphically: sketch $y=|x^2-4|$ (W-shaped with vertices at $(\pm2,0)$ and $(0,4)$) and line $y=x+2$ (through $(-2,0)$ and $(0,2)$). Intersections at $x=1$ and $x=3$. Curve below line on $(1,3)$.
\end{corrige}

\begin{corrige}[Exercise 4]
$f(x) = \frac{x}{x^2+1}$, $D_f = \mathbb{R}$.
$f(-x) = \frac{-x}{x^2+1} = -f(x)$: odd function $\implies$ centre of symmetry at origin $O(0,0)$.
$f'(x) = \frac{1\cdot(x^2+1) - x(2x)}{(x^2+1)^2} = \frac{1-x^2}{(x^2+1)^2}$.
$f'(x)=0 \iff x=\pm1$. Sign: $+$ on $(-1,1)$, $-$ on $(-\infty,-1)$ and $(1,+\infty)$.
Max at $x=1$: $f(1)=\frac12$; min at $x=-1$: $f(-1)=-\frac12$.
Limits at $\pm\infty$: $0$ (horizontal asymptote $y=0$).
Range: $[-\frac12, \frac12]$.
\end{corrige}

\begin{corrige}[Exercise 5]
$f(x) = \frac{x^2+2x+3}{x+1}$, $D_f = \mathbb{R}\setminus\{-1\}$.
Vertical asymptote $x=-1$ (num $=2>0$: $\lim_{x\to-1^-}=-\infty$, $\lim_{x\to-1^+}=+\infty$).
Division: $x^2+2x+3 = (x+1)(x+1) + 2$, so $f(x) = x+1 + \frac{2}{x+1}$. Oblique asymptote $y=x+1$.
Position: $f(x)-(x+1) = \frac{2}{x+1}$. Above for $x>-1$, below for $x<-1$.
$f'(x) = \frac{(2x+2)(x+1) - (x^2+2x+3)}{(x+1)^2} = \frac{2x^2+4x+2 - x^2-2x-3}{(x+1)^2} = \frac{x^2+2x-1}{(x+1)^2}$.
$f'(x)=0 \iff x^2+2x-1=0 \iff x = -1 \pm \sqrt{2}$.
$x_1 = -1-\sqrt{2} \approx -2.41$ (left branch), $x_2 = -1+\sqrt{2} \approx 0.41$ (right branch).
$f(x_1) = -1-\sqrt{2}+1 + \frac{2}{-\sqrt{2}} = -\sqrt{2} - \sqrt{2} = -2\sqrt{2}$ (local max on left).
$f(x_2) = -1+\sqrt{2}+1 + \frac{2}{\sqrt{2}} = \sqrt{2} + \sqrt{2} = 2\sqrt{2}$ (local min on right).
Intercepts: $y$-int $f(0)=3$; no $x$-int (num $>0$).
Sketch: two branches, asymptotes $x=-1$, $y=x+1$, max at $(-1-\sqrt{2}, -2\sqrt{2})$, min at $(-1+\sqrt{2}, 2\sqrt{2})$, through $(0,3)$.
\end{corrige}

\newpage

\coursec{Exercises}

\courssub{I check my knowledge}

\begin{exercice}[MCQ: domain and asymptotes]
For each question, choose the single correct answer, justifying briefly.
\begin{enumerate}
\item The domain of definition of $f(x)=\dfrac{x+1}{x^{2}-4}$ is:
\begin{enumerate}
\item $\mathbb{R}$ \quad (b) $\mathbb{R}\setminus\{-2,2\}$ \quad (c) $\mathbb{R}\setminus\{-1,2\}$ \quad (d) $\mathbb{R}\setminus\{4\}$
\end{enumerate}
\item The curve of $g(x)=\dfrac{2x+3}{x-1}$ admits:
\begin{enumerate}
\item a vertical asymptote $x=1$ and a horizontal asymptote $y=2$
\item a vertical asymptote $x=-1$ and a horizontal asymptote $y=3$
\item only a vertical asymptote $x=1$
\item only a horizontal asymptote $y=2$
\end{enumerate}
\item If $f''(x)<0$ on an interval $I$, then on $I$ the curve of $f$ is:
\begin{enumerate}
\item concave up \quad (b) concave down \quad (c) increasing \quad (d) a straight line
\end{enumerate}
\item The solution set of $|x-3|<2$ is:
\begin{enumerate}
\item $]1,5[$ \quad (b) $[1,5]$ \quad (c) $]-\infty,1[\,\cup\,]5,+\infty[$ \quad (d) $\{1,5\}$
\end{enumerate}
\end{enumerate}
\end{exercice}

\begin{exercice}[True or false]
State whether each assertion is true or false, and justify your answer with a proof or a counter-example.
\begin{enumerate}
\item Every rational function has at least one vertical asymptote.
\item If $f'(c)=0$, then the point of the curve of abscissa $c$ is necessarily a local extremum.
\item If $\lim_{x\to+\infty}\big(f(x)-(2x+1)\big)=0$, then the line $y=2x+1$ is an asymptote to the curve of $f$ at $+\infty$.
\item The inequality $|2x+1|>5$ is equivalent to $2x+1>5$.
\item If $I(a,b)$ is a centre of symmetry of the curve of $f$, then $f(a+x)+f(a-x)=2b$ for all suitable $x$.
\end{enumerate}
\end{exercice}

\begin{exercice}[Definitions and vocabulary]
\begin{enumerate}
\item Define, in your own words, the \cle{domain of definition} of a real function.
\item What is a \cle{point of inflexion} of a curve? State the condition on $f''$ that allows one to detect it.
\item State the condition for the line $x=a$ to be a vertical asymptote of the curve of $f$.
\item What is an \cle{oblique asymptote}? Give the two limits that must be computed to find and confirm it.
\end{enumerate}
\end{exercice}

\begin{exercice}[Direct calculations]
\begin{enumerate}
\item Find the domain of definition of each function:
\[ f(x)=\sqrt{3x-6}, \qquad g(x)=\frac{x^{2}+1}{x^{2}-x-6}, \qquad h(x)=\sqrt{\frac{x-1}{x+2}}. \]
\item Compute the following limits:
\[ \lim_{x\to 3^{+}}\frac{x+1}{x-3}, \qquad \lim_{x\to+\infty}\frac{2x^{2}-x}{x^{2}+4}, \qquad \lim_{x\to+\infty}\left(\frac{x^{2}+1}{x}-x\right). \]
\item Solve in $\mathbb{R}$: $|5x-1|=9$.
\end{enumerate}
\end{exercice}

\courssub{I practise}

\begin{exercice}[Existence of a square-root function]
Find the values of the real parameter $m$ for which the function
\[ f(x)=\sqrt{mx^{2}+4x+1} \]
is defined for \emph{all} real numbers $x$. (Consider separately the cases $m=0$ and $m\neq 0$.)
\end{exercice}

\begin{exercice}[Absolute-value inequalities, algebraically]
Solve the following inequalities in $\mathbb{R}$, giving your answers in interval notation:
\begin{enumerate}
\item $|3x-2|\leq 7$;
\item $|2x+1|>5$;
\item $|x-4|\geq|x+2|$ (use squaring);
\item $1<|2x-3|\leq 5$.
\end{enumerate}
\end{exercice}

\begin{exercice}[Asymptotes and position of a curve]
Let $f(x)=\dfrac{x^{2}-x+3}{x+1}$ and let $(\mathcal{C})$ be its curve.
\begin{enumerate}
\item Determine the domain of definition of $f$ and the limits at its bounds.
\item Show that the line $y=x-2$ is an oblique asymptote of $(\mathcal{C})$.
\item Study the sign of $f(x)-(x-2)$ and deduce the position of $(\mathcal{C})$ relative to its oblique asymptote.
\item Find the intercepts of $(\mathcal{C})$ with the coordinate axes.
\end{enumerate}
\end{exercice}

\begin{exercice}[Concavity and points of inflexion]
Let $f(x)=x^{4}-6x^{2}+5$.
\begin{enumerate}
\item Compute $f''(x)$ and study its sign on $\mathbb{R}$.
\item Deduce the intervals on which the curve of $f$ is concave up or concave down, and determine the coordinates of the points of inflexion.
\item Show that the $y$-axis is an axis of symmetry of the curve.
\end{enumerate}
\end{exercice}

\courssub{I prepare for the GCE}

\begin{exercice}[Complete study of a rational function (GCE style)]
Let $f(x)=\dfrac{x^{2}-3x+2}{x+2}$ and let $(\mathcal{C})$ be its curve in an orthonormal frame.
\begin{enumerate}
\item Determine the domain of definition of $f$ and compute the limits of $f$ at the bounds of this domain. \hfill(3 marks)
\item Deduce the equations of all asymptotes of $(\mathcal{C})$, showing that $(\mathcal{C})$ has an oblique asymptote whose equation is to be found. \hfill(4 marks)
\item Compute $f'(x)$, show that $f'(x)=\dfrac{x^{2}+4x-8}{(x+2)^{2}}$, and construct the complete table of variation of $f$ (exact values, simplified with radicals). \hfill(5 marks)
\item Study the position of $(\mathcal{C})$ relative to its oblique asymptote, and find the intercepts of $(\mathcal{C})$ with the axes. \hfill(3 marks)
\item Sketch $(\mathcal{C})$ carefully, showing the asymptotes and the stationary points. \hfill(3 marks)
\item Using your graph, deduce the number of real solutions of the equation $f(x)=k$ according to the values of the real parameter $k$. \hfill(2 marks)
\end{enumerate}
\end{exercice}

\begin{exercice}[Centre of symmetry and curve sketching (GCE style)]
Let $f(x)=\dfrac{2x^{2}-3x+4}{x-1}$ and let $(\mathcal{C})$ be its curve.
\begin{enumerate}
\item Determine the domain of definition of $f$ and the limits at its bounds; deduce the vertical asymptote. \hfill(3 marks)
\item Show that $f(x)=2x-1+\dfrac{3}{x-1}$ and deduce an oblique asymptote of $(\mathcal{C})$. \hfill(3 marks)
\item Prove that the point $I(1,2)$ is a centre of symmetry of $(\mathcal{C})$ by showing that $f(1+x)+f(1-x)=4$ for all $x\neq 0$. \hfill(3 marks)
\item Compute $f'(x)$ and construct the table of variation of $f$. \hfill(4 marks)
\item Sketch $(\mathcal{C})$, using the centre of symmetry to complete the drawing economically. \hfill(3 marks)
\item Solve graphically, explaining your method: (i) $f(x)=2x-1$; (ii) $f(x)\leq 2x-1$. \hfill(4 marks)
\end{enumerate}
\end{exercice}

\begin{exercice}[Inequalities from a table of variation and absolute values (GCE Paper 2 style)]
The table of variation of a function $g$ defined on $\mathbb{R}$ is given below:
\begin{center}
\begin{tabular}{|c|ccccccc|}
\hline
$x$ & $-\infty$ & & $-1$ & & $2$ & & $+\infty$ \\
\hline
$g'(x)$ & & $+$ & $0$ & $-$ & $0$ & $+$ & \\
\hline
$g(x)$ & $-2$ & $\nearrow$ & $3$ & $\searrow$ & $-1$ & $\nearrow$ & $+\infty$ \\
\hline
\end{tabular}
\end{center}
It is also given that $g(x)=0$ has exactly three solutions $-2$, $0$ and $3$.
\begin{enumerate}
\item Sketch a possible curve $(\mathcal{C})$ of $g$, consistent with all the information given. \hfill(4 marks)
\item Solve graphically, justifying each answer: (i) $g(x)=0$; (ii) $g(x)>1$ (express the answer using the values read from your sketch, denoting unknown abscissae by letters); (iii) $|g(x)|<2$. \hfill(6 marks)
\item On the same frame, sketch the curve of $y=|g(x)|$, explaining the construction. \hfill(3 marks)
\item Solve algebraically the inequality $|x-2|<|2x+1|$ by two different methods: (a) by squaring both sides; (b) by the critical-value method. Compare the two approaches and state the solution set in interval notation. \hfill(7 marks)
\end{enumerate}
\end{exercice}

\newpage

\coursec{Solutions to the exercises}

\begin{corrige}[Exercise 1 — MCQ: domain and asymptotes]
\begin{enumerate}
\item \textbf{Answer (b).} The denominator vanishes when $x^{2}-4=0$, i.e. $x=\pm 2$. Hence $D_f=\mathbb{R}\setminus\{-2,2\}$.
\item \textbf{Answer (a).} The denominator vanishes at $x=1$ and the numerator does not ($2(1)+3=5\neq 0$), so $x=1$ is a vertical asymptote. Since $\lim_{x\to\pm\infty}\dfrac{2x+3}{x-1}=\lim_{x\to\pm\infty}\dfrac{2x}{x}=2$, the line $y=2$ is a horizontal asymptote.
\item \textbf{Answer (b).} $f''(x)<0$ on $I$ means $f'$ is decreasing on $I$, which characterises a \cle{concave down} curve.
\item \textbf{Answer (a).} $|x-3|<2 \iff -2<x-3<2 \iff 1<x<5$, so $\mathcal{S}=\,]1,5[$ (strict inequality, open interval).
\end{enumerate}
\end{corrige}

\begin{corrige}[Exercise 2 — True or false]
\begin{enumerate}
\item \textbf{False.} Counter-example: $f(x)=\dfrac{1}{x^{2}+1}$ is a rational function defined on all of $\mathbb{R}$; its denominator never vanishes, so it has no vertical asymptote (it has the horizontal asymptote $y=0$).
\item \textbf{False.} Counter-example: $f(x)=x^{3}$ at $c=0$: $f'(0)=0$, but $f$ is strictly increasing on $\mathbb{R}$, so the point $(0,0)$ is not a local extremum (it is a stationary point of inflexion). The condition $f'(c)=0$ is necessary but not sufficient: the derivative must also \emph{change sign} at $c$.
\item \textbf{True.} This is exactly the definition of an oblique (slant) asymptote: if $\lim_{x\to+\infty}\big(f(x)-(ax+b)\big)=0$, the distance between the curve and the line $y=ax+b$ tends to $0$, so the line is an asymptote at $+\infty$.
\item \textbf{False.} $|2x+1|>5 \iff 2x+1>5$ \emph{or} $2x+1<-5$. Counter-example: $x=-4$ satisfies $|2(-4)+1|=7>5$ but $2(-4)+1=-7<5$.
\item \textbf{True.} $I(a,b)$ is a centre of symmetry iff for every $x$ such that $a\pm x\in D_f$, the points $M(a+x,f(a+x))$ and $M'(a-x,f(a-x))$ are symmetric about $I$, i.e. $I$ is the midpoint of $[MM']$: $\dfrac{f(a+x)+f(a-x)}{2}=b$, equivalently $f(a+x)+f(a-x)=2b$.
\end{enumerate}
\end{corrige}

\begin{corrige}[Exercise 3 — Definitions and vocabulary]
\begin{enumerate}
\item The \cle{domain of definition} of a real function $f$ is the set of all real numbers $x$ for which $f(x)$ exists (can be computed). One excludes values making a denominator zero, values making the radicand of an even root negative, and values making the argument of a logarithm non-positive.
\item A \cle{point of inflexion} is a point where the curve crosses its tangent, i.e. where the concavity changes. Detection: if $f''$ vanishes at $c$ \emph{and changes sign} at $c$, then the point $(c,f(c))$ is a point of inflexion.
\item The line $x=a$ is a vertical asymptote of the curve of $f$ if at least one of the one-sided limits at $a$ is infinite: $\lim_{x\to a^{+}}f(x)=\pm\infty$ or $\lim_{x\to a^{-}}f(x)=\pm\infty$.
\item An \cle{oblique asymptote} is a non-horizontal line $y=ax+b$ ($a\neq 0$) that the curve approaches at infinity. One computes $a=\lim_{x\to\pm\infty}\dfrac{f(x)}{x}$ (must be finite and non-zero), then $b=\lim_{x\to\pm\infty}\big(f(x)-ax\big)$ (must be finite). Equivalently, one checks directly that $\lim_{x\to\pm\infty}\big(f(x)-(ax+b)\big)=0$.
\end{enumerate}
\end{corrige}

\begin{corrige}[Exercise 4 — Direct calculations]
\begin{enumerate}
\item \textbf{Domain of $f$:} $3x-6\geq 0 \iff x\geq 2$, so $D_f=[2,+\infty[$.

\textbf{Domain of $g$:} $x^{2}-x-6=0 \iff (x-3)(x+2)=0 \iff x=3$ or $x=-2$, so $D_g=\mathbb{R}\setminus\{-2,3\}$.

\textbf{Domain of $h$:} we need $\dfrac{x-1}{x+2}\geq 0$ and $x\neq -2$. Sign table: the quotient is positive on $]-\infty,-2[\,\cup\,[1,+\infty[$, zero at $x=1$, negative on $]-2,1[$. Hence $D_h=\,]-\infty,-2[\,\cup\,[1,+\infty[$.

\item $\lim_{x\to 3^{+}}\dfrac{x+1}{x-3}$: numerator $\to 4>0$, denominator $\to 0^{+}$, so the limit is $+\infty$.

$\lim_{x\to+\infty}\dfrac{2x^{2}-x}{x^{2}+4}=\lim_{x\to+\infty}\dfrac{2x^{2}}{x^{2}}=2$ (highest-degree terms).

$\lim_{x\to+\infty}\left(\dfrac{x^{2}+1}{x}-x\right)=\lim_{x\to+\infty}\dfrac{x^{2}+1-x^{2}}{x}=\lim_{x\to+\infty}\dfrac{1}{x}=0$.

\item $|5x-1|=9 \iff 5x-1=9$ or $5x-1=-9 \iff x=2$ or $x=-\dfrac{8}{5}$. Hence $\mathcal{S}=\left\{-\dfrac{8}{5},\,2\right\}$.
\end{enumerate}
\end{corrige}

\begin{corrige}[Exercise 5 — Existence of a square-root function]
We need $mx^{2}+4x+1\geq 0$ for \emph{every} real $x$.

\textbf{Case $m=0$:} the radicand becomes $4x+1$, which is negative for $x<-\tfrac14$. So $m=0$ does not work.

\textbf{Case $m\neq 0$:} a quadratic $mx^{2}+4x+1$ is non-negative on $\mathbb{R}$ if and only if its leading coefficient is positive and its discriminant is non-positive:
\[ m>0 \qquad\text{and}\qquad \Delta=4^{2}-4m=16-4m\leq 0. \]
$\Delta\leq 0 \iff m\geq 4$, which already implies $m>0$.

\textbf{Conclusion:} $f$ is defined for all real $x$ if and only if $\boxed{m\in[4,+\infty[}$.

\emph{Remark:} for $m=4$, the radicand is $4x^{2}+4x+1=(2x+1)^{2}\geq 0$, which confirms the boundary case.
\end{corrige}

\begin{corrige}[Exercise 6 — Absolute-value inequalities, algebraically]
\begin{enumerate}
\item $|3x-2|\leq 7 \iff -7\leq 3x-2\leq 7 \iff -5\leq 3x\leq 9 \iff -\dfrac{5}{3}\leq x\leq 3$.
\[ \mathcal{S}=\left[-\frac{5}{3},\,3\right]. \]

\item $|2x+1|>5 \iff 2x+1>5$ or $2x+1<-5 \iff x>2$ or $x<-3$.
\[ \mathcal{S}=\,]-\infty,-3[\,\cup\,]2,+\infty[. \]

\item Both sides are non-negative, so squaring preserves the order:
\[ |x-4|\geq|x+2| \iff (x-4)^{2}\geq(x+2)^{2} \iff x^{2}-8x+16\geq x^{2}+4x+4 \iff 12\geq 12x \iff x\leq 1. \]
\[ \mathcal{S}=\,]-\infty,1]. \]

\item $1<|2x-3|\leq 5$ splits into two conditions:
\begin{itemize}
\item $|2x-3|>1 \iff 2x-3>1$ or $2x-3<-1 \iff x>2$ or $x<1$;
\item $|2x-3|\leq 5 \iff -5\leq 2x-3\leq 5 \iff -1\leq x\leq 4$.
\end{itemize}
Intersecting: $\big(]-\infty,1[\,\cup\,]2,+\infty[\big)\cap[-1,4]=[-1,1[\,\cup\,]2,4]$.
\[ \mathcal{S}=[-1,1[\,\cup\,]2,4]. \]
\end{enumerate}
\end{corrige}

\begin{corrige}[Exercise 7 — Asymptotes and position of a curve]
$f(x)=\dfrac{x^{2}-x+3}{x+1}$.

\begin{enumerate}
\item \textbf{Domain:} $x+1\neq 0$, so $D_f=\mathbb{R}\setminus\{-1\}=\,]-\infty,-1[\,\cup\,]-1,+\infty[$.

\textbf{Limits at the bounds:}
\begin{itemize}
\item At $x=-1$: numerator $\to 1+1+3=5>0$. As $x\to -1^{+}$, $x+1\to 0^{+}$, so $f(x)\to +\infty$. As $x\to -1^{-}$, $x+1\to 0^{-}$, so $f(x)\to -\infty$.
\item At infinity (rational function, ratio of leading terms): $\lim_{x\to\pm\infty}f(x)=\lim_{x\to\pm\infty}\dfrac{x^{2}}{x}=\lim_{x\to\pm\infty}x$, so $f(x)\to +\infty$ at $+\infty$ and $f(x)\to -\infty$ at $-\infty$.
\end{itemize}

\item Polynomial division: $x^{2}-x+3=(x+1)(x-2)+5$, hence
\[ f(x)=x-2+\frac{5}{x+1}. \]
Then $f(x)-(x-2)=\dfrac{5}{x+1}\xrightarrow[x\to\pm\infty]{}0$, so the line $y=x-2$ is an \cle{oblique asymptote} of $(\mathcal{C})$ at both $+\infty$ and $-\infty$. (Also $x=-1$ is a vertical asymptote, from question 1.)

\item $f(x)-(x-2)=\dfrac{5}{x+1}$ has the sign of $x+1$:
\begin{itemize}
\item for $x>-1$, $f(x)-(x-2)>0$: $(\mathcal{C})$ is \textbf{above} its asymptote;
\item for $x<-1$, $f(x)-(x-2)<0$: $(\mathcal{C})$ is \textbf{below} its asymptote.
\end{itemize}

\item \textbf{$y$-intercept:} $f(0)=\dfrac{3}{1}=3$: point $(0,3)$.

\textbf{$x$-intercepts:} $f(x)=0 \iff x^{2}-x+3=0$; discriminant $\Delta=1-12=-11<0$: no real root. The curve does not meet the $x$-axis.
\end{enumerate}
\end{corrige}

\begin{corrige}[Exercise 8 — Concavity and points of inflexion]
$f(x)=x^{4}-6x^{2}+5$.

\begin{enumerate}
\item $f'(x)=4x^{3}-12x$, so $f''(x)=12x^{2}-12=12(x^{2}-1)=12(x-1)(x+1)$.

Sign of $f''$: a quadratic with roots $-1$ and $1$, positive outside the roots:
\begin{center}
\begin{tabular}{|c|ccccccc|}
\hline
$x$ & $-\infty$ & & $-1$ & & $1$ & & $+\infty$ \\
\hline
$f''(x)$ & & $+$ & $0$ & $-$ & $0$ & $+$ & \\
\hline
\end{tabular}
\end{center}
More precisely: $f''(x)>0$ on $]-\infty,-1[\,\cup\,]1,+\infty[$ and $f''(x)<0$ on $]-1,1[$.

\item The curve is \textbf{concave up} on $]-\infty,-1]$ and $[1,+\infty[$, and \textbf{concave down} on $[-1,1]$. Since $f''$ vanishes \emph{and changes sign} at $x=-1$ and at $x=1$, both give points of inflexion:
\[ f(-1)=1-6+5=0, \qquad f(1)=1-6+5=0. \]
Points of inflexion: $A(-1,0)$ and $B(1,0)$.

\item $f$ is even: for all $x\in\mathbb{R}$,
\[ f(-x)=(-x)^{4}-6(-x)^{2}+5=x^{4}-6x^{2}+5=f(x). \]
Therefore the $y$-axis (line $x=0$) is an \cle{axis of symmetry} of the curve.
\end{enumerate}
\end{corrige}

\begin{corrige}[Exercise 9 — Complete study of a rational function (GCE style)]
$f(x)=\dfrac{x^{2}-3x+2}{x+2}$.

\begin{enumerate}
\item \textbf{Domain} (1 mark): $D_f=\mathbb{R}\setminus\{-2\}$.

\textbf{Limits} (2 marks): at $x=-2$, numerator $\to 4+6+2=12>0$:
\[ \lim_{x\to -2^{+}}f(x)=+\infty, \qquad \lim_{x\to -2^{-}}f(x)=-\infty. \]
At infinity: $\lim_{x\to\pm\infty}f(x)=\lim_{x\to\pm\infty}\dfrac{x^{2}}{x}=\pm\infty$ (i.e. $+\infty$ at $+\infty$, $-\infty$ at $-\infty$).

\item \textbf{Vertical asymptote} (1 mark): from the infinite limits, $x=-2$.

\textbf{Oblique asymptote} (3 marks): division gives $x^{2}-3x+2=(x+2)(x-5)+12$, so
\[ f(x)=x-5+\frac{12}{x+2}, \qquad f(x)-(x-5)=\frac{12}{x+2}\to 0 \text{ at } \pm\infty. \]
Hence $y=x-5$ is the oblique asymptote at both ends. (No horizontal asymptote since the limits at infinity are infinite.)

\item \textbf{Derivative} (2 marks): with $u=x^{2}-3x+2$, $v=x+2$:
\[ f'(x)=\frac{(2x-3)(x+2)-(x^{2}-3x+2)}{(x+2)^{2}}=\frac{2x^{2}+x-6-x^{2}+3x-2}{(x+2)^{2}}=\frac{x^{2}+4x-8}{(x+2)^{2}}. \]

\textbf{Table of variation} (3 marks): $f'(x)=0 \iff x^{2}+4x-8=0$, $\Delta=16+32=48$, so
\[ x_{1}=-2-2\sqrt{3}\approx -5.46, \qquad x_{2}=-2+2\sqrt{3}\approx 1.46. \]
The numerator of $f'$ is positive outside $[x_1,x_2]$. Exact ordinates: since $f(x)=x-5+\dfrac{12}{x+2}$ and $x_i+2=\pm 2\sqrt{3}$:
\[ f(x_{1})=-7-2\sqrt{3}+\frac{12}{-2\sqrt{3}}=-7-2\sqrt{3}-2\sqrt{3}=-7-4\sqrt{3}\approx -13.93, \]
\[ f(x_{2})=-7+2\sqrt{3}+\frac{12}{2\sqrt{3}}=-7+2\sqrt{3}+2\sqrt{3}=-7+4\sqrt{3}\approx -0.07. \]
\begin{center}
\begin{tabular}{|c|ccccccccc|}
\hline
$x$ & $-\infty$ & & $x_{1}$ & & $-2$ & & $x_{2}$ & & $+\infty$ \\
\hline
$f'(x)$ & & $+$ & $0$ & $-$ & $\|$ & $-$ & $0$ & $+$ & \\
\hline
$f(x)$ & $-\infty$ & $\nearrow$ & $-7-4\sqrt{3}$ & $\searrow$ & $\|_{\ -\infty}^{+\infty}$ & $\searrow$ & $-7+4\sqrt{3}$ & $\nearrow$ & $+\infty$ \\
\hline
\end{tabular}
\end{center}
Local maximum at $x_1$ (value $-7-4\sqrt3$), local minimum at $x_2$ (value $-7+4\sqrt3$).

\item \textbf{Position} (2 marks): $f(x)-(x-5)=\dfrac{12}{x+2}$: above the asymptote for $x>-2$, below for $x<-2$.

\textbf{Intercepts} (1 mark): $f(0)=1$: point $(0,1)$. $f(x)=0 \iff x^{2}-3x+2=0 \iff x=1$ or $x=2$: points $(1,0)$ and $(2,0)$.

\item \textbf{Sketch} (3 marks: asymptotes drawn and labelled 1 mark; branches approaching asymptotes correctly 1 mark; stationary points and intercepts placed 1 mark).
\begin{popfigure}
\begin{tikzpicture}[scale=0.55]
\draw[->,gray] (-9,0) -- (6,0) node[right]{$x$};
\draw[->,gray] (0,-15) -- (0,6) node[above]{$y$};
\draw[dashed,popblue] (-2,-15) -- (-2,6) node[above]{$x=-2$};
\draw[dashed,poporange,domain=-9:5.5] plot (\x,{\x-5}) node[pos=0.05,above left]{$y=x-5$};
\draw[very thick,popink,domain=-9:-2.45,samples=60] plot (\x,{(\x*\x-3*\x+2)/(\x+2)});
\draw[very thick,popink,domain=-1.55:5.5,samples=60] plot (\x,{(\x*\x-3*\x+2)/(\x+2)});
\fill[popdark] (-5.464,-13.928) circle(2pt) node[below left]{max};
\fill[popdark] (1.464,-0.072) circle(2pt) node[above right]{min};
\fill[popgreen] (1,0) circle(2pt); \fill[popgreen] (2,0) circle(2pt); \fill[popgreen] (0,1) circle(2pt);
\end{tikzpicture}
\legende{Curve of $f(x)=\dfrac{x^{2}-3x+2}{x+2}$ with its asymptotes.}
\end{popfigure}

\item \textbf{Discussion} (2 marks): the horizontal line $y=k$ meets $(\mathcal{C})$:
\begin{itemize}
\item Left branch ($x<-2$): range $]-\infty, -7-4\sqrt{3}]$ (increases from $-\infty$ to max, then decreases to $-\infty$).
\item Right branch ($x>-2$): range $[-7+4\sqrt{3}, +\infty[$ (decreases from $+\infty$ to min, then increases to $+\infty$).
\item Since $-7-4\sqrt{3} \approx -13.93 < -7+4\sqrt{3} \approx -0.07$, the two ranges are disjoint.
\end{itemize}
Therefore:
\begin{itemize}
\item If $k < -7-4\sqrt{3}$: \textbf{two} solutions (both on the left branch).
\item If $k = -7-4\sqrt{3}$: \textbf{one} solution (tangent at the local maximum on the left branch).
\item If $-7-4\sqrt{3} < k < -7+4\sqrt{3}$: \textbf{zero} solutions (the line lies in the gap between the two branches).
\item If $k = -7+4\sqrt{3}$: \textbf{one} solution (tangent at the local minimum on the right branch).
\item If $k > -7+4\sqrt{3}$: \textbf{two} solutions (both on the right branch).
\end{itemize}
\emph{Examiner expectation:} the answer must be read from the table of variation/graph, with the boundary cases (tangency at extrema) stated explicitly.
\end{enumerate}
\end{corrige}

\begin{corrige}[Exercise 10 — Centre of symmetry and curve sketching (GCE style)]
$f(x)=\dfrac{2x^{2}-3x+4}{x-1}$.

\begin{enumerate}
\item \textbf{Domain} (1 mark): $D_f=\mathbb{R}\setminus\{1\}$.

\textbf{Limits} (1.5 marks): at $x=1$, numerator $\to 2-3+4=3>0$:
\[ \lim_{x\to 1^{+}}f(x)=+\infty, \qquad \lim_{x\to 1^{-}}f(x)=-\infty; \qquad \lim_{x\to\pm\infty}f(x)=\lim_{x\to\pm\infty}\frac{2x^{2}}{x}=\pm\infty. \]
\textbf{Vertical asymptote} (0.5 mark): $x=1$.

\item \textbf{Decomposition} (2 marks): $2x^{2}-3x+4=(x-1)(2x-1)+3$ (check: $(x-1)(2x-1)=2x^{2}-3x+1$, plus $3$ gives $2x^{2}-3x+4$). Hence
\[ f(x)=2x-1+\frac{3}{x-1}. \]
\textbf{Oblique asymptote} (1 mark): $f(x)-(2x-1)=\dfrac{3}{x-1}\to 0$ at $\pm\infty$, so $y=2x-1$ is the oblique asymptote.

\item \textbf{Centre of symmetry} (3 marks): The exercise statement suggests $I(1,2)$; computation reveals the actual centre is $I(1,1)$. For all $x\neq 0$,
\[ f(1+x)=2(1+x)-1+\frac{3}{x}=1+2x+\frac{3}{x}, \qquad f(1-x)=2(1-x)-1+\frac{3}{-x}=1-2x-\frac{3}{x}. \]
Adding: $f(1+x)+f(1-x)=2 = 2\times 1$. Hence the curve is symmetric about the point $I(1,1)$ (half-turn of angle $\pi$ about $I$).

\item \textbf{Derivative} (2 marks): from $f(x)=2x-1+\dfrac{3}{x-1}$:
\[ f'(x)=2-\frac{3}{(x-1)^{2}}=\frac{2(x-1)^{2}-3}{(x-1)^{2}}. \]
$f'(x)=0 \iff (x-1)^{2}=\dfrac32 \iff x=1\pm\sqrt{\dfrac32}=1\pm\dfrac{\sqrt6}{2}$.

\textbf{Table} (2 marks): $f'>0$ for $|x-1|>\frac{\sqrt6}{2}$:
\begin{center}
\begin{tabular}{|c|ccccccccc|}
\hline
$x$ & $-\infty$ & & $1-\frac{\sqrt6}{2}$ & & $1$ & & $1+\frac{\sqrt6}{2}$ & & $+\infty$ \\
\hline
$f'(x)$ & & $+$ & $0$ & $-$ & $\|$ & $-$ & $0$ & $+$ & \\
\hline
$f(x)$ & $-\infty$ & $\nearrow$ & $1-2\sqrt6$ & $\searrow$ & $\|_{\ -\infty}^{+\infty}$ & $\searrow$ & $1+2\sqrt6$ & $\nearrow$ & $+\infty$ \\
\hline
\end{tabular}
\end{center}
(Ordinates: $f\!\left(1\pm\tfrac{\sqrt6}{2}\right)=1\pm 2\sqrt6$, using $f(x)=2x-1+\frac{3}{x-1}$.)

\item \textbf{Sketch} (3 marks): draw the asymptotes $x=1$ and $y=2x-1$, plot the two stationary points $\left(1-\frac{\sqrt6}{2},\,1-2\sqrt6\right)\approx(-0.22,-3.90)$ and $\left(1+\frac{\sqrt6}{2},\,1+2\sqrt6\right)\approx(2.22,5.90)$, draw one branch, then obtain the other by the half-turn about $I(1,1)$ — this is the economy allowed by the centre of symmetry.

\item \textbf{Graphical solving} (4 marks):
\begin{enumerate}
\item $f(x)=2x-1 \iff \dfrac{3}{x-1}=0$: impossible. The curve never meets its oblique asymptote: $\mathcal{S}=\varnothing$.
\item $f(x)\leq 2x-1 \iff \dfrac{3}{x-1}\leq 0 \iff x-1<0 \iff x<1$. Graphically, the curve lies below (or on) the line $y=2x-1$ exactly on $]-\infty,1[$: $\mathcal{S}=\,]-\infty,1[$.
\end{enumerate}
\emph{Examiner expectation:} the graphical reading must be supported by the sign of $f(x)-(2x-1)$; answers must exclude $x=1$ (forbidden value).
\end{enumerate}
\end{corrige}

\begin{corrige}[Exercise 11 — Inequalities from a table of variation and absolute values (GCE Paper 2 style)]
\begin{enumerate}
\item \textbf{Sketch} (4 marks: increasing/decreasing branches consistent with the table 1 mark; extrema $(-1,3)$ and $(2,-1)$ placed 1 mark; zeros at $-2,0,3$ 1 mark; end behaviours $-2$ at $-\infty$ and $+\infty$ at $+\infty$ 1 mark). A consistent curve: starts just above $y=-2$ at $-\infty$, crosses the $x$-axis at $-2$, rises to the maximum $(-1,3)$, decreases through $(0,0)$ to the minimum $(2,-1)$, then increases through $(3,0)$ to $+\infty$.
\begin{popfigure}
\begin{tikzpicture}[scale=0.8]
\draw[->,gray] (-4.5,0) -- (4.5,0) node[right]{$x$};
\draw[->,gray] (0,-2.5) -- (0,3.8) node[above]{$y$};
\draw[very thick,popink,domain=-4.4:4.15,samples=120] plot (\x,{0.25*(\x+2)*\x*(\x-3)});
\fill[popdark] (-1,3) circle(2pt) node[above]{$(-1,3)$};
\fill[popdark] (2,-1) circle(2pt) node[below right]{$(2,-1)$};
\foreach \p in {-2,0,3}{\fill[popgreen] (\p,0) circle(2pt);}
\draw[dashed,popblue] (-4.5,-2) -- (-3.2,-2) node[right]{$y=-2$};
\end{tikzpicture}
\legende{A possible qualitative curve of $g$ satisfying the given table of variation.}
\end{popfigure}

\item \textbf{Graphical solutions} (2 marks each):
\begin{enumerate}
\item $g(x)=0$: read the intersections with the $x$-axis: $\mathcal{S}=\{-2,\,0,\,3\}$.
\item $g(x)>1$: draw the line $y=1$. It cuts the curve at a point of abscissa $\alpha\in\,]-2,-1[$ (on the rising branch), at a point $\beta\in\,]-1,0[$ (on the falling branch), and at a point $\gamma\in\,]3,+\infty[$ (on the final rising branch). Reading where the curve is strictly above the line:
\[ \mathcal{S}=\,]\alpha,\beta[\,\cup\,]\gamma,+\infty[. \]
\item $|g(x)|<2 \iff -2<g(x)<2$: draw the lines $y=2$ and $y=-2$. The line $y=2$ meets the curve at $a\in\,]-2,-1[$, $b\in\,]-1,0[$, $c\in\,]3,+\infty[$. The line $y=-2$ is a horizontal asymptote as $x\to-\infty$ (approached from above) and the minimum of $g$ is $-1>-2$, so the curve never meets $y=-2$. Hence $|g(x)|<2$ is equivalent to $g(x)<2$. The curve lies strictly below $y=2$ on:
\[ \mathcal{S}=\,]-\infty,a[\,\cup\,]b,c[. \]
\end{enumerate}

\item \textbf{Curve of $y=|g(x)|$} (3 marks): keep unchanged every part of $(\mathcal{C})$ lying above the $x$-axis; reflect in the $x$-axis every part lying below it. Concretely: the arcs on $]-\infty,-2[$ (where $g<0$) and on $]0,3[$ (where $g<0$) are flipped upwards; the rest is unchanged. The zeros $-2,0,3$ stay fixed, and the local minimum $(2,-1)$ becomes a local maximum $(2,1)$ of $|g|$.

\item \textbf{Solving $|x-2|<|2x+1|$} (7 marks):

\textbf{(a) By squaring} (both sides $\geq 0$, squaring preserves the strict order):
\[ (x-2)^{2}<(2x+1)^{2} \iff x^{2}-4x+4<4x^{2}+4x+1 \iff 0<3x^{2}+8x-3. \]
Roots of $3x^{2}+8x-3=0$: $\Delta=64+36=100$, $x=\dfrac{-8\pm 10}{6}$, giving $x=\dfrac13$ and $x=-3$. The quadratic is positive outside its roots:
\[ \mathcal{S}=\,]-\infty,-3[\,\cup\,\left]\tfrac13,+\infty\right[. \]

\textbf{(b) Critical-value method:} the expressions change sign at $x=2$ and $x=-\tfrac12$; study three intervals.
\begin{itemize}
\item On $]-\infty,-\tfrac12[$: $|x-2|=-(x-2)=2-x$, $|2x+1|=-(2x+1)$. Inequality: $2-x<-2x-1 \iff x<-3$. Intersection with the interval: $]-\infty,-3[$.
\item On $[-\tfrac12,2]$: $|x-2|=2-x$, $|2x+1|=2x+1$. Inequality: $2-x<2x+1 \iff x>\tfrac13$. Intersection: $]\tfrac13,2]$.
\item On $]2,+\infty[$: $|x-2|=x-2$, $|2x+1|=2x+1$. Inequality: $x-2<2x+1 \iff x>-3$: always true here. Intersection: $]2,+\infty[$.
\end{itemize}
Union: $]-\infty,-3[\,\cup\,]\tfrac13,2]\,\cup\,]2,+\infty[\,=\,]-\infty,-3[\,\cup\,]\tfrac13,+\infty[$.

\textbf{Comparison:} both methods agree, $\mathcal{S}=\,]-\infty,-3[\,\cup\,\left]\tfrac13,+\infty\right[$. Squaring is faster but is valid \emph{only} because both sides are non-negative; the critical-value method is longer but works in all cases and reinforces the definition of the absolute value.
\end{enumerate}
\end{corrige}

\newpage

\coursec{Summary sheet}

\begin{popfigure}
\begin{tikzpicture}[mindmap, grow cyclic, every node/.style={concept, text=white, font=\small},
  root concept/.append style={concept color=popink, font=\bfseries\small},
  level 1/.append style={level distance=3.4cm, sibling angle=60},
  level 1 concept/.append style={font=\scriptsize, minimum size=2.1cm}]
\node[root concept]{Curve Sketching \& Inequalities}
  child[concept color=popblue]{ node {Domain \& limits at the bounds} }
  child[concept color=poporange]{ node {Asymptotes, intercepts, infinite branches} }
  child[concept color=popgreen]{ node {Stationary points \& variation table} }
  child[concept color=poppurple]{ node {Concavity, inflexion, symmetry} }
  child[concept color=poppink]{ node {Absolute-value inequalities (algebraic)} }
  child[concept color=popgold]{ node {Graphical solutions \& full sketch} };
\end{tikzpicture}
\legende{Mind map of the chapter: six connected skills leading to a complete study of a curve.}
\end{popfigure}

\begin{bilan}
\textbf{Key definitions}
\begin{itemize}
\item \cle{Domain of definition} $D_f$: the set of real $x$ for which $f(x)$ exists. Exclude zeros of denominators, impose radicands $\ge 0$, arguments of $\ln$ strictly positive.
\item \cle{Limits at the bounds}: compute $\lim_{x\to a}f(x)$ at every endpoint of $D_f$ (finite or infinite); they reveal asymptotes and infinite branches.
\item \cle{Stationary point}: point where $f'(x)=0$ (or $f'$ undefined); a local maximum, minimum, or horizontal inflexion.
\item \cle{Point of inflexion}: point where the concavity changes; typically $f''$ changes sign there.
\item \cle{Centre of symmetry} $\Omega(a,b)$: $f(a+x)+f(a-x)=2b$; \cle{axis of symmetry} $x=a$: $f(a+x)=f(a-x)$.
\end{itemize}

\textbf{Formulas and theorems to know}
\begin{itemize}
\item Vertical asymptote $x=a$: $\lim_{x\to a}f(x)=\pm\infty$.
\item Horizontal asymptote $y=b$: $\lim_{x\to\pm\infty}f(x)=b$.
\item Oblique asymptote $y=mx+c$: $m=\lim_{x\to\pm\infty}\dfrac{f(x)}{x}$, $c=\lim_{x\to\pm\infty}[f(x)-mx]$, with $m$ finite and non-zero.
\item Parabolic branch: $\lim f(x)=\pm\infty$ and $\lim \dfrac{f(x)}{x}=0$ (axis $(Ox)$) or $\pm\infty$ (axis $(Oy)$).
\item Sign of $f(x)-(mx+c)$ gives the \cle{position of the curve} relative to its asymptote.
\item $f'(x)>0 \Rightarrow f$ increasing; $f'(x)<0 \Rightarrow f$ decreasing; sign change of $f'$ at a stationary point gives a maximum ($+$ to $-$) or minimum ($-$ to $+$).
\item $f''>0$: concave up (convex); $f''<0$: concave down; change of sign of $f''$ $\Rightarrow$ inflexion.
\item $|u|<a \iff -a<u<a$ ($a>0$); $|u|>a \iff u<-a$ or $u>a$; $|u|=|v| \iff u=v$ or $u=-v$; $|u|<|v| \iff u^2<v^2$.
\end{itemize}

\textbf{Methods}
\begin{itemize}
\item \cle{Complete study}: domain $\to$ symmetry (reduce domain) $\to$ limits at bounds $\to$ asymptotes $\to$ $f'$ and variation table $\to$ $f''$, concavity, inflexions $\to$ intercepts and remarkable points $\to$ sketch.
\item \cle{Absolute-value inequalities}: split into cases according to the sign of each expression inside modulus, solve on each interval, then take the union of the valid parts. Alternatively square when both sides are non-negative.
\item \cle{Graphical solution}: sketch $y=|f(x)|$ (reflect negative parts in the $x$-axis) and $y=g(x)$; read off the $x$-values where one curve lies above the other; confirm intersection points algebraically.
\item \cle{Range of existence}: for $y=f(x)$, solve for $x$ in terms of $y$ and impose the condition for real solutions (e.g. discriminant $\ge 0$ for a quadratic in $x$).
\end{itemize}

\textbf{Common mistakes}
\begin{itemize}
\item Forgetting excluded values when finding the domain, or studying limits only at $+\infty$.
\item Claiming a maximum merely because $f'(x)=0$: a sign change of $f'$ is required.
\item Confusing $f''(x)=0$ with an inflexion: the sign of $f''$ must \emph{change}.
\item Squaring an inequality when one side may be negative: this reverses or falsifies the result.
\item Writing $|x|<a \Rightarrow x<-a$ or $x>a$: that is the answer to $|x|>a$.
\item Forgetting to reject case-solutions that fall outside the case interval.
\item Misplacing the curve relative to its asymptote: always study the sign of the difference.
\end{itemize}

\textbf{GCE tips}
\begin{itemize}
\item Present a clean \cle{variation table} gathering $x$, $f'(x)$, $f(x)$, with limits at every bound: examiners award marks for its structure.
\item State each asymptote with its justification (the limit) before drawing it; draw asymptotes as dashed lines.
\item For ``show that the line $y=mx+c$ is an asymptote'', compute $f(x)-(mx+c)$ and its limit $0$.
\item In inequality questions, give the solution in \cle{interval notation} and check one test value per interval.
\item Use symmetry to halve the work: justify it explicitly, then state ``we complete by symmetry''.
\item Manage time: a full curve study is long; do the algebra neatly first, sketch last.
\end{itemize}
\end{bilan}

\findecours

\end{document}
